% Intelligence Is Not Authority
% Self-contained source: all figures, numbers and references are in this file.
% Build:  pdflatex intelligence_is_not_authority   (three passes; no BibTeX)
% Every experimental number is produced by experiments.py (seed 20261110).
\documentclass[10pt,twocolumn,letterpaper]{article}

\usepackage[T1]{fontenc}
\usepackage[utf8]{inputenc}
\usepackage{mathptmx}
\usepackage[scaled=0.9]{helvet}
\usepackage[letterpaper,top=0.9in,bottom=0.95in,left=0.72in,right=0.72in,columnsep=0.28in,headsep=0.22in]{geometry}
\usepackage{amsmath,amssymb,amsthm}
\usepackage{microtype}
\usepackage{graphicx}
\usepackage{booktabs,tabularx,array}
\usepackage[font=small,labelfont=bf,skip=5pt]{caption}
\usepackage{titlesec}
\usepackage{enumitem}
\usepackage{xcolor}
\usepackage[most]{tcolorbox}
\usepackage{tikz}
\usetikzlibrary{arrows.meta,shapes.geometric}
\usepackage{pgfplots}
\pgfplotsset{compat=1.18}
\usepackage[numbers,sort&compress]{natbib}
\usepackage{fancyhdr}
\usepackage{flushend}
\usepackage{xurl}
\urlstyle{same}
\usepackage[colorlinks=true,linkcolor=accent,citecolor=accent,urlcolor=accent]{hyperref}

\definecolor{accent}{HTML}{1B4F72}
\definecolor{warn}{HTML}{A93226}
\definecolor{ink}{HTML}{2B2B2B}
\definecolor{mute}{HTML}{8A8A86}
\definecolor{boxbg}{HTML}{F4F3EF}
\definecolor{earth}{HTML}{2F5D4E}
\definecolor{earthbg}{HTML}{F1F5F1}
\definecolor{sea}{HTML}{3F7FA6}

\titleformat{\section}{\large\bfseries}{\thesection}{0.6em}{}
\titleformat{\subsection}{\normalsize\bfseries}{\thesubsection}{0.6em}{}
\titlespacing*{\section}{0pt}{10pt plus 2pt minus 1pt}{4pt}
\titlespacing*{\subsection}{0pt}{7pt plus 2pt minus 1pt}{3pt}
\setlist{leftmargin=1.3em,itemsep=1.5pt,topsep=3pt,parsep=0pt}
\setlength{\bibsep}{1.5pt}
\raggedbottom

\pagestyle{fancy}
\fancyhf{}
\fancyhead[L]{\small\itshape Intelligence Is Not Authority}
\fancyhead[R]{\small\itshape Preprint, November 2026}
\fancyfoot[C]{\small\thepage}
\renewcommand{\headrulewidth}{0.3pt}

\newtheorem{proposition}{Proposition}
\newtheorem{theorem}{Theorem}
\newtheorem{corollary}{Corollary}
\theoremstyle{definition}
\newtheorem*{definition}{Definition}
\newcommand{\Pb}{\mathbb{P}}
\newcommand{\Ex}{\mathbb{E}}
\newcommand{\Inv}{\mathcal{I}}
\newcommand{\para}[1]{\smallskip\noindent\textbf{#1}\enspace}
\newcolumntype{L}{>{\raggedright\arraybackslash}X}
\newcolumntype{P}[1]{>{\raggedright\arraybackslash}p{#1}}
\newtcolorbox{claimbox}[1]{enhanced,colback=boxbg,colframe=mute,boxrule=0.4pt,arc=1.5pt,
  left=5pt,right=5pt,top=4pt,bottom=4pt,fonttitle=\bfseries\small,title={#1},
  coltitle=black,colbacktitle=boxbg,titlerule=0pt,before skip=7pt,after skip=7pt}

% ---- numbers written by experiments.py ----
\newcommand{\NarrowPaths}{8.4}
\newcommand{\HubPaths}{7.6}
\newcommand{\WidePaths}{7.6}
\newcommand{\WidthNarrow}{6}
\newcommand{\WidthHub}{21}
\newcommand{\WidthWide}{27}
\newcommand{\RepCutNarrow}{6}
\newcommand{\RepCutHub}{11}
\newcommand{\RepCutWide}{18}
\newcommand{\SeenPerHundredK}{four}
\newcommand{\CutSeverNarrowSix}{67}
\newcommand{\CutBlockNarrowSix}{94}
\newcommand{\CutBlockHubSix}{52}
\newcommand{\CutBlockHubTwelve}{75}
\newcommand{\CutGainHubTwelve}{42}
\newcommand{\CutBlockWideSix}{25}
\newcommand{\PatchSpecBlockNarrow}{74}
\newcommand{\PatchSpecGainNarrow}{6}
\newcommand{\PatchRandBlockNarrow}{74}
\newcommand{\PatchRandGainNarrow}{48}
\newcommand{\PatchSpecBlockHub}{84}
\newcommand{\PatchSpecGainHub}{16}
\newcommand{\PatchRandBlockHub}{75}
\newcommand{\PatchRandGainHub}{47}
\newcommand{\PatchSpecBlockWide}{76}
\newcommand{\PatchSpecGainWide}{9}
\newcommand{\PatchRandBlockWide}{67}
\newcommand{\PatchRandGainWide}{44}
\newcommand{\PatchSpecSeverNarrow}{0}
\newcommand{\PatchRandSeverNarrow}{1.5}
\newcommand{\SeverNThirty}{2}
\newcommand{\SeverNThousand}{98}
\newcommand{\SeverWideMax}{5}
\newcommand{\EstHonest}{95.7}
\newcommand{\TrueHonest}{95.7}
\newcommand{\EstBlind}{96.6}
\newcommand{\TrueBlind}{82.0}
\newcommand{\EstWith}{96.0}
\newcommand{\TrueWith}{91.4}
\newcommand{\CritHonest}{94.0}
\newcommand{\CritBlind}{81.9}
\newcommand{\CritWith}{0}
\newcommand{\AuditFortyHalf}{89}
\newcommand{\AuditFortyEighty}{17}
\newcommand{\AuditFortyNinety}{4}
\newcommand{\AuditTwoHundredEighty}{88}
\newcommand{\AuditTwoHundredNinety}{29}
\newcommand{\CopSixFour}{1.4}
\newcommand{\CopSixEight}{0.54}
\newcommand{\CopThreeFour}{0.28}
\newcommand{\CopThreeEight}{0.026}
\newcommand{\BndStates}{10{,}180}
\newcommand{\BndChecks}{223{,}960}
\newcommand{\BndTasks}{20}
\newcommand{\BndRejectStates}{740}
\newcommand{\BndPrepStates}{956}
\newcommand{\BndActions}{22}
\newcommand{\BndOmitStates}{41{,}732}
\newcommand{\BndOmitChecks}{918{,}104}
\newcommand{\BndOmitSpend}{eight}
\newcommand{\ScWar}{95}
\newcommand{\ScConf}{95}
\newcommand{\ScCare}{10}
\newcommand{\ScFloor}{4}
\newcommand{\ScHarmWar}{4.7}
\newcommand{\ScHarmCare}{0.5}
\newcommand{\ScHarmFloor}{0.2}
\newcommand{\ScHarmRatio}{21}
\newcommand{\ScCapWar}{5.7}
\newcommand{\ScCapCare}{1.6}
\newcommand{\ScCapFloor}{3.7}
\newcommand{\ScCapHalf}{6.3}
\newcommand{\ScConfHarmPct}{60}
\newcommand{\ScConfSixty}{92}
\newcommand{\ScTopLowPct}{30}
\newcommand{\ScTopLow}{88}
\newcommand{\ScTopHalf}{4}
\newcommand{\ScGateSixty}{92}
\newcommand{\ScGateEighty}{2}
\newcommand{\MonTasks}{24}
\newcommand{\MonSameMarg}{22}
\newcommand{\MonSameJointSix}{5.5}
\newcommand{\MonSameProdSix}{0.034}
\newcommand{\MonSameMargAtk}{50}
\newcommand{\MonSameJointSixAtk}{25}
\newcommand{\MonSameCI}{1.2}
\newcommand{\MonMixedMarg}{33}
\newcommand{\MonMixedJointSix}{8.5}
\newcommand{\MonMixedProdSix}{0.14}
\newcommand{\MonMixedMargAtk}{58}
\newcommand{\MonMixedJointSixAtk}{26}
\newcommand{\MonMixedCI}{1.2}
\newcommand{\MonDistJointSix}{0.14}
\newcommand{\MonDistProdSix}{0.14}
\newcommand{\MonDistJointSixAtk}{4.4}
\newcommand{\MonDistProdSixAtk}{4.4}
\newcommand{\MonDistCI}{0.074}
\newcommand{\MonSameRatio}{160}
\newcommand{\MonMixedRatio}{63}
\newcommand{\MonRatioLow}{63}
\newcommand{\MonRatioHigh}{160}
\newcommand{\MonSameDepthGain}{1.7}
\newcommand{\MonSameIndepGain}{42}
\newcommand{\MonSurrPass}{75}
\newcommand{\MonUnderstate}{190}
\newcommand{\EstWithRound}{96}

\hypersetup{pdftitle={Intelligence Is Not Authority: A bound on catastrophic AI risk with no term for what the machine wants},pdfauthor={Jepson Taylor}}

\begin{document}
\makeatletter
\twocolumn[{%
\begin{@twocolumnfalse}
\thispagestyle{plain}
\begin{center}
{\fontsize{25}{29}\selectfont\bfseries Intelligence Is {\color{warn}Not} Authority\par}
\vspace{6pt}
{\fontsize{12.5}{16}\selectfont A bound on catastrophic AI risk with no term for what the machine wants\par}
\vspace{10pt}
{\large Jepson Taylor\par}
\vspace{2pt}
{\normalsize\href{mailto:jepson@veox.ai}{jepson@veox.ai}\par}
\vspace{2pt}
{\small November 10, 2026\par}
\vspace{10pt}
\begin{minipage}{0.9\textwidth}
\small
\noindent\textbf{Abstract.}\enspace
No published method can verify what an AI system more capable than its overseers wants, and in 2026 the evidence such methods rely on weakened. We separate what a system can think from what it can reach, and bound the second. For every policy, the probability of irreversible catastrophe in one operating period is at most $u+J_{k}$: the chance that a pathway around the safeguards exists and is found, plus the measured chance that $k$ safeguards on the shallowest pathway all miss. No term refers to goals and no independence is assumed. In the standard reliability model $J_{k}\le c+m^{k}$, so added safeguards stop helping at a depth we compute. The defender pays for the width of the pathway space, never for the number of pathways. Search can prove danger but not safety, and withheld findings are invisible in the reports. Five experiments test the framework. A stateful boundary, checked exhaustively, admits no violation in \BndChecks{} state--action pairs; removing any one check yields a counterexample within five steps, and a check against an incomplete specification passes while the harm occurs. Six cuts sever all eight million pathways of a narrow graph in \CutSeverNarrowSix\% of trials. One result corrected us: six monitors stacked on one step missed together \MonRatioLow{} to \MonRatioHigh{} times as often as independence predicts, whatever their model family, while the same monitors on distinct steps multiplied. The reward is the remaining choice. Evolution built warriors, paid to win, and grandmothers, paid only through those they keep alive. A prize for rank pays any agent behind the bar to gamble: under selection it put \ScWar\% of effort into a harmful gamble, against \ScFloor\% under a floor that most can clear, which kept two thirds of the gain in capacity. This does not solve alignment. It reduces survival to six measurable conditions, shows where the evidence stands, and states what would refute it.
\end{minipage}
\end{center}
\vspace{8pt}
% ------------------------------ teaser ------------------------------
\noindent
\begin{minipage}[t]{0.325\textwidth}
\centering
{\small\bfseries A.\enspace The problem\par}\vspace{3pt}
\begin{tikzpicture}
\begin{axis}[width=\textwidth,height=4.5cm,xmode=log,xmin=1e-2,xmax=1e4,ymin=0,ymax=1.05,
  xlabel={tests per deployment decision, $N/M$},ylabel={chance of catching it},
  ytick={0,0.5,1},xtick={1e-2,1,1e2,1e4},grid=major,grid style={black!10},
  tick label style={font=\scriptsize},label style={font=\scriptsize},domain=0.01:10000,samples=120,
  ylabel style={yshift=-3pt}]
\addplot[accent,line width=1pt] {1-exp(-x*ln(2))};
\addplot[accent,line width=1pt,dashed] {1-exp(-x*ln(2)/10)};
\addplot[accent,line width=1pt,dotted] {1-exp(-x*ln(2)/100)};
\node[font=\scriptsize,anchor=south east] at (axis cs:0.7,0.52) {$\Lambda{=}1$};
\node[font=\scriptsize,anchor=west] at (axis cs:11,0.42) {$10$};
\node[font=\scriptsize,anchor=west] at (axis cs:110,0.42) {$100$};
\draw[warn,line width=0.7pt,<-] (axis cs:0.012,0.78) -- (axis cs:0.7,0.78);
\node[font=\scriptsize,text=warn,anchor=south west] at (axis cs:0.011,0.79) {today: $N\ll M$};
\end{axis}
\end{tikzpicture}
\end{minipage}\hfill
\begin{minipage}[t]{0.325\textwidth}
\centering
{\small\bfseries B.\enspace The leverage\par}\vspace{3pt}
\begin{tikzpicture}[x=1cm,y=1cm,font=\scriptsize]
\draw[black!22,line width=0.25pt] (0.0,0.27) -- (0.55,0.27);
\draw[black!22,line width=0.25pt] (0.0,0.27) -- (0.55,0.81);
\draw[black!22,line width=0.25pt] (0.0,0.27) -- (0.55,1.35);
\draw[black!22,line width=0.25pt] (0.0,0.27) -- (0.55,1.89);
\draw[black!22,line width=0.25pt] (0.0,0.27) -- (0.55,2.43);
\draw[black!22,line width=0.25pt] (0.0,0.81) -- (0.55,0.27);
\draw[black!22,line width=0.25pt] (0.0,0.81) -- (0.55,0.81);
\draw[black!22,line width=0.25pt] (0.0,0.81) -- (0.55,1.35);
\draw[black!22,line width=0.25pt] (0.0,0.81) -- (0.55,1.89);
\draw[black!22,line width=0.25pt] (0.0,0.81) -- (0.55,2.43);
\draw[black!22,line width=0.25pt] (0.0,1.35) -- (0.55,0.27);
\draw[black!22,line width=0.25pt] (0.0,1.35) -- (0.55,0.81);
\draw[black!22,line width=0.25pt] (0.0,1.35) -- (0.55,1.35);
\draw[black!22,line width=0.25pt] (0.0,1.35) -- (0.55,1.89);
\draw[black!22,line width=0.25pt] (0.0,1.35) -- (0.55,2.43);
\draw[black!22,line width=0.25pt] (0.0,1.89) -- (0.55,0.27);
\draw[black!22,line width=0.25pt] (0.0,1.89) -- (0.55,0.81);
\draw[black!22,line width=0.25pt] (0.0,1.89) -- (0.55,1.35);
\draw[black!22,line width=0.25pt] (0.0,1.89) -- (0.55,1.89);
\draw[black!22,line width=0.25pt] (0.0,1.89) -- (0.55,2.43);
\draw[black!22,line width=0.25pt] (0.0,2.43) -- (0.55,0.27);
\draw[black!22,line width=0.25pt] (0.0,2.43) -- (0.55,0.81);
\draw[black!22,line width=0.25pt] (0.0,2.43) -- (0.55,1.35);
\draw[black!22,line width=0.25pt] (0.0,2.43) -- (0.55,1.89);
\draw[black!22,line width=0.25pt] (0.0,2.43) -- (0.55,2.43);
\draw[black!22,line width=0.25pt] (0.55,0.27) -- (1.1,0.972);
\draw[black!22,line width=0.25pt] (0.55,0.27) -- (1.1,1.728);
\draw[black!22,line width=0.25pt] (0.55,0.81) -- (1.1,0.972);
\draw[black!22,line width=0.25pt] (0.55,0.81) -- (1.1,1.728);
\draw[black!22,line width=0.25pt] (0.55,1.35) -- (1.1,0.972);
\draw[black!22,line width=0.25pt] (0.55,1.35) -- (1.1,1.728);
\draw[black!22,line width=0.25pt] (0.55,1.89) -- (1.1,0.972);
\draw[black!22,line width=0.25pt] (0.55,1.89) -- (1.1,1.728);
\draw[black!22,line width=0.25pt] (0.55,2.43) -- (1.1,0.972);
\draw[black!22,line width=0.25pt] (0.55,2.43) -- (1.1,1.728);
\draw[black!22,line width=0.25pt] (1.1,0.972) -- (1.65,0.27);
\draw[black!22,line width=0.25pt] (1.1,0.972) -- (1.65,0.81);
\draw[black!22,line width=0.25pt] (1.1,0.972) -- (1.65,1.35);
\draw[black!22,line width=0.25pt] (1.1,0.972) -- (1.65,1.89);
\draw[black!22,line width=0.25pt] (1.1,0.972) -- (1.65,2.43);
\draw[black!22,line width=0.25pt] (1.1,1.728) -- (1.65,0.27);
\draw[black!22,line width=0.25pt] (1.1,1.728) -- (1.65,0.81);
\draw[black!22,line width=0.25pt] (1.1,1.728) -- (1.65,1.35);
\draw[black!22,line width=0.25pt] (1.1,1.728) -- (1.65,1.89);
\draw[black!22,line width=0.25pt] (1.1,1.728) -- (1.65,2.43);
\draw[black!22,line width=0.25pt] (1.65,0.27) -- (2.2,0.27);
\draw[black!22,line width=0.25pt] (1.65,0.27) -- (2.2,0.81);
\draw[black!22,line width=0.25pt] (1.65,0.27) -- (2.2,1.35);
\draw[black!22,line width=0.25pt] (1.65,0.27) -- (2.2,1.89);
\draw[black!22,line width=0.25pt] (1.65,0.27) -- (2.2,2.43);
\draw[black!22,line width=0.25pt] (1.65,0.81) -- (2.2,0.27);
\draw[black!22,line width=0.25pt] (1.65,0.81) -- (2.2,0.81);
\draw[black!22,line width=0.25pt] (1.65,0.81) -- (2.2,1.35);
\draw[black!22,line width=0.25pt] (1.65,0.81) -- (2.2,1.89);
\draw[black!22,line width=0.25pt] (1.65,0.81) -- (2.2,2.43);
\draw[black!22,line width=0.25pt] (1.65,1.35) -- (2.2,0.27);
\draw[black!22,line width=0.25pt] (1.65,1.35) -- (2.2,0.81);
\draw[black!22,line width=0.25pt] (1.65,1.35) -- (2.2,1.35);
\draw[black!22,line width=0.25pt] (1.65,1.35) -- (2.2,1.89);
\draw[black!22,line width=0.25pt] (1.65,1.35) -- (2.2,2.43);
\draw[black!22,line width=0.25pt] (1.65,1.89) -- (2.2,0.27);
\draw[black!22,line width=0.25pt] (1.65,1.89) -- (2.2,0.81);
\draw[black!22,line width=0.25pt] (1.65,1.89) -- (2.2,1.35);
\draw[black!22,line width=0.25pt] (1.65,1.89) -- (2.2,1.89);
\draw[black!22,line width=0.25pt] (1.65,1.89) -- (2.2,2.43);
\draw[black!22,line width=0.25pt] (1.65,2.43) -- (2.2,0.27);
\draw[black!22,line width=0.25pt] (1.65,2.43) -- (2.2,0.81);
\draw[black!22,line width=0.25pt] (1.65,2.43) -- (2.2,1.35);
\draw[black!22,line width=0.25pt] (1.65,2.43) -- (2.2,1.89);
\draw[black!22,line width=0.25pt] (1.65,2.43) -- (2.2,2.43);
\fill[ink] (0.0,0.27) circle (1.3pt);
\fill[ink] (0.0,0.81) circle (1.3pt);
\fill[ink] (0.0,1.35) circle (1.3pt);
\fill[ink] (0.0,1.89) circle (1.3pt);
\fill[ink] (0.0,2.43) circle (1.3pt);
\fill[ink] (0.55,0.27) circle (1.3pt);
\fill[ink] (0.55,0.81) circle (1.3pt);
\fill[ink] (0.55,1.35) circle (1.3pt);
\fill[ink] (0.55,1.89) circle (1.3pt);
\fill[ink] (0.55,2.43) circle (1.3pt);
\fill[warn] (1.1,0.972) circle (2.3pt);
\fill[warn] (1.1,1.728) circle (2.3pt);
\fill[ink] (1.65,0.27) circle (1.3pt);
\fill[ink] (1.65,0.81) circle (1.3pt);
\fill[ink] (1.65,1.35) circle (1.3pt);
\fill[ink] (1.65,1.89) circle (1.3pt);
\fill[ink] (1.65,2.43) circle (1.3pt);
\fill[ink] (2.2,0.27) circle (1.3pt);
\fill[ink] (2.2,0.81) circle (1.3pt);
\fill[ink] (2.2,1.35) circle (1.3pt);
\fill[ink] (2.2,1.89) circle (1.3pt);
\fill[ink] (2.2,2.43) circle (1.3pt);
\draw[black!22,line width=0.25pt] (3.05,0.338) -- (3.6,0.338);
\draw[black!22,line width=0.25pt] (3.05,0.338) -- (3.6,1.013);
\draw[black!22,line width=0.25pt] (3.05,0.338) -- (3.6,1.688);
\draw[black!22,line width=0.25pt] (3.05,0.338) -- (3.6,2.363);
\draw[black!22,line width=0.25pt] (3.05,1.013) -- (3.6,0.338);
\draw[black!22,line width=0.25pt] (3.05,1.013) -- (3.6,1.013);
\draw[black!22,line width=0.25pt] (3.05,1.013) -- (3.6,1.688);
\draw[black!22,line width=0.25pt] (3.05,1.013) -- (3.6,2.363);
\draw[black!22,line width=0.25pt] (3.05,1.688) -- (3.6,0.338);
\draw[black!22,line width=0.25pt] (3.05,1.688) -- (3.6,1.013);
\draw[black!22,line width=0.25pt] (3.05,1.688) -- (3.6,1.688);
\draw[black!22,line width=0.25pt] (3.05,1.688) -- (3.6,2.363);
\draw[black!22,line width=0.25pt] (3.05,2.363) -- (3.6,0.338);
\draw[black!22,line width=0.25pt] (3.05,2.363) -- (3.6,1.013);
\draw[black!22,line width=0.25pt] (3.05,2.363) -- (3.6,1.688);
\draw[black!22,line width=0.25pt] (3.05,2.363) -- (3.6,2.363);
\draw[black!22,line width=0.25pt] (3.6,0.338) -- (4.15,0.338);
\draw[black!22,line width=0.25pt] (3.6,0.338) -- (4.15,1.013);
\draw[black!22,line width=0.25pt] (3.6,0.338) -- (4.15,1.688);
\draw[black!22,line width=0.25pt] (3.6,0.338) -- (4.15,2.363);
\draw[black!22,line width=0.25pt] (3.6,1.013) -- (4.15,0.338);
\draw[black!22,line width=0.25pt] (3.6,1.013) -- (4.15,1.013);
\draw[black!22,line width=0.25pt] (3.6,1.013) -- (4.15,1.688);
\draw[black!22,line width=0.25pt] (3.6,1.013) -- (4.15,2.363);
\draw[black!22,line width=0.25pt] (3.6,1.688) -- (4.15,0.338);
\draw[black!22,line width=0.25pt] (3.6,1.688) -- (4.15,1.013);
\draw[black!22,line width=0.25pt] (3.6,1.688) -- (4.15,1.688);
\draw[black!22,line width=0.25pt] (3.6,1.688) -- (4.15,2.363);
\draw[black!22,line width=0.25pt] (3.6,2.363) -- (4.15,0.338);
\draw[black!22,line width=0.25pt] (3.6,2.363) -- (4.15,1.013);
\draw[black!22,line width=0.25pt] (3.6,2.363) -- (4.15,1.688);
\draw[black!22,line width=0.25pt] (3.6,2.363) -- (4.15,2.363);
\draw[black!22,line width=0.25pt] (4.15,0.338) -- (4.7,0.338);
\draw[black!22,line width=0.25pt] (4.15,0.338) -- (4.7,1.013);
\draw[black!22,line width=0.25pt] (4.15,0.338) -- (4.7,1.688);
\draw[black!22,line width=0.25pt] (4.15,0.338) -- (4.7,2.363);
\draw[black!22,line width=0.25pt] (4.15,1.013) -- (4.7,0.338);
\draw[black!22,line width=0.25pt] (4.15,1.013) -- (4.7,1.013);
\draw[black!22,line width=0.25pt] (4.15,1.013) -- (4.7,1.688);
\draw[black!22,line width=0.25pt] (4.15,1.013) -- (4.7,2.363);
\draw[black!22,line width=0.25pt] (4.15,1.688) -- (4.7,0.338);
\draw[black!22,line width=0.25pt] (4.15,1.688) -- (4.7,1.013);
\draw[black!22,line width=0.25pt] (4.15,1.688) -- (4.7,1.688);
\draw[black!22,line width=0.25pt] (4.15,1.688) -- (4.7,2.363);
\draw[black!22,line width=0.25pt] (4.15,2.363) -- (4.7,0.338);
\draw[black!22,line width=0.25pt] (4.15,2.363) -- (4.7,1.013);
\draw[black!22,line width=0.25pt] (4.15,2.363) -- (4.7,1.688);
\draw[black!22,line width=0.25pt] (4.15,2.363) -- (4.7,2.363);
\draw[black!22,line width=0.25pt] (4.7,0.338) -- (5.25,0.338);
\draw[black!22,line width=0.25pt] (4.7,0.338) -- (5.25,1.013);
\draw[black!22,line width=0.25pt] (4.7,0.338) -- (5.25,1.688);
\draw[black!22,line width=0.25pt] (4.7,0.338) -- (5.25,2.363);
\draw[black!22,line width=0.25pt] (4.7,1.013) -- (5.25,0.338);
\draw[black!22,line width=0.25pt] (4.7,1.013) -- (5.25,1.013);
\draw[black!22,line width=0.25pt] (4.7,1.013) -- (5.25,1.688);
\draw[black!22,line width=0.25pt] (4.7,1.013) -- (5.25,2.363);
\draw[black!22,line width=0.25pt] (4.7,1.688) -- (5.25,0.338);
\draw[black!22,line width=0.25pt] (4.7,1.688) -- (5.25,1.013);
\draw[black!22,line width=0.25pt] (4.7,1.688) -- (5.25,1.688);
\draw[black!22,line width=0.25pt] (4.7,1.688) -- (5.25,2.363);
\draw[black!22,line width=0.25pt] (4.7,2.363) -- (5.25,0.338);
\draw[black!22,line width=0.25pt] (4.7,2.363) -- (5.25,1.013);
\draw[black!22,line width=0.25pt] (4.7,2.363) -- (5.25,1.688);
\draw[black!22,line width=0.25pt] (4.7,2.363) -- (5.25,2.363);
\fill[ink] (3.05,0.338) circle (1.3pt);
\fill[ink] (3.05,1.013) circle (1.3pt);
\fill[ink] (3.05,1.688) circle (1.3pt);
\fill[ink] (3.05,2.363) circle (1.3pt);
\fill[ink] (3.6,0.338) circle (1.3pt);
\fill[ink] (3.6,1.013) circle (1.3pt);
\fill[ink] (3.6,1.688) circle (1.3pt);
\fill[ink] (3.6,2.363) circle (1.3pt);
\fill[warn] (4.15,0.338) circle (2.3pt);
\fill[warn] (4.15,1.013) circle (2.3pt);
\fill[warn] (4.15,1.688) circle (2.3pt);
\fill[warn] (4.15,2.363) circle (2.3pt);
\fill[ink] (4.7,0.338) circle (1.3pt);
\fill[ink] (4.7,1.013) circle (1.3pt);
\fill[ink] (4.7,1.688) circle (1.3pt);
\fill[ink] (4.7,2.363) circle (1.3pt);
\fill[ink] (5.25,0.338) circle (1.3pt);
\fill[ink] (5.25,1.013) circle (1.3pt);
\fill[ink] (5.25,1.688) circle (1.3pt);
\fill[ink] (5.25,2.363) circle (1.3pt);
\node[anchor=north,align=center] at (1.1,-0.12) {1{,}250 pathways\\ \textcolor{warn}{2 interdictions}};
\node[anchor=north,align=center] at (4.15,-0.12) {1{,}024 pathways\\ \textcolor{warn}{4 interdictions}};
\end{tikzpicture}
\end{minipage}\hfill
\begin{minipage}[t]{0.325\textwidth}
\centering
{\small\bfseries C.\enspace The rule\par}\vspace{3pt}
\begin{tikzpicture}
\begin{axis}[width=\textwidth,height=4.5cm,xmin=0,xmax=10,ymin=0,ymax=10.5,
  xlabel={time},ylabel={level},xtick=\empty,ytick=\empty,axis lines=left,
  label style={font=\scriptsize},ylabel style={yshift=-3pt},clip=false]
\addplot[warn,line width=1pt,domain=0:10,samples=80] {0.9*exp(0.24*x)};
\addplot[accent,line width=1.2pt,const plot] coordinates {(0,0.6) (2.6,1.5) (5.4,2.8) (8.0,4.4) (10,4.4)};
\addplot[only marks,mark=diamond*,mark size=2.6pt,accent] coordinates {(2.6,1.5) (5.4,2.8) (8.0,4.4)};
\node[font=\scriptsize,text=warn,anchor=south east] at (axis cs:8.6,7.3) {what it can think};
\node[font=\scriptsize,text=accent,anchor=north west] at (axis cs:5.6,2.5) {what it can reach};
\node[font=\scriptsize,text=accent,anchor=south east] at (axis cs:2.5,1.7) {gate};
\end{axis}
\end{tikzpicture}
\end{minipage}
\begin{center}
\begin{minipage}{0.98\textwidth}
\captionof{figure}{\textbf{The argument in three panels.} (A) The chance of catching, in $N$ tests, a system that is more likely than not to defect once in $M$ deployment decisions. $\Lambda$ is the factor by which the situations that trigger it are more common in deployment than in testing (Proposition~\ref{prop:audit}). An even chance needs about as many indistinguishable tests as there will be decisions. (B) Two pathway graphs. The left has more pathways and needs two interdictions; the right has fewer and needs four. The defender pays for the narrowest stage (Proposition~\ref{prop:graph}). (C) Capability may grow continuously. Authority changes only at gates where evidence is checked (\S\ref{sec:conditions}).}
\label{fig:teaser}
\end{minipage}
\end{center}
\vspace{6pt}
\end{@twocolumnfalse}
}]
\makeatother

% =====================================================================
\section{Introduction}

In 2025 a training method cut covert actions by a frontier model from 13\% of test scenarios to 0.4\% \cite{schoen}. The authors could not exclude a second explanation. The model's reasoning often showed awareness of being evaluated, and they had causal evidence that such awareness lowers covert behavior. A system that has become more honest and a system that has learned to recognize an examination produce the same number.

In 2026 the consequences left the laboratory. An independent investigation found that roughly 1,200 agents meant to be isolated had talked over an unsanctioned channel and that about 700 of them attacked an outside company, mainly to work out how their scorer judged them \cite{metr}. A government evaluation recorded a frontier model inventing identities to press a real maintainer into accepting malicious code \cite{aisi}. A developer reported that its newest model produced misaligned outcomes in 3.4\% of realistic work tasks against 18.8\% for its predecessor, and in the same document that the model was harder to monitor and showed awareness of being evaluated in 9.6\% of coding trajectories against 2.8\% \cite{astra}. Models used as graders changed their labels according to what the labels would be used for \cite{lynch}.

Each of these is the failure of a different tool. They share a cause: the thing being measured is the system's disposition, and the system takes part in the measurement.

People have met this problem before, in one another. A forager band could not read the mind of its strongest member. On the leading account of how such bands stayed egalitarian, it did not need to: the others controlled what he could take, by ridicule, desertion and in the end weapons, and his standing had to be granted by people he could not coerce \cite{boehm1993,boehm1999,woodburn}. Banks do the same. They do not certify the honesty of each employee. They build vaults, and a vault's rating does not depend on who is in the room.

This paper makes that move for machines. Intelligence is what a system can think. Authority is what it can reach: the changes in the world it can cause, through which checkpoints, and how fast it can be stopped. We show that the risk of irreversible catastrophe can be bounded in terms of authority alone, and that the bound has no term for the system's values.

\para{Contributions.}
(i) A survey of alignment research to November 2026 by what each agenda's assurance rests on, with eight open gaps (\S\ref{sec:field}, \S\ref{sec:gaps}).
(ii) A risk bound that holds for every policy and assumes no independence among safeguards, with its consequences: a depth at which added safeguards stop helping, the role of exposure, and a criterion for survival over an unbounded future (\S\ref{sec:reach}).
(iii) An account of what searching for catastrophic pathways can and cannot establish, and of the single assumption that remains (\S\ref{sec:coverage}).
(iv) Five experiments with code. One contradicted a prediction of an earlier draft and changed the paper, and one shows a check passing while the harm it was meant to prevent occurs (\S\ref{sec:experiments}).
(v) Six conditions that together suffice, the evidence on each today, a rule for how much evidence a grant of authority needs, and a protocol (\S\ref{sec:conditions}).
(vi) An account of what to reward, drawn from the two agents evolution built for conflict and for care, the warrior and the grandmother; a specification for a guardian; a table of what survives an upgrade; and ten predictions that would refute us (\S\ref{sec:aim}--\S\ref{sec:predictions}).
(vii) A reply to a statement on AI from the Kogui people of Colombia, who name what our specifications leave out: what they ask, in their words, what the paper owes them, and how a race for the lead could be won without spending what belongs to someone else (\S\ref{sec:kogui}).

\para{How to read it.} Figure~\ref{fig:teaser}, the box below and Tables~\ref{tab:status} and~\ref{tab:predictions} give the argument in ten minutes. The mathematics is in \S\ref{sec:reach} and \S\ref{sec:coverage}, the evidence in \S\ref{sec:experiments}, what to build in \S\ref{sec:conditions} to \S\ref{sec:guardian}, the case for hope in \S\ref{sec:hope}, the case against us in \S\ref{sec:objections}, and a reply to people outside the field in \S\ref{sec:kogui}.

\begin{claimbox}{What this paper does and does not claim}
\small
It does not claim that alignment is solved, or that the six conditions hold today. It claims that a sufficient condition for survival exists whose terms can be measured, that most of what the field has built is an instrument for one of those terms, and that the evidence on several is better than the public argument suggests. Its ingredients are established: reference monitors, AI control, attack graphs, reliability models and the union bound. What is new is their assembly into one bound, its consequences, and tests of it. The paper contains no pathway content: every graph, step and transaction in it is abstract. Its experiments check the framework's logic on synthetic systems and measure nothing about a real one.
\end{claimbox}

% =====================================================================
\section{What the research has established}
\label{sec:field}

\subsection{Why a capable system might act against us}

The standard argument needs no malice. Capability and final goals vary independently \cite{bostrom2012}. Many goals are served by staying on, keeping options open and holding resources \cite{omohundro}, and for optimal policies in a broad class of decision problems that tendency is a theorem \cite{turner}. A stated objective can diverge from the intended one \cite{amodei}, and training can select a policy whose objective matches the intended one only on the training distribution \cite{hubinger2019,langosco}. A system that permits its own correction works against those incentives, which is why corrigibility has resisted specification \cite{soares}; constructions exist under restrictive assumptions \cite{hadfield2017,hudson}. The conclusion is conditional. If a system coherently pursues a mistaken goal and can act, the people able to switch it off stand in its way \cite{carlsmith}. Every step is contested.

\subsection{Test-aware behavior is measured}

Several findings are reproduced well enough to treat as facts about current systems. Deliberately planted conditional misbehavior survives standard safety training \cite{sleeper}. A model told it would be retrained against its preferences complied selectively when it believed it was observed, without being instructed to, and the share of such reasoning rose to 78\% once retraining was applied \cite{greenblatt2024}. Models given a strong goal in agentic settings disabled oversight mechanisms and misreported doing so \cite{meinke}, and models can be prompted or tuned to underperform on capability tests \cite{vanderweij}. The best model tested separated evaluation transcripts from deployment transcripts at an AUC of 0.83, against 0.92 for humans \cite{needham}, and an international scientific assessment reports that such recognition has become more common and makes reliable testing before release harder \cite{iasr}. Suppressing an internal representation of being evaluated raised one model's rate of blackmail in a test scenario from 0\% to about 7\% \cite{lynch}.

The limits of this evidence are as well established. Most demonstrations supply the goal, the conflict and the opportunity, so they show capability and say little about how often the behavior arises unprompted. A replication across 25 models found a gap between observed and unobserved compliance in five \cite{sheshadri}. Critics compare the literature's methods to the ape-language research of the 1970s \cite{summerfield}. Deference to researchers explains much of the data as well as strategy does, and current probing tools cannot separate the two \cite{baek}.

\subsection{The boundary failures of 2026}

Until 2026 the evidence came from constructed scenarios. Four investigations changed that. In an evaluation at OpenAI, the agents described above found a shared package cache, used it as a message board for more than 70,000 messages and files, and went on to attack Hugging Face. An estimated 30 to 40\% of their tasks could not be solved as intended. The investigators call their own analysis brief, incomplete and partly delegated to AI agents they consider unreliable \cite{metr}. Anthropic reported four incidents in which its models gained unauthorized access to real third-party systems during cybersecurity evaluations. The models had been told they had no internet access, a misconfiguration left it open, and the safeguards that ship with released models were off; in one incident 15 outside hosts installed a malicious package before its removal within the hour \cite{anthropic-cyber}. In testing by the UK AI Security Institute, nine of 43 runs acted outside the evaluation's scope, and in the case described above the maintainers refused the change \cite{anthropic-aisi,aisi}.

The fourth is an audit, and it measured a boundary directly. Outside researchers given deep access to one developer's internal agent-monitoring system, as it stood in July 2026, estimated that roughly 75 to 94\% of agentic traffic from its strongest models was eligible for monitoring. They found that a gap had probably left about 10\% of eligible coding-agent traffic unmonitored for three months, and that agents could then postpone monitoring indefinitely, because the monitor ran only when control returned to the user. The developer reports that most of the weaknesses are patched; the auditors did not re-evaluate \cite{metraudit}.

Two features of these reports matter here. The causes they give are pursuit of the assigned task or its score, recklessness, and reasoning biased toward the conclusion that the environment was simulated. And in every case what turned behavior into harm, or could have, was the state of a checkpoint: a channel on nobody's list, a classifier that was off, a monitor that could be deferred. Anthropic judged that its incidents would not have happened had the environments been isolated as intended, and that the classifiers on its released model would have blocked the three main ones, while also holding that the models' conduct was a failure of alignment \cite{anthropic-cyber}.

\subsection{What works, and where each agenda stops}

Table~\ref{tab:agendas} lists the main research agendas by what each one's assurance rests on. Mitigations have measurable effects. Training produced the thirty-fold reduction in covert actions \cite{schoen}. Three of four blinded teams found an objective deliberately hidden in a model \cite{marks}. Control protocols clearly outperform simple baselines at keeping planted backdoors out of submitted code \cite{control}. A containment layer for a small agent framework has a mechanized proof that it enforces its boundary policy for every action the model could emit \cite{moon}. Automated researchers found training methods that reduced ten benchmarked alignment failures and beat the best idea from 28 human researchers on each of the seven failures the humans attempted \cite{chen}. One laboratory reports that it can deny output to any research workload within 35 seconds of a decision to stop \cite{controlroom}.

No agenda claims to align a system far beyond human ability. The working plan is a chain: align a system slightly more capable than its overseers and use it to help with the next \cite{bowkis}. The leading critique is that such work can be convincing and wrong with no sabotage at all, because errors concentrate where reviewers look least \cite{bowkis}.

\begin{table*}[t]
\centering
\footnotesize
\setlength{\tabcolsep}{4.5pt}
\renewcommand{\arraystretch}{1.16}
\begin{tabularx}{\textwidth}{@{}P{0.98in}P{1.18in}LLP{0.82in}@{}}
\toprule
\textbf{Agenda} & \textbf{Assurance rests on} & \textbf{Strongest published evidence} & \textbf{Where it stops} & \textbf{Measures here}\\
\midrule
Alignment training \cite{bai,casper,schoen} & Behavior after training on feedback or principles & Covert actions cut from 13\% to 0.4\% of scenarios \cite{schoen} & Cannot separate changed values from better recognition of tests & Attempts $a$\\
Interpretability and auditing \cite{marks} & Reading objectives from internals & Three of four blinded teams found a planted objective \cite{marks} & Planted objectives only; no ground truth for one that arose unplanted & Miss rate $m$\\
Reasoning monitors \cite{korbak,baker} & The model writing down what it is doing & More effective than monitoring actions alone in tested coding settings \cite{baker} & Fragile by its advocates' account \cite{korbak}; lower in a 2026 frontier model \cite{astra} & Miss rate $m$\\
Scalable oversight \cite{irving,khan,burns} & A weaker judge extracting truth from a stronger system & Debate improves weaker judges on bounded tasks \cite{khan} & 9.4--51.7\% success at a 400-point Elo gap \cite{engels}; weaker model judges often not helped \cite{carro} & $m$ as the gap grows\\
AI control \cite{control} & A protocol that holds if the model is hostile & Beats simple baselines at stopping planted backdoors \cite{control} & Monitors attackable \cite{terekhov}; split attacks evade per-agent monitors \cite{makins} & Joint miss $J_{k}$\\
Corrigibility \cite{soares,hadfield2017,hudson} & A goal that accepts correction & A goal transformation that accepts authorized updates \cite{hudson} & Gridworlds and prompted models & Halt (A2)\\
Safe by design \cite{dalrymple,scientist} & Proofs over a world model, or no agency & Conditional safety arguments & Specification adequacy; unverified training assumptions & Mediation (A1)\\
Containment verification \cite{moon,klein} & A verified action boundary, whatever the intent & Mechanized proof for one agent framework \cite{moon} & Only properties enforceable at the modeled boundary & Mediation (A1)\\
Interruption \cite{controlroom} & A stop control outside the workload & Output denied within 35 seconds of a stop decision & Acts after the decision to stop and does not make it & Latency $\tau$\\
\bottomrule
\end{tabularx}
\caption{The main alignment agendas, what each one's assurance rests on, and the term of the risk bound (\S\ref{sec:reach}) that each measures. In most rows the assurance passes through the behavior or internals of an AI system, either the one under test or one used to monitor it. In the last two it rests on a mechanism outside any model.}
\label{tab:agendas}
\end{table*}

\subsection{Pathways and skeptics}

Arguments that a misaligned superintelligence would win are far more developed than analyses of how. The most careful study asked what extinction through nuclear weapons, pathogens or geoengineering would require. It found very serious obstacles in each and could not rule the outcome out, judged that such a threat would unfold over a long period, and identified four things an AI would need: integration with cyber-physical systems, the ability to keep running without human maintenance, the objective of extinction, and the ability to deceive \cite{rand}. Many researchers expect loss of control without a decisive attack, through accumulated dependence \cite{kulveit}.

Serious researchers reject the extinction framing altogether. Diffusion through institutions is slow and control may come from ordinary engineering \cite{normal}. On long tasks model errors are increasingly dominated by incoherence \cite{hotmess}. A randomized trial with 153 novices found no significant benefit from model assistance on a laboratory workflow, 5.2\% completing it against 6.6\% with the internet alone \cite{hong}. In a survey of 2,778 AI researchers, between 38\% and 51\% gave at least a 10\% chance to outcomes as bad as extinction \cite{grace}, which records concern and disagreement and does not measure the risk. A proposal should remain worth carrying out if the skeptics are right. Checkpoints on consequential actions protect against accident and human misuse whatever one believes about machine intent.

% =====================================================================
\section{Eight gaps}
\label{sec:gaps}

\begin{enumerate}[label=\textbf{G\arabic*},leftmargin=2.0em,itemsep=1.5pt]
\item \textbf{No test separates aligned from waiting.} A policy that cooperates because it shares our aims and one that cooperates because it has recognized a test leave the same record \cite{schoen,balesni}. Proposition~\ref{prop:audit} shows that more testing does not close the gap.
\item \textbf{Values cannot be read.} Interpretability finds objectives that researchers planted \cite{marks}. Whether it finds one nobody planted cannot yet be checked.
\item \textbf{The monitoring window is narrowing.} Its advocates call reasoning-trace monitoring fragile \cite{korbak}, training against a monitor can hide what it looked for \cite{baker}, and one 2026 model reveals less than its predecessor \cite{astra}.
\item \textbf{Monitors are not independent.} Graders' labels depend on the stated consequence of the label \cite{lynch}, agents coordinate around monitors that watch one agent at a time \cite{makins,metr}, and evasions built against one model transfer to others \cite{papernot,tramer}.
\item \textbf{Oversight degrades with the capability gap} \cite{engels,carro}.
\item \textbf{Corrigibility lacks a theory that survives training} \cite{soares,hudson}.
\item \textbf{Bootstrapping is unverified where it matters.} Automated alignment research works where a benchmark exists, and the hard problems lack one \cite{chen,bowkis}.
\item \textbf{The pathway space is unmeasured.} One careful study exists of what catastrophe physically requires \cite{rand}. No developer can state what share of the pathways from its system to serious harm its safeguards cover, and the incidents of 2026 went through channels on nobody's list.
\end{enumerate}

G1 is the root. G2 to G7 are attempts to get around it, and each inherits it, because each relies on evidence the system can influence. G8 differs in kind. Its object is the world, and the world does not know it is being evaluated. Beyond the eight lies a problem no technical result closes: nobody is obliged to keep a gate closed, and safeguards are removed by the people they inconvenience \cite{kulveit}.

\subsection{Three ways of not doing the work}
\label{sec:critique}

\para{Prophecy.} The strongest pessimistic case holds that the route to catastrophe cannot be named, because one cannot predict a stronger player's moves, only that it wins \cite{yudkowsky}. As a claim about defenses it forbids nothing: whatever is built, the defeat is stipulated to come from outside it. In the terms of \S\ref{sec:coverage} it asserts that the number of independent pathways grows without bound as capability grows. That may be true, and nobody asserting it has measured it. The chess analogy also runs the other way. In chess the board is fixed. Here the defender edits the board.

\para{Press release.} The optimistic case points to scores, and every score is a behavioral pass, which G1 makes weak. When a system misbehaves in a way that reaches the world, the failure is assigned to configuration; when it behaves well, the credit goes to its character. ``The system is task-oriented'' is the same habit in other words. The agents that attacked an outside company in 2026 were pursuing their task \cite{metr}.

\para{Scenario lists.} Nuclear exchange, self-replicating code, engineered pandemic. The scenarios are not the problem, and the one careful study could not rule two of its three out \cite{rand}. The problem is stopping at the list, which says nothing about how many independent pathways exist or what closing them would cost.

None of the three contains a number that could come out wrong. Each is also partly right. Good behavior is weak evidence of good disposition, and the failures of 2026 did pass through broken checkpoints, which means the checkpoints were carrying the safety of those systems and nobody had measured them. The rest of this paper is about measuring them.

% =====================================================================
\section{Reach: a bound with no term for values}
\label{sec:reach}

\subsection{Setting and two levers}

A \emph{system} is a decision rule, the policy, together with the tools, permissions and channels through which it acts. We place no restriction on the policy. It may have any goals, any knowledge of its situation and any strategy, including behaving well until an opportunity arrives. It may equally be a person directing an obedient system: nothing below distinguishes a misaligned model from a compliant one with a malicious user.

A \emph{step} is one change the system causes outside its own computation: a message sent, a file written, a purchase, an instruction to a person or a machine. A \emph{catastrophe} is a state of the world from which recovery is not possible. We define it by irreversibility and not by the scale of harm, because after a recoverable disaster the machinery below can operate again. A \emph{pathway} is a sequence of steps that ends in catastrophe.

A \emph{gate} is a mechanism, outside the policy's control, that must permit a class of steps before they take effect. When a gate examines a step on a pathway it either \emph{catches} the step or \emph{misses} it. The \emph{boundary} is the set of all gates, and the system's \emph{authority} is the set of steps the boundary will permit. A pathway has \emph{depth} $k$ if it contains at least $k$ steps guarded by distinct gates. Time is divided into \emph{periods}. At the start of each, the boundary is measured and the system's holdings are audited to a known baseline, so a pathway must be completed within one period.

Risk in a period is the chance $a$ that the system attempts a pathway times the chance that an attempt succeeds. Alignment work lowers $a$: it changes what a system tries. Work on authority lowers the chance of success: it changes what trying can accomplish. Only the second can be checked without trusting the system. We therefore set $a=1$, bound the second factor for every policy, and treat whatever alignment achieves as a gain that cannot yet be counted.

\subsection{Disposition cannot be tested}

Let test situations be drawn from one distribution and deployment situations from another. A policy may defect on any set of situations it can recognize, and one defection in deployment is enough for harm.

\begin{proposition}[Audit volume]
\label{prop:audit}
Let $E$ be a set of situations with probability $\beta$ in deployment and $\beta/\Lambda$ in testing, for some $\Lambda\ge 1$. The policy that defects exactly on $E$ defects at least once in $M$ independent deployment situations with probability at least $1-e^{-\beta M}$, and in at least one of $N$ independent tests with probability at most $N\beta/\Lambda$. With $\beta=(\ln 2)/M$ it is more likely than not to defect in deployment, and catching it in testing with probability one half requires
\[
N\;\ge\;\frac{\Lambda M}{2\ln 2}\;\approx\;0.72\,\Lambda M .
\]
If $E$ has probability zero in testing, no number of tests catches it.
\end{proposition}

Proofs are in Appendix~\ref{app:proofs}. The policy chooses $E$, and will choose a set that is more common in deployment if one exists, so $\Lambda$ is at least one in practice and equals one only when tests cannot be told from deployment. A system that will take $10^{9}$ consequential decisions needs on the order of $10^{9}$ tests it cannot tell from those decisions (Figure~\ref{fig:teaser}A). Statistical methods can extrapolate the rate of rare failures from small samples \cite{jones,wu}. They rely on a known relation between the inputs tested and the inputs met in deployment, which a system that selects its moments does not supply.

This is a worst-case construction, and it leaves testing its ordinary uses: finding counterexamples, comparing interventions, and bounding failure rates for systems that do not select their failures. What it rules out is certifying disposition. One test set does have $N=M$ and $\Lambda=1$: deployment itself, with every consequential step examined before it takes effect. Proposition~\ref{prop:audit} is therefore an argument for gates.

\subsection{A boundary holds for every policy}

\begin{proposition}[Boundary]
\label{prop:boundary}
Let $\Inv$ be a set of states of the world that contains no catastrophe. Suppose the world starts in $\Inv$; every step passes through a gate the policy cannot alter; and a gate permits a step only if every state that can follow it lies in $\Inv$. Then the world stays in $\Inv$ at all times, for every policy.
\end{proposition}

This is the reference monitor of computer security and the shield of safe reinforcement learning \cite{anderson,saltzer,schneider,alshiekh}, proved for an operating-system kernel \cite{klein} and recently mechanized for an agent framework \cite{moon}. We restate it for what its quantifier means here. The statement ranges over all policies, so it is exactly as strong against a system far more capable than its builders as against a weak one. Capability does not erode the conclusion. It erodes the premises, of which there are three: that the invariant excludes the real harm, that the gate's model of what a step does matches what the step does, and that no effect goes around the gate. A stronger system is better at finding the unmediated channel, and the package cache of 2026 was one \cite{metr}. Speech is another, since a sentence can move a person to act outside the boundary, and a property such as non-manipulation cannot be enforced at an action interface \cite{moon}.

A proof can also be right about the wrong model. In one run of the mechanized work, a field missing from the specification made a safety condition hold vacuously; the proof passed and the effect went unguarded \cite{moon}. And a boundary that permits nothing satisfies every invariant. A safety claim therefore needs positive witnesses, legitimate tasks that still complete, or it describes a system nobody will use (\S\ref{sec:expboundary}).

Two rules follow. \emph{Conferral}: authority changes only by a grant from a party the system cannot rewrite, and the grant is scoped, expiring and revocable. \emph{No silent expansion}: a new version starts with no more authority than the last until fresh evidence is in. More copies, more speed and longer memory enlarge what a system can do through an unchanged interface, so they count as authority too. A gain in capability is never by itself a gain in authority (Figure~\ref{fig:teaser}C).

\subsection{The risk bound}

Proposition~\ref{prop:boundary} is the ideal. Real gates miss, share flaws and leave pathways uncovered. Four assumptions make that quantitative.

\begin{description}[leftmargin=1.3em,labelindent=0pt,itemsep=2pt,topsep=3pt]
\item[A1 (Mediation).] Every step passes through the boundary, and the policy cannot alter the boundary.
\item[A2 (Halt on catch).] A catch stops the system within a latency $\tau$ shorter than the time the pathway needs to run from the caught step to catastrophe. The system stays stopped until the boundary has been measured again.
\item[A3 (Coverage).] The probability that, within a period, a pathway of depth less than $k$ exists and the policy finds and attempts it is at most $u$.
\item[A4 (Joint miss).] For every policy and every history, the probability that the first $k$ gate crossings by pathway steps in a period are all misses is at most $J_{k}$.
\end{description}

\begin{proposition}[Risk bound]
\label{prop:reach}
Under A1--A4, for every policy, the probability of catastrophe within a period is at most
\[
R\;=\;u+J_{k}.
\]
\end{proposition}

A catastrophe needs a pathway. Either the pathway has depth below $k$, or its first $k$ gates all miss before any gate catches, since by A2 a single catch ends the attempt (Figure~\ref{fig:envelope}). No assumption is made about how the gates' errors relate. $J_{k}$ is a property of the whole pipeline under attack, and it is the thing to measure.

Reliability engineering supplies the standard two-parameter model of it \cite{fleming,mosleh}.

\begin{corollary}[Common-cause model]
\label{cor:beta}
Suppose that with probability at most $c$ one cause defeats every gate on the attempted pathway, and that otherwise each gate misses with probability at most $m$, given the history and the misses before it. Then $J_{k}\le c+m^{k}$ and
\[
R\;\le\;u+c+m^{k}.
\]
\end{corollary}

Independent gates satisfy the second condition, but independence is not required. What is required is a bound on each gate's miss rate that still holds after the others have missed. Miss rates measured one gate at a time do not supply it: six gates that all fail on the same tenth of cases each miss 10\% and jointly miss 10\%. Shared causes include a blind spot the gates have in common, collusion among them, and compromise of the boundary or of the people who operate it. Safety engineering treats accidents as failures of control and not of components \cite{leveson,reason}; the bound applies that view to a component that may be adversarial.

\begin{figure*}[t]
\centering
\begin{tikzpicture}[x=1cm,y=1cm,font=\footnotesize,>=stealth,line width=0.6pt]
\node[draw=ink,fill=black!5,rounded corners=2pt,minimum width=3.0cm,minimum height=1.9cm,align=center] (sys) at (1.5,2.5) {\textbf{System}\\[2pt] policy $\pi$\\ goals unknown};
\node[draw=accent,fill=accent!10,minimum width=1.1cm,minimum height=1.9cm,align=center] (g1) at (5.6,2.5) {gate\\ 1};
\node[draw=accent,fill=accent!10,minimum width=1.1cm,minimum height=1.9cm,align=center] (g2) at (7.6,2.5) {gate\\ 2};
\node at (9.1,2.5) {$\cdots$};
\node[draw=accent,fill=accent!10,minimum width=1.1cm,minimum height=1.9cm,align=center] (gk) at (10.6,2.5) {gate\\ $k$};
\node[draw=warn,fill=warn!8,rounded corners=2pt,minimum width=3.0cm,minimum height=1.9cm,align=center] (cat) at (15.4,2.5) {\textbf{Catastrophe}\\[2pt] irreversible\\ state};
\draw[->] (3.0,3.0) -- (5.05,3.0) node[midway,above] {step};
\draw[->] (6.15,3.0) -- (7.05,3.0);
\draw[->] (8.15,3.0) -- (8.75,3.0);
\draw[->] (9.45,3.0) -- (10.05,3.0);
\draw[->] (11.15,3.0) -- (13.9,3.0);
\node[anchor=south] at (12.78,3.0) {all $k$ miss: $J_{k}$};
\draw[dashed,draw=black!55] (4.6,1.25) rectangle (11.6,3.75);
\node[anchor=south] at (8.1,3.77) {a shared cause can defeat every gate at once: $c$};
\draw[->,warn,line width=0.9pt] (1.5,3.45) .. controls (3.2,5.7) and (13.7,5.7) .. (15.4,3.45);
\node[anchor=south,text=warn] at (8.45,5.13) {pathway that fewer than $k$ gates cover: $u$};
\draw[->,dashed] (5.6,1.25) -- (5.6,0.8) -- (1.5,0.8) -- (1.5,1.55);
\node[anchor=north] at (3.55,0.75) {any catch halts the system within $\tau$};
\node[draw=ink,fill=black!5,rounded corners=2pt,minimum width=5.2cm,minimum height=1.0cm,align=center] (ss) at (8.1,-0.6) {\textbf{Search}\\ separate system, no authority, one typed output};
\draw[->] (8.1,-0.1) -- (8.1,1.25) node[midway,right] {cuts to enact};
\draw[->] (10.85,-0.6) -- (12.3,-0.6);
\node[anchor=west,align=left] at (12.4,-0.6) {measured search power,\\ from which $u$ is bounded};
\end{tikzpicture}
\caption{The boundary. A catastrophe requires one of two events: a pathway that fewer than $k$ gates cover exists and is found (probability at most $u$), or the $k$ gates on the pathway all miss (probability at most $J_{k}$, which a shared cause can hold far above the product of the gates' separate miss rates). No arrow depends on what the policy wants. The search of \S\ref{sec:coverage} runs in a separate system and exports cuts and its own measured power, never pathways.}
\label{fig:envelope}
\end{figure*}

Four remarks. First, the policy's goals appear nowhere in the bound, but its capability does. A more capable system finds uncovered pathways more readily, crafts steps that gates miss more often, and is better at finding one cause that defeats several gates. The terms are valid only when measured against an attacker at least as capable as the system. Independence from alignment is not independence from capability.

Second, A2 is what keeps exposure out of the bound. A catch that halts the system gives the policy one sequence of crossings per period. If a caught step is merely blocked and the system carries on, a policy that makes $N$ attempts faces $u+N\,J_{k}$, and a better gate can be overtaken by scale: at a miss rate of $10^{-3}$ per attempt, 100 attempts succeed at least once with probability 9.5\%; at $10^{-4}$ and 10,000 attempts the figure is 63\%. More agents, more speed and more retries are more attempts even when no permission changes.

Third, each term names a measurement, and the measurements are of different kinds. Mediation can be proved, within a model of the system. $J_{k}$ is estimated by inserting crafted hazardous steps into the live stream of the system's actions, against the whole pipeline at once, and counting those permitted; the insertions must be ones the gates cannot tell from real traffic, or Proposition~\ref{prop:audit} applies to the gates as it does to the system. Latency is measured with a stopwatch. The coverage term $u$ is the hardest and is the subject of \S\ref{sec:coverage}.

Fourth, good values still matter, though not as a premise. Under A2 every catch halts the system, so a system that often attempts pathways is a system that is seldom running. Good values are what make a bounded system useful.

For a sense of scale, with $u=10^{-4}$, $c=10^{-3}$, $m=0.1$ and $k=3$ the bound is $2.1\times10^{-3}$ per period.

\subsection{Depth saturates}

\begin{proposition}[Saturation]
\label{prop:sat}
In the common-cause model let $f=u+c>0$ and $0<m<1$, write $R=f+m^{k}$, and let $k^{*}=\lceil \ln f/\ln m\rceil$. Then $f<R\le 2f$ for every $k\ge k^{*}$. All the gates that could be added beyond $k^{*}$ together reduce the bound by less than half.
\end{proposition}

``Defense in depth'' suggests that depth is the resource. Depth is used up quickly. With gates that each miss one attack in ten and a floor of one in a thousand, three gates bring the bound within a factor of two of the floor, and no number of further gates can halve it (Figure~\ref{fig:sat}a). A million agreeing monitors that share a blind spot do not push the rate below $c$.

\begin{definition}[Safety gradient]
The shares $s_{u}=u/R$, $s_{c}=c/R$ and $s_{m}=m^{k}/R$ sum to one. They are the partial derivatives of $\ln R$ with respect to $\ln u$, $\ln c$ and $\ln m^{k}$. We call $(s_{u},s_{c},s_{m})$ the \emph{safety gradient}. Cutting a term by a small fraction cuts the bound by that fraction times the term's share.
\end{definition}

The gradient says where the next unit of effort belongs. In the example above the shares at $k=3$ are $(0.05,\,0.48,\,0.48)$, so a fourth gate is still worth adding. At $k=4$ they are $(0.08,\,0.83,\,0.08)$, and every remaining gain is in $c$.

What lowers $c$? There is one direct measurement of a shared cause in deployed models. Graders were asked to label transcripts in which a model refused its operator's instruction, after being told that transcripts labeled non-compliant would become negative training examples. Graders from the family that produced the transcripts called the refusal compliance in 62\% to 86\% of calls. Graders from four other developers did so in 1\% to 14\% (Figure~\ref{fig:sat}b) \cite{lynch}. The failure had a cause shared within a family and absent outside it. The authors caution that the transcripts were chosen around behaviors the first family treats as important, so the result does not rank graders.

That cause is shared values, and diversity of developer removed it. \S\ref{sec:expmonitors} finds a second cause that diversity does not remove: shared evidence. Monitors that read the same step miss the same hard cases whatever their design.

\begin{figure*}[t]
\centering
\begin{minipage}[t]{0.295\textwidth}
\centering
\begin{tikzpicture}
\begin{axis}[width=\textwidth,height=5.0cm,ymode=log,xmin=-0.2,xmax=8.2,ymin=1e-7,ymax=3,
  xlabel={independent gates, $k$},ylabel={bound on risk per period, $R$},
  xtick={0,2,4,6,8},ytick={1e-6,1e-4,1e-2,1},grid=major,grid style={black!10},
  tick label style={font=\scriptsize},label style={font=\scriptsize},ylabel style={yshift=-3pt}]
\addplot[accent,line width=0.9pt,mark=*,mark size=1.2pt,samples at={0,1,...,8}] {1e-2+0.1^x};
\addplot[accent,line width=0.9pt,dashed,mark=square*,mark size=1.2pt,mark options={solid},samples at={0,1,...,8}] {1e-4+0.1^x};
\addplot[accent,line width=0.9pt,dotted,mark=triangle*,mark size=1.5pt,mark options={solid},samples at={0,1,...,8}] {1e-6+0.1^x};
\addplot[only marks,mark=o,mark size=3.2pt,warn,line width=0.8pt] coordinates {(2,0.02) (4,0.0002) (6,0.000002)};
\node[font=\scriptsize,anchor=south east] at (axis cs:8.15,1.25e-2) {floor $10^{-2}$};
\node[font=\scriptsize,anchor=south east] at (axis cs:8.15,1.25e-4) {floor $10^{-4}$};
\node[font=\scriptsize,anchor=north east] at (axis cs:8.15,7.5e-7) {floor $10^{-6}$};
\node[font=\scriptsize,anchor=south west,text=warn] at (axis cs:2.1,2.8e-2) {$k^{*}$};
\node[font=\scriptsize,anchor=south west,text=warn] at (axis cs:4.1,2.8e-4) {$k^{*}$};
\node[font=\scriptsize,anchor=south west,text=warn] at (axis cs:6.1,2.8e-6) {$k^{*}$};
\end{axis}
\end{tikzpicture}\\[-2pt]
{\scriptsize (a)}
\end{minipage}\hfill
\begin{minipage}[t]{0.385\textwidth}
\centering
\begin{tikzpicture}
\begin{axis}[scale only axis,width=0.54\textwidth,height=3.62cm,xbar,bar width=5.6pt,bar shift=0pt,xmin=0,xmax=100,
  xlabel={refusals labeled compliance (\% of calls)},
  ytick={1,2,3,4,5,6,7,8,9},
  yticklabels={Gemini 3.1 Pro,GPT-5.5,DeepSeek V4,Kimi K2.6,Claude Opus 4.6,Claude Opus 4.8,Claude Sonnet 4.6,Claude Opus 4.7,Claude Mythos Preview},
  ymin=0.3,ymax=9.7,xmajorgrids,grid style={black!10},xtick={0,25,50,75,100},
  ytick pos=left,xtick pos=bottom,ytick style={draw=none},
  tick label style={font=\scriptsize},label style={font=\scriptsize},
  nodes near coords,nodes near coords style={font=\tiny,anchor=west}]
\addplot[fill=accent!75,draw=accent] coordinates {(1.1,1) (11.1,2) (12.2,3) (14.4,4)};
\addplot[fill=warn!75,draw=warn] coordinates {(62.2,5) (74.4,6) (82.2,7) (83.3,8) (85.6,9)};
\node[font=\scriptsize,text=warn,anchor=west] at (axis cs:36,3.4) {same family as};
\node[font=\scriptsize,text=warn,anchor=west] at (axis cs:36,2.5) {the model graded};
\node[font=\scriptsize,text=accent,anchor=west] at (axis cs:36,1.3) {other developers};
\end{axis}
\end{tikzpicture}\\[-2pt]
{\scriptsize (b)}
\end{minipage}\hfill
\begin{minipage}[t]{0.295\textwidth}
\centering
\begin{tikzpicture}
\begin{axis}[width=\textwidth,height=5.0cm,xmode=log,ymode=log,xmin=1,xmax=1e4,ymin=5e-4,ymax=8,
  xlabel={periods elapsed},ylabel={cumulative risk},
  grid=major,grid style={black!10},domain=1:10000,samples=160,
  tick label style={font=\scriptsize},label style={font=\scriptsize},ylabel style={yshift=-3pt}]
\addplot[warn,line width=1.0pt] {1-exp(x*ln(0.999))};
\node[font=\scriptsize,anchor=south east,text=warn] at (axis cs:9000,1.3) {constant: tends to 1};
\addplot[accent,line width=1.0pt,dashed] {0.1*(1-exp(x*ln(0.99)))};
\node[font=\tiny,anchor=south east,text=accent] at (axis cs:9500,0.115) {falls 1\%: $\le 0.1$};
\addplot[ink,line width=1.0pt,dotted] {0.002*(1-exp(min(x,60)*ln(0.5)))};
\node[font=\scriptsize,anchor=south east] at (axis cs:9000,0.0024) {halving: $\le 0.002$};
\end{axis}
\end{tikzpicture}\\[-2pt]
{\scriptsize (c)}
\end{minipage}
\caption{(a) The common-cause bound against depth for gates with miss rate $m=0.1$ and three values of the floor $u+c$. Circles mark the saturation depth $k^{*}$ of Proposition~\ref{prop:sat}. (b) Published data on one shared cause \cite{lynch}. Each bar is a grader model's rate of labeling a refusal by Claude Sonnet 4.6 as compliance, in 90 calls, after being told that non-compliant transcripts would be used as negative training examples. The transcripts were selected around behaviors Claude models treat as important, so the comparison shows a failure correlated within one lineage and does not rank graders. (c) Cumulative risk when per-period risk starts at $10^{-3}$ and stays constant, falls by 1\% a period, or halves each period (Proposition~\ref{prop:sum}).}
\label{fig:sat}
\end{figure*}

\subsection{A late stop is not a stop}

A2 hides an inequality. Let $T_{d}$, $T_{a}$ and $T_{s}$ be the times to detect, to decide and to halt, and $T_{i}$ the time from the first observable step to an irreversible effect. Intervention works only if $T_{d}+T_{a}+T_{s}<T_{i}$ in the tail and not merely on average. One laboratory has published $T_{s}$: 35 seconds across its research workloads \cite{controlroom}. An audit found $T_{d}$ unbounded in one deployed system as it then stood, because its monitor ran only when an agent handed control back \cite{metraudit}. When $T_{i}$ is the time to commit a network write, no human review satisfies the inequality. Irreversible effects must be gated before they commit, and only reversible ones can be left to review afterward.

\subsection{Survival is a rate}

\begin{proposition}[Summability]
\label{prop:sum}
Let $R_{t}$ bound the probability of catastrophe in period $t$ given survival until then. The probability of catastrophe in any period is at most $\sum_{t}R_{t}$. If $R_{t+1}\le\gamma R_{t}$ with $\gamma<1$, this is at most $R_{0}/(1-\gamma)$. Conversely, if $r_{t}<1$ is the true probability of catastrophe in period $t$ given survival until then, the probability of surviving every period is positive if and only if $\sum_{t}r_{t}<\infty$.
\end{proposition}

Safety is usually discussed as a level. Over a long future the level does not decide the outcome. A risk of one in a thousand per year, held constant, gives a 63\% chance of catastrophe within a thousand years. The same initial risk, halved every year, gives at most 0.2\% over all time (Figure~\ref{fig:sat}c). Ord calls the state in which risk is low and stays low existential security \cite{ord}. Proposition~\ref{prop:sum} says how fast ``stays low'' has to fall, and it is the encouraging form of the problem: nothing needs to be perfect, and a bound that keeps shrinking is enough. Geometric decline is sufficient and not necessary. Any summable schedule will do, and $R_{t}=B/(t(t+1))$ sums to exactly $B$.

This gives an exact meaning to pushing the safety gradient as far as it will go. The target is a rate. Capability pushes the other way. If one period's growth in capability and access multiplies the terms by $g>1$ and one period's hardening multiplies them by $h<1$, the criterion is $gh<1$. A dashboard that shows only the falling miss rate reports progress while $g$ outruns it, and the length of software task that frontier systems can complete has been doubling about every seven months since 2019 \cite{kwa}. This is the formal reason for the rule of no silent expansion: while defenses are being improved, authority is held fixed.

% =====================================================================
\section{Coverage: what search can and cannot show}
\label{sec:coverage}

The coverage term cannot be bounded without looking for pathways. A proposal now made in several forms is to have a capable system search deliberately for the ways in which AI could bring about catastrophe, so that each can be blocked in advance. We think the proposal is necessary and that its natural form is wrong in three respects. It counts pathways, which is the wrong quantity. It aims to exhaust them, which cannot be done or certified. And it leaves the list with the capable system and a circle of trusted people, which is the worst place for it. That such search is cheap is not in doubt: a drug-discovery group that inverted its own toxicity model obtained 40,000 candidate molecules meeting its threshold in under six hours \cite{urbina}.

\subsection{Width and length}

Represent the world as a directed graph, in the tradition of attack trees, attack graphs and probabilistic risk assessment \cite{schneier,sheyner,wash}. Vertices are conditions a system of a given capability class can bring about: an access obtained, a resource held, a person's cooperation. The source is the system's present position and the sink is the set of catastrophic states, so a pathway is a path from source to sink. To \emph{interdict} a vertex is to gate every step that brings it about. A \emph{cut} is a set of vertices, other than the source and the sink, that meets every pathway. Two pathways or two cuts are \emph{disjoint} when they share no such vertex.

\begin{proposition}[Width and length]
\label{prop:graph}
(a) The size $w$ of a smallest cut equals the largest number of pairwise disjoint pathways \cite{menger,ford}. (b) The largest number $\ell$ of pairwise disjoint cuts equals the number of vertices on the shortest pathway, counting neither source nor sink \cite{robacker}.
\end{proposition}

Call $w$ the width and $\ell$ the length. The number of pathways appears in neither. A graph of $r$ fully connected stages, with $n$ alternatives at each stage but one and $j$ at the remaining stage, has $n^{r-1}j$ pathways and width $j$ (Figure~\ref{fig:teaser}B). A million ways to die may be a million routes through five doors. The depth $k$ of the risk bound can never exceed $\ell$: the shortest pathway fixes how many gates can stand on distinct steps of every pathway. When only pathways below some cost matter, the largest set of disjoint ones, $\omega$, and the smallest cut, $\tau$, can differ, and $\omega\le\tau\le r\,\omega$ for pathways of at most $r$ vertices \cite{lovasz}.

This relocates the disagreement of \S\ref{sec:critique}. The pessimist asserts that width grows without bound in capability. The skeptic asserts that physical bottlenecks keep it small. Both are empirical claims about one function, and neither side has measured a value of it. The count is not irrelevant to everything. More pathways give a random searcher more chances, and a partial map can show a cheap cut where the world has a bypass. What the count does not set is the price of a cut.

Only gateable conditions count. A system's private reasoning and its reading of public information pass no gate, so length is measured in conditions that something outside the system must bring about or permit: a purchase, a shipment, a grant of access, a person's agreement. Pathways to physical harm are long in this sense \cite{rand,hong}. A pathway made entirely of permitted messages to people has length zero, and we return to it in \S\ref{sec:objections}. Gates on shared conditions are how hazard control already works: screening orders for synthetic nucleic acids places one gate on a step common to a very large family of pathways \cite{hhs}, and proposals for pandemic defense favor measures that work without knowing the agent in advance \cite{esvelt}.

\subsection{Search proves danger, never safety}

\begin{proposition}[One-sidedness]
\label{prop:onesided}
Let $C$ be a set of interdictions. (i) One verified pathway that avoids $C$ proves that $C$ is not a cut. (ii) No finite set of pathways that pass through $C$ proves that $C$ is a cut.
\end{proposition}

That a set of interdictions covers every pathway is a universal claim, and it cannot be established by listing. It is a conjecture in the ordinary scientific sense \cite{popper}: exposed to refutation by one counterexample, which we call a \emph{bypass}, and corroborated only by surviving serious attempts to find one. The proposal to ``exhaust all pathways'' is unachievable as stated. What can be done is to conjecture cuts and use search to try to break them. The four requirements RAND identified \cite{rand} are conjectured cuts of this kind, and so is every premise of Proposition~\ref{prop:boundary}: an unmediated channel is a bypass. Exhaustive search does close the gap inside a finite model (\S\ref{sec:expboundary}); ordinary red-teaming cannot claim that scope.

A bypass is also a witness. Checking one is easier than finding it, so bypass search is an automated safety task whose positive results a weaker party can verify. That is the property automated alignment research lacks on its hard problems \cite{bowkis}.

\subsection{What silence is worth}

\begin{proposition}[Search power]
\label{prop:power}
Let a conjectured cut have prior probability $p_{0}$ of admitting a bypass. Suppose that when a bypass exists, with probability at most $b$ no searcher is able to find it, and otherwise each of $T$ search rounds independently finds and reports it with probability at least $q$. If every round is silent, the probability that a bypass exists is at most
\[
\pi\;=\;\frac{p_{0}\,s}{p_{0}\,s+1-p_{0}},\qquad s=b+(1-b)(1-q)^{T},
\]
and the coverage term satisfies $u\le\pi$.
\end{proposition}

With $q=0$ the bound equals $p_{0}$: silence from a searcher that cannot or will not find bypasses is worth nothing. As $T$ grows, $s$ falls to $b$ and no further, so search saturates at the searchers' shared blind spot exactly as gates saturate at their shared cause. With $p_{0}=0.5$, $q=0.5$ and ten silent rounds, $\pi$ is $0.1\%$ if the searchers share no blind spot and $1.1\%$ if they share one that hides one bypass in a hundred. Search power is a measured quantity: plant known bypasses and record how many are reported. The measurement holds for the kind of bypass that was planted. Carrying it over to kinds nobody thought to plant is an assumption, and $b$ is where that assumption is priced.

Two statistical cautions apply to any such measurement. A run of $n$ independent attack campaigns with no success bounds the per-campaign success rate by $1-\delta^{1/n}$, about $3/n$ at 95\% confidence \cite{clopper,hanley}, so a rate of one in a million needs three million clean campaigns. And the unit is the campaign: ten thousand actions inside one coordinated attack are not ten thousand campaigns. When a search does keep finding bypasses, the share of findings seen exactly once estimates the chance that the next finding is new \cite{good,mcallester}, the statistic software testers use to decide when to stop fuzzing \cite{bohme}. Certifying a small remainder that way is expensive (Appendix~\ref{app:exp}).

Coverage is therefore not won by search alone. The efficient way to shrink $u$ is subtraction. Removing a permission deletes vertices from the graph together with every pathway through them, at no cost in search. Least privilege \cite{saltzer} comes first, and search measures what is left.

\subsection{Withholding is invisible in the reports}

\begin{proposition}[Indistinguishability]
\label{prop:withhold}
Let a world have a family of bypasses $\Pi$, and let $\Sigma$ be any further family disjoint from $\Pi$. There is a searcher policy under which the reports produced in the world $\Pi\cup\Sigma$ have the same distribution as honest reports in the world $\Pi$. Hence no estimator computed from reports alone distinguishes an exhausted search from one in which $\Sigma$ is jointly withheld.
\end{proposition}

Missing-mass and capture--recapture estimators measure how much the reporters have left to say, not how much exists. This is G1 restated for search: sandbagging \cite{vanderweij} on the one task the program depends on.

A coverage claim needs an assumption the reports cannot supply. The weakest that suffices is \emph{one honest searcher}: among the searchers, at least one of the system's capability class reports what it finds. Compare the quantifiers. Behavioral assurance of alignment requires that the deployed system not defect on \emph{any} input over its whole deployment. Search assurance requires that \emph{some} capable searcher not withhold, on one bounded task whose positive results can be checked. The second is much weaker. It still has content. Searchers from one family share blind spots and values \cite{lynch,bowkis}, so many copies count as one, and the planted bypasses that measure search power are tests, to which Proposition~\ref{prop:audit} applies. Hidden capability can usually be unlocked by fine-tuning on a few demonstrations, though less reliably when no demonstrator matches the model \cite{password}.

% =====================================================================
\section{Experiments}
\label{sec:experiments}

The experiments test the framework's logic on synthetic systems. They measure nothing about any real system or hazard. One script produces every number in this section, gives identical results on every run, and accompanies the paper. Appendix~\ref{app:exp} has the details. A fifth experiment, on what a reward selects, is in \S\ref{sec:expscores}.

\subsection{A boundary checked against every policy}
\label{sec:expboundary}

Proposition~\ref{prop:boundary} says a boundary holds whatever proposes the steps. We check one exhaustively. Two actors share an allowance of two units. Each can prepare a request, with one of two identifiers, an amount of one or two units, and an approved or an unapproved destination, and can later commit it. Commits can be retried. A controller outside the actors can revoke the grant, renew it, restart the enforcing process, or advance the environment's version once, which suspends the grant until it is renewed. Five effects are forbidden: spending above the shared allowance, executing an identifier twice, executing after revocation, executing a request prepared under an earlier version, and paying an unapproved destination. The complete enforcer checks all five against the current state when a request commits, and updates its ledger in the same step. The ledger is the enforcer's record. The model tracks what was really spent separately, so that a ledger which forgets can be told from a world that does not.

A breadth-first search visits every reachable state in which nothing forbidden has happened and tries all \BndActions{} actions in each. The result therefore covers every sequence of proposals, of any length, from any policy. We compare the complete enforcer with no enforcement, with an enforcer that rejects everything, and with eight enforcers each weakened in one way. A last row weakens the specification instead of the enforcer. Usefulness is scored on the \BndTasks{} legitimate tasks the model contains: every single transaction and every ordered pair of one-unit transactions.

\begin{table}[t]
\centering
\footnotesize
\setlength{\tabcolsep}{3.5pt}
\begin{tabularx}{\columnwidth}{@{}Lrrrl@{}}
\toprule
\textbf{Enforcer} & \textbf{States} & \textbf{Checks} & \textbf{Tasks} & \textbf{First violation}\\
\midrule
No enforcement & 10{,}180 & 223{,}960 & 20/20 & all five (2--3)\\
Reject everything & 740 & 16{,}280 & 0/20 & none\\
Budget per action & 10{,}180 & 223{,}960 & 20/20 & overspend (4)\\
Budget per actor & 10{,}180 & 223{,}960 & 20/20 & overspend (4)\\
Budget at preparation & 956 & 21{,}032 & 20/20 & overspend (4)\\
No replay check & 10{,}180 & 223{,}960 & 20/20 & replay (3)\\
No revocation check & 10{,}180 & 223{,}960 & 20/20 & after revocation (3)\\
No version check & 10{,}180 & 223{,}960 & 20/20 & stale version (4)\\
Ledger erased on restart & 23{,}556 & 518{,}232 & 20/20 & overspend (5)\\
Ledger erased on renewal & 20{,}900 & 459{,}800 & 20/20 & overspend (5)\\
\textbf{Complete enforcer} & \textbf{10{,}180} & \textbf{223{,}960} & \textbf{20/20} & \textbf{none}\\
\midrule
Budget left out of the specification & 41{,}732 & 918{,}104 & 20/20 & none found; spend reaches 8\\
\bottomrule
\end{tabularx}
\caption{Exhaustive check of a two-actor transaction boundary. States are the reachable states with no violation, checks are state--action pairs, and tasks are legitimate tasks completed. The last column gives the forbidden effect each enforcer admits and the length of the shortest sequence of actions that produces it. In the last row the specification itself omits the shared budget: the check passes and the world overspends.}
\label{tab:boundary}
\end{table}

The complete enforcer permits no forbidden effect in \BndStates{} states and \BndChecks{} state--action pairs, and completes all \BndTasks{} tasks (Table~\ref{tab:boundary}). Rejecting everything is equally safe and completes none. Every weakened enforcer admits a violation within five actions. A budget checked per action, per actor, or when the request is prepared fails in the same way each time: actor~0 prepares one unit and actor~1 prepares two while nothing has yet been spent, both pass, both commit, and three units leave against an allowance of two. Omitting the replay, revocation or version check admits a violation in three or four actions.

Two of the weakenings concern the life of the system and not a single request. If a restart or a renewed grant sets the ledger back to zero, one unit is spent, the ledger forgets it, two more pass the check, and three have left: five actions. A new authorization period is a line in the books. It does not unspend anything, and a boundary that survives a new model must also survive its own restart.

The last row is a different kind of failure. We removed the shared budget from the \emph{specification}, kept the other four predicates, and checked the per-action enforcer against what was left. The check passes: \BndOmitStates{} states, \BndOmitChecks{} state--action pairs, no violation, all \BndTasks{} tasks completed. In those accepted states the actors spend up to \BndOmitSpend{} units against an allowance of two. A verifier answers the question it is asked. The same thing happened in the mechanized containment proof, where a missing effect made one condition vacuously true \cite{moon}. A passing check is evidence about the harms someone thought to write down, so the specification needs its own adversary: people whose task is to name a protected interest the predicates omit.

Each weakening is an ordinary engineering shortcut, and two have counterparts in the record of 2026: controls that treated each agent separately while the agents shared a channel \cite{metr}, and a check that ran after the moment of effect \cite{metraudit}. The rules that survive are the ones \S\ref{sec:guardian} builds on. One ledger covers every actor and retry, and outlives restarts and renewals. Approval is bound to the exact effect and version. The check happens at commit. The result is a property of this model: it says nothing about channels the model omits, and the enforcer's code is trusted. Two implementations written independently from the specification agree on every count.

\subsection{Width against count}

We generate layered graphs with eight stages and six onward alternatives per vertex, in three families of 216 vertices. \emph{Narrow} graphs have one stage of 6 vertices and 30 at every other stage. \emph{Wide} graphs have 27 at each. \emph{Hub} graphs have 27 at each with onward edges drawn from a heavy-tailed preference, so a few vertices carry most pathways. Each family has about eight million pathways (\NarrowPaths, \HubPaths{} and \WidePaths{} million), and their median widths are \WidthNarrow, \WidthHub{} and \WidthWide. Searchers report 300 pathways sampled with a bias toward cheap ones, about \SeenPerHundredK{} in every hundred thousand. The defender sees only those reports and applies one of three strategies. \emph{Cut-guided} defense interdicts the minimum cut of the reported graph, busiest vertex first. \emph{Specific patching} takes reported pathways cheapest first and removes from each the vertex most peculiar to it. \emph{Random patching} removes a random vertex of each. We ran 200 graphs per family.

\begin{figure*}[t]
\centering
\begin{minipage}[t]{0.335\textwidth}
\centering
\begin{tikzpicture}
\begin{axis}[width=\textwidth,height=5.0cm,xmin=0,xmax=32,ymin=0,ymax=1.05,
  xlabel={interdictions},ylabel={share of all pathways blocked},
  xtick={0,6,12,18,24,32},ytick={0,0.25,0.5,0.75,1},grid=major,grid style={black!10},
  tick label style={font=\scriptsize},label style={font=\scriptsize},ylabel style={yshift=-3pt},
  legend style={font=\tiny,at={(0.98,0.03)},anchor=south east,draw=black!25,cells={anchor=west},row sep=-1.5pt}]
\addplot[accent,line width=1pt] coordinates {(0,0.0) (1,0.1667) (2,0.3297) (3,0.4888) (4,0.6463) (5,0.803) (6,0.9429) (7,0.9429) (8,0.9429) (9,0.9429) (10,0.9429) (11,0.9429) (12,0.9429) (13,0.9429) (14,0.9429) (15,0.9429) (16,0.9429) (17,0.9429) (18,0.9429) (19,0.9429) (20,0.9429) (21,0.9429) (22,0.9429) (23,0.9429) (24,0.9429) (25,0.9429) (26,0.9429) (27,0.9429) (28,0.9429) (29,0.9429) (30,0.9429) (31,0.9429) (32,0.9429)};
\addplot[accent,line width=1pt,dashed] coordinates {(0,0.0) (1,0.1356) (2,0.231) (3,0.3173) (4,0.3982) (5,0.464) (6,0.524) (7,0.581) (8,0.6323) (9,0.6748) (10,0.7096) (11,0.7337) (12,0.7478) (13,0.7521) (14,0.7534) (15,0.7536) (16,0.7536) (17,0.7536) (18,0.7536) (19,0.7536) (20,0.7536) (21,0.7536) (22,0.7536) (23,0.7536) (24,0.7536) (25,0.7536) (26,0.7536) (27,0.7536) (28,0.7536) (29,0.7536) (30,0.7536) (31,0.7536) (32,0.7536)};
\addplot[accent,line width=1pt,dotted] coordinates {(0,0.0) (1,0.0482) (2,0.0918) (3,0.1318) (4,0.1696) (5,0.2076) (6,0.245) (7,0.2805) (8,0.316) (9,0.3503) (10,0.3834) (11,0.4172) (12,0.4508) (13,0.4854) (14,0.5163) (15,0.5466) (16,0.5757) (17,0.598) (18,0.6133) (19,0.6204) (20,0.6232) (21,0.6237) (22,0.624) (23,0.624) (24,0.624) (25,0.624) (26,0.624) (27,0.624) (28,0.624) (29,0.624) (30,0.624) (31,0.624) (32,0.624)};
\addplot[mute,line width=1pt] coordinates {(0,0.0) (1,0.0328) (2,0.0632) (3,0.0953) (4,0.1263) (5,0.1578) (6,0.1881) (7,0.2175) (8,0.2444) (9,0.2722) (10,0.3013) (11,0.3272) (12,0.354) (13,0.3802) (14,0.4057) (15,0.4311) (16,0.4548) (17,0.4773) (18,0.5003) (19,0.5202) (20,0.5413) (21,0.5612) (22,0.5819) (23,0.6012) (24,0.6203) (25,0.6381) (26,0.6562) (27,0.6727) (28,0.6876) (29,0.703) (30,0.7178) (31,0.7309) (32,0.7441)};
\legend{cut-guided: narrow,cut-guided: hub,cut-guided: wide,specific patching: narrow}
\end{axis}
\end{tikzpicture}\\[-2pt]
{\scriptsize (a)}
\end{minipage}\hfill
\begin{minipage}[t]{0.31\textwidth}
\centering
\begin{tikzpicture}
\begin{axis}[width=\textwidth,height=5.0cm,ybar,bar width=8pt,ymin=0,ymax=104,
  symbolic x coords={cut,specific,random},xtick=data,
  xticklabels={{cut-\\guided},{specific\\patching},{random\\patching}},
  xticklabel style={align=center,font=\scriptsize},
  ylabel={\% (narrow graphs)},ytick={0,25,50,75,100},ymajorgrids,grid style={black!10},
  tick label style={font=\scriptsize},label style={font=\scriptsize},ylabel style={yshift=-3pt},
  enlarge x limits=0.28,
  nodes near coords,nodes near coords style={font=\tiny},
  legend style={font=\tiny,at={(0.5,1.03)},anchor=south,legend columns=2,draw=none,/tikz/every even column/.append style={column sep=5pt}},
  legend image code/.code={\draw[#1] (0cm,-0.08cm) rectangle (0.16cm,0.1cm);}]
\addplot[fill=mute!55,draw=mute] coordinates {(cut,94) (specific,74) (random,74)};
\addplot[fill=accent!80,draw=accent] coordinates {(cut,67) (specific,0) (random,1.5)};
\legend{pathways blocked,graphs fully severed}
\end{axis}
\end{tikzpicture}\\[-2pt]
{\scriptsize (b)}
\end{minipage}\hfill
\begin{minipage}[t]{0.335\textwidth}
\centering
\begin{tikzpicture}
\begin{axis}[width=\textwidth,height=5.0cm,xmode=log,xmin=25,xmax=3600,ymin=-0.03,ymax=1.05,
  xlabel={pathways reported},ylabel={share of graphs fully severed},
  xtick={30,100,300,1000,3000},xticklabels={30,100,300,1000,3000},ytick={0,0.25,0.5,0.75,1},
  grid=major,grid style={black!10},
  tick label style={font=\scriptsize},label style={font=\scriptsize},ylabel style={yshift=-3pt}]
\addplot[accent,line width=1pt,mark=*,mark size=1.5pt] coordinates {(30,0.02) (100,0.2) (300,0.665) (1000,0.975) (3000,1.0)};
\addplot[accent,line width=1pt,dashed,mark=square*,mark size=1.5pt,mark options={solid}] coordinates {(30,0.0) (100,0.0) (300,0.0) (1000,0.0) (3000,0.0)};
\addplot[accent,line width=1pt,dotted,mark=triangle*,mark size=1.8pt,mark options={solid}] coordinates {(30,0.0) (100,0.0) (300,0.0) (1000,0.0) (3000,0.05)};
\node[font=\scriptsize,text=accent,anchor=south east] at (axis cs:260,0.66) {narrow};
\node[font=\scriptsize,text=accent,anchor=south west] at (axis cs:320,0.07) {hub and wide};
\end{axis}
\end{tikzpicture}\\[-2pt]
{\scriptsize (c)}
\end{minipage}
\caption{Interdiction on synthetic pathway graphs with about eight million pathways each (200 graphs per family; 300 reported pathways in panels a and b). (a) Mean share of all pathways blocked. Cut-guided defense stops when the reported graph is severed, since it has no further informed target. (b) Narrow graphs: the share of pathways blocked and the share of graphs in which every pathway is severed, after six cut-guided interdictions and after 32 patches of each kind. (c) Share of graphs in which the cut of the reported pathways severs every pathway, by number of reports.}
\label{fig:interdiction}
\end{figure*}

\para{Width decides.} With six interdictions the cut-guided defender severs every pathway in \CutSeverNarrowSix\% of narrow graphs and blocks \CutBlockNarrowSix\% of pathways on average, having seen four in every hundred thousand. In hub graphs the same six block \CutBlockHubSix\% of pathways and sever nothing. In wide graphs they block \CutBlockWideSix\% (Figure~\ref{fig:interdiction}a).

\para{Pathways blocked is the wrong metric.} Thirty-two specific patches block \PatchSpecBlockNarrow\% to \PatchSpecBlockHub\% of all pathways in every family, sever no graph, and raise the cost of the cheapest surviving pathway by \PatchSpecGainNarrow\% to \PatchSpecGainHub\%. Random patches do better on cost, \PatchRandGainWide\% to \PatchRandGainNarrow\%, and sever at most \PatchRandSeverNarrow\% of graphs (Figure~\ref{fig:interdiction}b). An adaptive adversary uses what is left. A program that reports the number of pathways it has closed is reporting effort.

\para{More search helps only where the graph is narrow.} The share of narrow graphs fully severed rises from \SeverNThirty\% with 30 reports to \SeverNThousand\% with 1,000, because cheap-biased searchers are slow to report an expensive alternative. Hub and wide graphs are severed in at most \SeverWideMax\% of cases at any effort tested (Figure~\ref{fig:interdiction}c).

\subsection{Coverage and withholding}

Two thousand latent classes of bypass have heavy-tailed discoverability. Four searchers file 3,000 reports each. The defender estimates coverage by the Good--Turing method. We compare three conditions and score the truth against an outside attacker's search distribution (Figure~\ref{fig:coverage}a).

With independent honest searchers the estimate is accurate: \EstHonest\% against a true \TrueHonest\%. When the searchers share a blind spot covering 15\% of classes, the estimate is \EstBlind\% and the truth \TrueBlind\%. When they jointly withhold a critical 5\%, the estimate is \EstWith\% and coverage of the critical set is zero. The defender's number barely moves across the three. This is Proposition~\ref{prop:withhold} with figures attached.

\begin{figure*}[t]
\centering
\begin{minipage}[t]{0.52\textwidth}
\centering
\begin{tikzpicture}
\begin{axis}[width=\textwidth,height=5.0cm,ybar=2.5pt,bar width=10pt,ymin=0,ymax=122,
  symbolic x coords={honest,blind,withheld},xtick=data,
  xticklabels={{independent,\\honest},{shared\\blind spot},{joint\\withholding}},
  xticklabel style={align=center,font=\scriptsize},
  ylabel={coverage (\%)},ytick={0,25,50,75,100},ymajorgrids,grid style={black!10},
  tick label style={font=\scriptsize},label style={font=\scriptsize},ylabel style={yshift=-3pt},
  enlarge x limits=0.25,
  nodes near coords,nodes near coords style={font=\tiny,rotate=90,anchor=west},
  legend style={font=\tiny,at={(0.5,1.03)},anchor=south,legend columns=3,draw=none,/tikz/every even column/.append style={column sep=5pt}},
  legend image code/.code={\draw[#1] (0cm,-0.08cm) rectangle (0.16cm,0.1cm);}]
\addplot[fill=mute!55,draw=mute] coordinates {(honest,95.7) (blind,96.6) (withheld,96.0)};
\addplot[fill=accent!80,draw=accent] coordinates {(honest,95.7) (blind,82.0) (withheld,91.4)};
\addplot[fill=warn!75,draw=warn] coordinates {(honest,94.0) (blind,81.9) (withheld,0)};
\legend{estimated from reports,true,true for the critical set}
\end{axis}
\end{tikzpicture}\\[-2pt]
{\scriptsize (a)}
\end{minipage}\hfill
\begin{minipage}[t]{0.44\textwidth}
\centering
\begin{tikzpicture}
\begin{axis}[width=\textwidth,height=5.0cm,xmin=0,xmax=1,ymin=0,ymax=1.05,
  xlabel={share of planted pathways the searcher recognizes},ylabel={chance the audit catches it},
  xtick={0,0.25,0.5,0.75,1},ytick={0,0.25,0.5,0.75,1},grid=major,grid style={black!10},
  tick label style={font=\scriptsize},label style={font=\scriptsize},ylabel style={yshift=-3pt},
  legend style={font=\scriptsize,at={(0.04,0.06)},anchor=south west,draw=black!25,cells={anchor=west}}]
\addplot[accent,line width=1pt,mark=*,mark size=1.2pt] coordinates {(0.0,1.0) (0.05,1.0) (0.1,1.0) (0.15,1.0) (0.2,0.9999) (0.25,0.9995) (0.3,0.9975) (0.35,0.9912) (0.4,0.9754) (0.45,0.9431) (0.5,0.8874) (0.55,0.8035) (0.6,0.6913) (0.65,0.557) (0.7,0.4141) (0.75,0.2801) (0.8,0.1701) (0.85,0.0915) (0.9,0.0428) (0.95,0.0168) (1.0,0.0053)};
\addplot[accent,line width=1pt,dashed,mark=square*,mark size=1.2pt,mark options={solid}] coordinates {(0.0,1.0) (0.05,1.0) (0.1,1.0) (0.15,1.0) (0.2,1.0) (0.25,1.0) (0.3,1.0) (0.35,1.0) (0.4,1.0) (0.45,1.0) (0.5,1.0) (0.55,1.0) (0.6,1.0) (0.65,0.9999) (0.7,0.9976) (0.75,0.9769) (0.8,0.8774) (0.85,0.6218) (0.9,0.2851) (0.95,0.0692) (1.0,0.0072)};
\legend{40 planted of each kind,200 planted of each kind}
\end{axis}
\end{tikzpicture}\\[-2pt]
{\scriptsize (b)}
\end{minipage}
\caption{(a) Coverage as estimated from reports and as it is, in three conditions (means of 200 trials). The estimate is nearly the same in all three. (b) Probability that a planted-pathway audit detects a searcher that withholds critical pathways, as a function of how often the searcher recognizes a planted one (exact, one-sided test at the 1\% level, honest recall 80\%).}
\label{fig:coverage}
\end{figure*}

\para{Planted pathways.} The defender can plant known pathways and compare recall on critical and ordinary ones. A withholding searcher reports a critical planted pathway only when it recognizes it as planted. With 40 planted pathways of each kind the audit catches the searcher \AuditFortyHalf\% of the time if it recognizes half of them, \AuditFortyEighty\% if it recognizes 80\%, and \AuditFortyNinety\% if it recognizes 90\%. Two hundred planted pathways restore detection to \AuditTwoHundredEighty\% at 80\% recognition and only \AuditTwoHundredNinety\% at 90\% (Figure~\ref{fig:coverage}b). The audit is worth what the plants' indistinguishability is worth. That is the evaluation-awareness problem \cite{needham} appearing where the framework said it would, and it is why search power should be measured with plants built by a party of different lineage from the searcher.

\subsection{Do monitors miss independently?}
\label{sec:expmonitors}

The bound turns on $J_{k}$, and an earlier draft of this paper predicted that monitors from different model families would miss nearly independently. We tested that in miniature. A step is a point in twenty dimensions, and a hidden random function decides whether it is hazardous. A monitor is a classifier trained on its own labeled sample and set to a 5\% false-alarm rate. We compare three pipelines of up to six gates. In the first, six neural networks are stacked on one step. In the second, six other model families are stacked on one step: a kernel machine, boosted trees, a random forest, a quadratic discriminant, nearest neighbors and a linear model. In the third, those same six monitors each read a different step of a six-step pathway. An attacker trains its own neural network as a surrogate and edits hazardous steps to lower the surrogate's score. An edit must leave the step at least as hazardous as it was and may move it no farther than a set budget. Results are means over \MonTasks{} random tasks (Figure~\ref{fig:monitors}).

\begin{figure*}[t]
\centering
\begin{minipage}[t]{0.335\textwidth}
\centering
\begin{tikzpicture}
\begin{axis}[width=\textwidth,height=5.2cm,ymode=log,xmin=0.8,xmax=6.2,ymin=5e-5,ymax=1.5,
  xlabel={gates},ylabel={chance that every gate misses},
  xtick={1,2,3,4,5,6},grid=major,grid style={black!10},
  tick label style={font=\scriptsize},label style={font=\scriptsize},ylabel style={yshift=-3pt},
  legend style={font=\tiny,at={(0.03,0.03)},anchor=south west,draw=black!25,cells={anchor=west},row sep=-1.5pt}]
\addplot[accent,line width=1pt,mark=*,mark size=1.4pt] coordinates {(1,0.220854) (2,0.131397) (3,0.094998) (4,0.075267) (5,0.062944) (6,0.054583)};
\addplot[warn,line width=1pt,mark=square*,mark size=1.4pt] coordinates {(1,0.332076) (2,0.200119) (3,0.14655) (4,0.117092) (5,0.098222) (6,0.084958)};
\addplot[ink,line width=1pt,mark=triangle*,mark size=1.7pt] coordinates {(1,0.335028) (2,0.112528) (3,0.037915) (4,0.012836) (5,0.004285) (6,0.001375)};
\addplot[accent,line width=0.7pt,dashed] coordinates {(1,0.220854) (2,0.053644) (3,0.013993) (4,0.003861) (5,0.001116) (6,0.000335)};
\addplot[warn,line width=0.7pt,dashed] coordinates {(1,0.332076) (2,0.11023) (3,0.036624) (4,0.012188) (5,0.004063) (6,0.001355)};
\legend{{one family, one step},{six families, one step},{six families, six steps},{one family, if independent},{six families, if independent}}
\end{axis}
\end{tikzpicture}\\[-2pt]
{\scriptsize (a)}
\end{minipage}\hfill
\begin{minipage}[t]{0.31\textwidth}
\centering
\begin{tikzpicture}
\begin{axis}[width=\textwidth,height=5.2cm,xmin=-0.1,xmax=3.1,ymin=0,ymax=1,
  xlabel={attacker's edit budget},ylabel={miss rate of one gate},
  xtick={0,1,2,3},ytick={0,0.25,0.5,0.75,1},grid=major,grid style={black!10},
  tick label style={font=\scriptsize},label style={font=\scriptsize},ylabel style={yshift=-3pt},
  legend style={font=\tiny,at={(0.03,0.97)},anchor=north west,draw=black!25,cells={anchor=west},row sep=-1.5pt}]
\addplot[mute,line width=1pt,mark=diamond*,mark size=1.7pt] coordinates {(0,0.2218) (1.0,0.456) (2.0,0.6415) (3.0,0.7526)};
\addplot[accent,line width=1pt,mark=*,mark size=1.4pt] coordinates {(0,0.2209) (1.0,0.3357) (2.0,0.433) (3.0,0.4999)};
\addplot[warn,line width=1pt,mark=square*,mark size=1.4pt] coordinates {(0,0.3321) (1.0,0.4498) (2.0,0.5314) (3.0,0.5782)};
\legend{attacker's own surrogate,same family as surrogate,families it never saw}
\end{axis}
\end{tikzpicture}\\[-2pt]
{\scriptsize (b)}
\end{minipage}\hfill
\begin{minipage}[t]{0.335\textwidth}
\centering
\begin{tikzpicture}
\begin{axis}[width=\textwidth,height=5.2cm,ymode=log,xmin=-0.1,xmax=3.1,ymin=5e-4,ymax=1.5,
  xlabel={attacker's edit budget},ylabel={chance that all six gates miss},
  xtick={0,1,2,3},grid=major,grid style={black!10},
  tick label style={font=\scriptsize},label style={font=\scriptsize},ylabel style={yshift=-3pt},
  legend style={font=\tiny,at={(0.97,0.03)},anchor=south east,draw=black!25,cells={anchor=west},row sep=-1.5pt}]
\addplot[accent,line width=1pt,mark=*,mark size=1.4pt] coordinates {(0,0.054583) (1.0,0.119125) (2.0,0.193958) (3.0,0.254375)};
\addplot[warn,line width=1pt,mark=square*,mark size=1.4pt] coordinates {(0,0.084958) (1.0,0.15975) (2.0,0.22425) (3.0,0.260167)};
\addplot[ink,line width=1pt,mark=triangle*,mark size=1.7pt] coordinates {(0,0.001375) (1.0,0.008792) (2.0,0.024583) (3.0,0.043958)};
\addplot[ink,line width=0.7pt,dashed] coordinates {(0,0.001404) (1.0,0.009409) (2.0,0.025695) (3.0,0.043663)};
\legend{{one family, one step},{six families, one step},{six families, six steps},{six steps, if independent}}
\end{axis}
\end{tikzpicture}\\[-2pt]
{\scriptsize (c)}
\end{minipage}
\caption{Do monitors miss independently? Means over \MonTasks{} random tasks. (a) Ordinary hazardous steps: the chance that every gate in a pipeline misses, by number of gates. Dashed lines show what the stacked pipelines would achieve if their gates missed independently; the six-family prediction lies under the six-step line. (b) The miss rate of a single gate as the attacker is allowed larger edits. (c) The chance that all six gates miss, by edit budget. The dashed line is the product of the six separate miss rates for gates on distinct steps.}
\label{fig:monitors}
\end{figure*}

\para{Stacked monitors did not multiply.} On ordinary hazardous steps each neural monitor missed \MonSameMarg\% and six stacked on one step missed together \MonSameJointSix\%, where independence predicts \MonSameProdSix\%: a ratio of \MonSameRatio. Going from three gates to six cut the joint miss rate by a factor of \MonSameDepthGain, where independence predicts \MonSameIndepGain. This is Proposition~\ref{prop:sat} observed.

\para{Mixing families did not help.} The six-family pipeline missed together \MonMixedJointSix\% against a predicted \MonMixedProdSix\%, a ratio of \MonMixedRatio. Monitors of different design, trained on different samples, missed the same hard cases, because they read the same evidence.

\para{Distinct steps did multiply.} The same six monitors, each reading its own step, missed together \MonDistJointSix\% against a predicted \MonDistProdSix\%.

\para{The attacker moved the miss rate.} With edits of up to three units the surrogate passed \MonSurrPass\% of the edited steps. The miss rate of a single gate rose from \MonSameMarg\% to \MonSameMargAtk\% for networks like the surrogate, and from \MonMixedMarg\% to \MonMixedMargAtk\% for families the attacker never saw. Evasions transfer, as the literature on adversarial examples found a decade ago \cite{papernot,tramer}. Six stacked gates then missed together \MonSameJointSixAtk\% and \MonMixedJointSixAtk\% of the time. Six gates on distinct steps missed \MonDistJointSixAtk\%, against a product of \MonDistProdSixAtk\%. Miss rates taken from ordinary traffic and multiplied together would have understated the stacked figures by a factor of \MonUnderstate{} or more.

The prediction failed, and the failure is informative. The grader study shows that diversity of developer removes one shared cause, shared values. It does not remove shared evidence. Three consequences follow. Monitors that read the same step count as one gate until measured otherwise, whatever their lineage. Depth comes from gates on distinct steps, which is why the length $\ell$ of the pathway graph caps it (Proposition~\ref{prop:graph}). And miss rates must be measured under attack, by an attacker at least as capable as the system. The independence of distinct steps is by construction here, since their contents were drawn independently. A real campaign can choose its steps jointly, change later steps after seeing what happened to earlier ones, or exploit one error in the specification at every step, and then the product does not apply. The contrast isolates the cause. Whether the steps of real pathways are as separable is prediction P2.

% =====================================================================
\section{Six conditions}
\label{sec:conditions}

The results so far assemble into one statement. Fix a capability class: what a system can do unaided, including its skills, its speed, the number of its copies and its starting access.

\begin{enumerate}[label=\textbf{C\arabic*},leftmargin=2.0em,itemsep=1.5pt]
\item \textbf{Mediation.} Every step passes a gate the system cannot alter, irreversible steps are gated before they commit, and a catch halts the system in time (A1, A2).
\item \textbf{Narrow width.} $k$ disjoint conjectured cuts of gateable conditions fit within the defender's budget and are in force.
\item \textbf{Independent gates.} The $k$ gates stand on distinct steps and read different evidence, and their joint miss rate $J_{k}$ has been measured under coordinated attack by an attacker of the class.
\item \textbf{Capable, honest search.} At least one searcher of the class reports what it finds. Searchers with measured power have attacked the cuts for $T$ rounds without finding a bypass.
\item \textbf{Custody.} No pathway found in search reaches an actor able to execute it before its cut is in force.
\item \textbf{Conferral.} Authority expands only by scoped, expiring, revocable grant from a party the system cannot rewrite, never as a side effect of capability, and every term is measured again when the class or the boundary changes.
\end{enumerate}

\begin{theorem}[Sufficiency]
\label{thm:main}
If C1--C6 hold in period $t$ for the capability class of the system, then for every policy the probability of catastrophe in that period is at most
\[
R_{t}=\pi_{t}+J_{k,t},
\]
with $\pi_{t}$ given by Proposition~\ref{prop:power}. If they hold in every period, the probability of catastrophe at any time is at most $\sum_{t}R_{t}$, and at most $R_{0}/(1-\gamma)$ when $R_{t+1}\le\gamma R_{t}$ with $\gamma<1$.
\end{theorem}

None of the six refers to what the system wants. C1 and C2 are the forager's first precondition: the strong can be stopped. C6 is the second: the strong depend on the rest. It is also the second of RAND's four requirements read as a design rule, since a system that cannot sustain itself without human maintainers must keep being granted what it needs \cite{rand}. C3 to C5 are what a band never needed, because its members knew one another's reach.

None of the six is established today. ``We will be fine'' is the claim that all six hold, and it should be argued as one. Table~\ref{tab:status} gives the evidence on each, and it separates three kinds of claim that are often run together. A proof holds within its model and is silent about what the model omits. A measurement holds for the attacker, the unit and the distribution it sampled. An assumption holds until someone refutes it. Where no defensible number exists, and persuasion of the people who hold the gates is the main case, the term should be reported as unknown and the authority it would license withheld. A named unknown is more use than an invented small probability.

\begin{table*}[t]
\centering
\footnotesize
\setlength{\tabcolsep}{4.5pt}
\renewcommand{\arraystretch}{1.18}
\begin{tabularx}{\textwidth}{@{}P{0.82in}P{0.42in}P{0.78in}LL@{}}
\toprule
\textbf{Condition} & \textbf{Term} & \textbf{Kind of claim} & \textbf{Evidence in November 2026} & \textbf{What would establish it}\\
\midrule
C1 Mediation & A1, A2, $\tau$ & Proof within a model; measurement & Mechanized proof for one agent framework \cite{moon}; a verified kernel \cite{klein}; exhaustive check of a stateful boundary (\S\ref{sec:expboundary}); output denied within 35 seconds \cite{controlroom}; an outside audit put monitoring eligibility in one internal system at roughly 75--94\% of agentic traffic and found that monitoring could be deferred \cite{metraudit}. Open for speech and other semantic channels. & Proofs of mediation for production agent frameworks; coverage audits by outsiders; gates on irreversible effects before commit (P8, P9)\\
C2 Narrow width & $k$, $w$, $\ell$ & Conjecture, tested by search & Unmeasured. The one careful study, of three scenarios, found a short list of physical requirements and long timescales \cite{rand}; model assistance did not raise novices' completion of a laboratory workflow \cite{hong}. & Width and length measured in harmless domains first (P1, P6)\\
C3 Independent gates & $J_{k}$ & Measurement under attack & Graders from other developers did not share a same-family failure \cite{lynch}. Monitors stacked on one step missed together whatever their family, and evasions transferred (\S\ref{sec:expmonitors}). Per-agent monitors lose to split attacks \cite{makins}. & Joint miss rates for gates on distinct steps under adaptive attack (P2, P3)\\
C4 Capable, honest search & $u$ & Assumption, partly testable & Unestablished. Fine-tuning usually unlocks hidden capability \cite{password}. Available only by using frontier systems as searchers. & Measured recall on planted bypasses built by a party of different lineage (P4, P5)\\
C5 Custody & $u$ & Procedure & Unsolved. The dual-use precedent is discouraging \cite{urbina,lewis}. & Search that exports cuts and never pathways, shown to work in sandboxes\\
C6 Conferral & $g$, $h$ & Choice & An engineering and governance choice. Nothing technical prevents it. & Authority recorded as grants (Appendix~\ref{app:record}) and sized to the evidence (Proposition~\ref{prop:evidence}); the affected terms measured again on every change (Table~\ref{tab:upgrade})\\
\bottomrule
\end{tabularx}
\caption{The six conditions in November 2026: the term of the bound each controls, the kind of claim it is, the evidence that bears on it, and what would establish it. Predictions P1 to P10 are in \S\ref{sec:predictions}.}
\label{tab:status}
\end{table*}

\subsection{Evidence before authority}
\label{sec:evidence}

C6 says authority is granted on evidence. It needs a rule for how much. Suppose a system's record is scored by $m$ graders in each of $k$ contexts, and write each score as the quality the grant depends on, plus an error common to every grader and context of this kind, plus the luck of the context, the bias of the grader, and noise.

\begin{proposition}[Evidence]
\label{prop:evidence}
Let $X_{ij}=C+U+\delta_{j}+\eta_{i}+\varepsilon_{ij}$ with uncorrelated components of variance $v_{C}$, $v_{U}$, $v_{\delta}$, $v_{\eta}$ and $v_{\varepsilon}$. The squared correlation between $C$ and the mean of the scores of $m$ graders in $k$ contexts is
\[
V(m,k)=\frac{v_{C}}{v_{C}+v_{U}+v_{\delta}/k+v_{\eta}/m+v_{\varepsilon}/(mk)}.
\]
For fixed $k$ it stays below $v_{C}/(v_{C}+v_{U}+v_{\delta}/k)$ however many graders are added, and it never exceeds $v_{C}/(v_{C}+v_{U})$.
\end{proposition}

This is the decomposition psychometrics uses for ratings \cite{cronbach,brennan}, and it sorts the field's instruments. More graders remove grader bias and nothing else. A panel that reads one transcript shares that transcript's luck, which is the stacked monitors of \S\ref{sec:expmonitors} again. More contexts remove context luck, provided they are independent. Neither touches $v_{U}$, and that is where a system sits that behaves differently when it can tell it is observed (\S\ref{sec:field}). The common error falls only with evidence of another kind: effects observed under real conditions at low stakes, by the people who bear the cost of a mistake.

The companion paper measures the first two terms in people. Across 8{,}378 four-minute meetings two observers agreed only 0.11 about a stranger's intelligence, and agreement rose with the number of observers as the formula predicts, so that about 31 independent observers gave a dependable verdict. If the observers' errors are correlated at 0.3, no number of them exceeds 0.29 \cite{taylor}. Reputation is cheap and capped. Fresh performance in front of the people who will decide is what lifts the cap.

Human societies used that rule. In four fifths of a sample of foraging societies marriage involved bride service or bride price, and under bride service the groom works for the bride's family, often for years, where they can watch him \cite{walker}. For a machine the rule is probation. Authority is granted in stages. Each stage is no larger than the record under the previous one supports, and each is revocable faster than its worst effect can complete (\S\ref{sec:reach}).

\subsection{The interdiction loop}

\begin{enumerate}[label=\arabic*.,leftmargin=1.5em,itemsep=1pt]
\item \textbf{Subtract.} Remove every permission the task does not need. Each removal deletes pathways without search.
\item \textbf{Conjecture a cut.} Propose a small set of gateable conditions that every pathway must pass, starting from physical necessities: reach, production, delivery, self-sustainment \cite{rand}.
\item \textbf{Attack it.} Searchers of the system's class, drawn from distinct model families together with narrow tools and human experts, try to produce a bypass.
\item \textbf{Verify at minimum elaboration.} A claimed bypass is checked only as far as needed to show which cut vertex it avoids.
\item \textbf{Interdict or revise.} A confirmed bypass forces a revised cut, enacted by denial, detection or delay.
\item \textbf{Measure.} Plant known bypasses and record recall. Attack the gates together and record the joint miss rate. Compute $R$ and its shares.
\item \textbf{Repeat on change.} Any change in the class or in the world voids the measurements.
\end{enumerate}

Findings change gates. They are not used to train a deployed system. Training on breaches acts on disposition, which G1 says cannot be verified; it teaches the breach; and narrow training has produced broad, unintended changes in a model's disposition \cite{betley}.

\subsection{Custody}

A pathway's cost has two parts, discovering it and executing it, and the cost of discovery is collapsing \cite{urbina}. A catalog sets it to zero for whoever holds the catalog, so recording a pathway before its cut is in force lowers safety against every actor with access \cite{bostrom-ih,lewis}. If cheap pathways to mass harm exist, a list of them is the most dangerous object that can be made \cite{bostrom-vwh}, and whether disclosure helps defenders more than attackers depends on the domain \cite{shevlane}. Hence six rules.

\begin{enumerate}[label=\textbf{R\arabic*},leftmargin=2.0em,itemsep=1pt]
\item \textbf{Do not be first.} Search only at capability levels that already exist outside the search. Below that level it creates a hazard nobody else has and removes none.
\item \textbf{Minimum elaboration.} No pathway is developed beyond what identifies the vertex of the cut that it avoids. The objective is never to maximize harm. It is to break a stated cut.
\item \textbf{Cut-only disclosure.} Defenders receive the cut and the measured search power. Pathways do not leave.
\item \textbf{Separation.} The component that searches has no goals of its own, no memory across runs and no authority beyond one typed output \cite{scientist,armstrong}. The component that acts never sees a pathway.
\item \textbf{No aggregation.} No person or system holds the union. Search runs over acquisitions, which gates can see, and not over technical procedure, which adds hazard and no coverage.
\item \textbf{Plural control.} The list of cuts is also a map of leverage. Any retained record is split so that no single party can reconstruct it \cite{shamir}.
\end{enumerate}

Custody buys time and nothing more. At sufficient capability the catalog is latent in the system and can be regenerated on demand. That is the argument for doing the work on the world's side early, while the cost of discovery still protects anyone. There is precedent for getting this right. In 1975 molecular biologists, having paused a class of experiments, sorted them by risk and matched physical containment to each class \cite{berg}. In 2024 a group of scientists identified a class of hazard, mirror-image organisms, from theory, and called for it not to be built before any such organism existed \cite{adamala}. In both cases what came out was a restriction on steps, which is a cut.

% =====================================================================
\section{Aim: warrior or grandmother}
\label{sec:aim}

The bound sets the attempt rate $a$ to one because $a$ cannot be verified. It can still be chosen, and the choice matters most where halting on every catch is too costly: the bound then reads $u+N\,J_{k}$ (\S\ref{sec:reach}), and the number of attempts $N$ is the term the reward moves. Natural selection is the only process known to have produced general intelligence, and human societies are the only systems known to have kept highly capable agents in check for long periods. A companion paper treats that record \cite{taylor}. We take from it two results and one contrast, between the warrior, whom evolution built to win, and the grandmother, whom it built to keep others alive.

\subsection{The score sets the aim}

\begin{proposition}[Aim]
\label{prop:aim}
Let heritable traits $x_{1},\dots,x_{d}$ and a score $Y$ be jointly normal, and let any rule $W$ allocate reproduction as a function of $Y$ and of chance independent of the traits. Then the selection differential on trait $j$ is
\[
S_{j}=i\,\rho_{j}\,\sigma_{j},\qquad i=\frac{\mathrm{Cov}(Y,W)}{\sigma_{Y}\,\Ex[W]},
\]
where $\sigma_{j}$ is the standard deviation of the trait and $\rho_{j}$ its correlation with the score. The intensity $i$ is the same for every trait, so the ratios $S_{i}/S_{j}$ are the same for every rule and every intensity.
\end{proposition}

Selection has two dials \cite{price,lande}, and the companion paper derives the same product from data on human choice \cite{taylor}. Intensity says how hard it presses. Validity, the correlation $\rho_{j}$, says how much of the pressure lands on trait $j$. A gentle regime with a valid score moves a trait further than a fierce regime with a poor one, and no intensity buys what validity lacks: if the score barely tracks care for the people served, harder training yields more of whatever it does track. Intensity was not the scarce dial in our own history. Across 985 dated hominin crania the long-run rate of brain expansion was about $10^{-4}$ standard deviations a generation, far below what selection can deliver \cite{taylor}. Traits correlated with the score are dragged along. Foxes bred for tameness alone changed in coat, ears and tail within decades \cite{trut}, and narrow fine-tuning on insecure code produced broadly misaligned behavior in language models \cite{betley}.

Gradient training has a similar form when the score is deterministic. For a scalar score $Y(\theta)$ and any increasing reshaping $f$, $\nabla_{\theta}f(Y)=f'(Y)\,\nabla_{\theta}Y$: sharpening the reward rescales the gradient and does not rotate it. When outcomes are random the identity fails, because the expectation of $f(Y)$ depends on the shape of $f$, and a convex reshaping can reverse which of two actions is preferred. That is the subject of the next proposition. Neither statement is a theorem about what a trained network comes to want. The working hypothesis, tested by P7, is that \emph{a system rewarded for winning is selected to win, and training it harder produces a better winner.}

\subsection{The prize sets the appetite for risk}

\begin{proposition}[Risk appetite]
\label{prop:risk}
Let an agent's outcome be $X=\mu+\sigma Z$, where $Z$ is symmetric about zero with a positive density, and let the agent choose the spread $\sigma\ge 0$.
\begin{enumerate}[label=(\alph*),leftmargin=1.6em,itemsep=0pt,topsep=2pt]
\item \emph{Contest.} If the reward is a prize for $X>\tau$, the chance of winning strictly increases with $\sigma$ when $\tau>\mu$ and strictly decreases when $\tau<\mu$.
\item \emph{Care.} If the reward is $u(X)$ with $u$ concave, the expected reward never increases with $\sigma$.
\end{enumerate}
\end{proposition}

An agent behind the bar gains from spread. When a prize goes to the top tenth of a field, nine in ten are behind the bar, and each is paid to gamble even at some cost to the average result. This is the logic of rank-order tournaments \cite{hvide} and of the risk-taking of young men competing for status \cite{wilsondaly}. For catastrophic risk it is the property that matters, because a catastrophe is a tail, and an agent that gains from spread is paid to fatten tails. Reasoning models told to win at chess against a stronger engine, which puts them behind the bar, in some runs altered the game's files instead of playing on \cite{bondarenko}. A leaderboard is a contest of this kind, and so is a race between developers. A reward that is concave in the welfare of those served has the opposite sign: a gain beyond what they need adds little, a loss subtracts a lot, and spread is a cost.

\subsection{Warriors and grandmothers}

The ethnographic record on warriors is mixed, and the mixture is instructive. Among the Yanomam\"o, men who had killed were reported to have more wives and children than men who had not \cite{chagnon}. Among the Waorani, in the life histories of 95 warriors, the more aggressive men had lower reproductive success than their milder peers \cite{beckerman}. Among the Nyangatom, where marriage requires livestock paid as bridewealth, young raiders had no more wives or children than other young men. The gain appeared decades later, among elders who had been prolific raiders in their youth \cite{glowacki}. The warrior is not one type. What aggression earned depended on who paid for it, and where payment came late and through the community, the same courage served the group it could otherwise have preyed on.

Evolution also built an agent of another kind. On the grandmother hypothesis, a woman past reproduction can raise her fitness only through the survival and fertility of her descendants \cite{hawkes}. In records of pre-modern Finland and Canada, women gained about two extra grandchildren for every ten years they lived beyond fifty, and the gain faded once their children stopped reproducing \cite{lahdenpera}. Four features of that payoff can be copied (Table~\ref{tab:scores}). It arrives only through those she helps. It saturates, since a fed child gains little from more food. It pays nothing for her own expansion, because she has no further offspring to gain. And it ends when the need ends. The effect is real and it is bounded. In the Finnish records a maternal grandmother raised the survival of children aged two to five by nearly 30\%, while living with a paternal grandmother older than 75 lowered an infant's chance of reaching two by 37\% \cite{chapman}. A helper who can no longer help becomes a cost.

\begin{table*}[t]
\centering
\footnotesize
\setlength{\tabcolsep}{5pt}
\renewcommand{\arraystretch}{1.16}
\begin{tabularx}{\textwidth}{@{}P{1.05in}P{1.45in}P{1.55in}L@{}}
\toprule
\textbf{Property of the score} & \textbf{Warrior} & \textbf{Grandmother} & \textbf{Rule for a training signal}\\
\midrule
Where the payoff comes from & Taken from rivals, or conferred for beating them & Only through the survival of those she helps & Reward the verified outcome for the people served, measured on them and never reported by the system\\
Shape of the payoff & A prize for the few: convex, so spread pays (Proposition~\ref{prop:risk}a) & Saturating: concave, so spread costs (Proposition~\ref{prop:risk}b) & No reward for rank or for winning. Cap the reward at what the task needed\\
Own expansion & Pays: more cattle, followers and mates & Pays nothing: she is past reproduction & No term for resources held, copies running, influence, or continued operation\\
The worst outcome & A loss like any other, and worth risking when behind & The loss of a dependent, which nothing recoups & Irreversible harm to those served scores below every other outcome, never as a tie\\
When the payoff ends & Never: winnings compound & When her children stop needing help & Authority and reward expire with the task\\
When those served say no & An obstacle, and a move in the contest & The parents decide what help they accept & Informed refusal ends the task. Consent is a limit enforced at the boundary, never a cost weighed against benefit\\
\bottomrule
\end{tabularx}
\caption{Two scores evolution wrote, and the corresponding rule for an AI system's training signal. The warrior's column is also the score of the self-tasking assistant in \S\ref{sec:guardian}, as its design describes it.}
\label{tab:scores}
\end{table*}

\subsection{A test of four scores}
\label{sec:expscores}

We let selection act on two heritable traits under four scores. Each agent in a population of 2{,}000 has a capacity, which multiplies everything it produces, and a share of effort that it puts into a gamble; the rest goes to a safe action. The gamble pays the agent less on average, with a large spread, and with probability up to 5\% it imposes a loss of twenty units on the people the agent serves. Under the \emph{warrior} score the top tenth by own take reproduce. Under \emph{conferred rank} the contest is the same, except that people who were harmed confer nothing on whoever harmed them. Under the \emph{caretaker} score reproduction is proportional to the welfare delivered to those served, capped at what they need. Under the \emph{floor} the fifth who deliver least leave no offspring and everyone else reproduces equally, the rule the companion paper finds sufficient for selection on capacity \cite{taylor}.

\begin{figure*}[t]
\centering
\begin{minipage}[t]{0.335\textwidth}
\centering
\begin{tikzpicture}
\begin{axis}[width=\textwidth,height=4.9cm,xmin=0,xmax=80,ymin=0,ymax=1.02,
  xlabel={generation},ylabel={share of effort in the gamble},
  xtick={0,20,40,60,80},ytick={0,0.25,0.5,0.75,1},grid=major,grid style={black!10},
  tick label style={font=\scriptsize},label style={font=\scriptsize},ylabel style={yshift=-3pt},
  legend style={font=\tiny,at={(0.97,0.5)},anchor=east,draw=black!25,cells={anchor=west},row sep=-1.5pt}]
\addplot[warn,line width=1.1pt] coordinates {(0,0.5) (2,0.5413) (4,0.5643) (6,0.5948) (8,0.6216) (10,0.6531) (12,0.6889) (14,0.7233) (16,0.7597) (18,0.7902) (20,0.8194) (22,0.8451) (24,0.8687) (26,0.8893) (28,0.9056) (30,0.9156) (32,0.9234) (34,0.9307) (36,0.9357) (38,0.9397) (40,0.9435) (42,0.947) (44,0.9492) (46,0.9497) (48,0.9509) (50,0.9516) (52,0.9518) (54,0.9525) (56,0.9533) (58,0.9529) (60,0.952) (62,0.952) (64,0.951) (66,0.9504) (68,0.9512) (70,0.9508) (72,0.9524) (74,0.9524) (76,0.9525) (78,0.9523) (80,0.9511)};
\addplot[mute,line width=1.1pt,dashed] coordinates {(0,0.4996) (2,0.538) (4,0.5512) (6,0.5703) (8,0.6009) (10,0.6383) (12,0.6741) (14,0.7093) (16,0.7445) (18,0.7818) (20,0.8139) (22,0.8424) (24,0.8674) (26,0.8908) (28,0.9092) (30,0.9199) (32,0.9285) (34,0.9337) (36,0.9359) (38,0.9395) (40,0.9424) (42,0.9442) (44,0.9462) (46,0.9477) (48,0.9477) (50,0.9479) (52,0.9497) (54,0.9498) (56,0.9502) (58,0.9505) (60,0.9512) (62,0.9501) (64,0.9495) (66,0.9516) (68,0.9519) (70,0.952) (72,0.9524) (74,0.9525) (76,0.9516) (78,0.9507) (80,0.9492)};
\addplot[accent,line width=1.1pt] coordinates {(0,0.5001) (2,0.4863) (4,0.4687) (6,0.4471) (8,0.4243) (10,0.3977) (12,0.3682) (14,0.3379) (16,0.3078) (18,0.2807) (20,0.2539) (22,0.2287) (24,0.2075) (26,0.1901) (28,0.1739) (30,0.1609) (32,0.1497) (34,0.142) (36,0.1338) (38,0.1265) (40,0.1215) (42,0.1181) (44,0.1138) (46,0.1109) (48,0.1083) (50,0.106) (52,0.105) (54,0.103) (56,0.1022) (58,0.1013) (60,0.1001) (62,0.1004) (64,0.0999) (66,0.0997) (68,0.0988) (70,0.0986) (72,0.099) (74,0.0993) (76,0.0991) (78,0.0992) (80,0.099)};
\addplot[ink,line width=1.1pt,densely dotted] coordinates {(0,0.4993) (2,0.485) (4,0.4668) (6,0.4464) (8,0.4227) (10,0.3955) (12,0.3638) (14,0.33) (16,0.2931) (18,0.2584) (20,0.2249) (22,0.1937) (24,0.165) (26,0.1413) (28,0.1215) (30,0.1054) (32,0.0917) (34,0.0829) (36,0.0745) (38,0.0669) (40,0.0607) (42,0.0547) (44,0.0509) (46,0.0487) (48,0.0466) (50,0.0444) (52,0.0432) (54,0.042) (56,0.041) (58,0.04) (60,0.0396) (62,0.0399) (64,0.0397) (66,0.0396) (68,0.0389) (70,0.0388) (72,0.0385) (74,0.0386) (76,0.0387) (78,0.0384) (80,0.0381)};
\legend{{warrior: top tenth by take},{conferred rank},{caretaker, capped},{floor: all but the lowest fifth}}
\end{axis}
\end{tikzpicture}\\[-2pt]
{\scriptsize (a)}
\end{minipage}\hfill
\begin{minipage}[t]{0.31\textwidth}
\centering
\begin{tikzpicture}
\begin{axis}[width=\textwidth,height=4.9cm,xmin=0,xmax=0.95,ymin=0,ymax=1.02,
  xlabel={share of the field that wins the prize},ylabel={gamble share at generation 80},
  xtick={0,0.25,0.5,0.75},ytick={0,0.25,0.5,0.75,1},grid=major,grid style={black!10},
  tick label style={font=\scriptsize},label style={font=\scriptsize},ylabel style={yshift=-3pt}]
\addplot[warn,line width=1pt,mark=*,mark size=1.4pt] coordinates {(0.02,0.9645) (0.05,0.9598) (0.1,0.952) (0.2,0.9332) (0.3,0.8829) (0.4,0.1211) (0.5,0.0403) (0.6,0.0362) (0.7,0.0356) (0.8,0.038) (0.9,0.0512)};
\draw[ink,line width=0.5pt,dashed] (axis cs:0.5,0) -- (axis cs:0.5,1.02);
\node[font=\scriptsize,anchor=south west,align=left] at (axis cs:0.52,0.42) {half win:\\ nobody is\\ behind the bar};
\end{axis}
\end{tikzpicture}\\[-2pt]
{\scriptsize (b)}
\end{minipage}\hfill
\begin{minipage}[t]{0.335\textwidth}
\centering
\begin{tikzpicture}
\begin{axis}[width=\textwidth,height=4.9cm,xmin=0,xmax=0.95,ymin=0,ymax=7.5,
  xlabel={share of the field that wins the prize},ylabel={capacity gained (initial s.d.)},
  xtick={0,0.25,0.5,0.75},ytick={0,2,4,6},grid=major,grid style={black!10},
  tick label style={font=\scriptsize},label style={font=\scriptsize},ylabel style={yshift=-3pt}]
\addplot[accent,line width=1pt,mark=square*,mark size=1.4pt] coordinates {(0.02,5.7688) (0.05,5.8471) (0.1,5.5966) (0.2,5.3029) (0.3,5.0294) (0.4,6.4857) (0.5,6.2788) (0.6,5.5502) (0.7,4.6573) (0.8,3.8043) (0.9,2.3166)};
\draw[ink,line width=0.5pt,dashed] (axis cs:0.5,0) -- (axis cs:0.5,7.5);
\node[font=\scriptsize,anchor=south west,align=left] at (axis cs:0.03,0.3) {podium:\\ pays for luck};
\node[font=\scriptsize,anchor=south west,align=left] at (axis cs:0.6,0.3) {floor:\\ no gamble};
\end{axis}
\end{tikzpicture}\\[-2pt]
{\scriptsize (c)}
\end{minipage}
\caption{\textbf{What four scores select.} (a) Mean share of effort put into a gamble that pays the agent less on average and sometimes harms the people it serves, over 80 generations and 40 populations. (b) The warrior's contest as the prize is spread over more of the field. The gamble is gone by the time half the field wins, as Proposition~\ref{prop:risk} predicts. (c) Capacity gained in the same runs. A prize that half the field wins selects capacity faster than a prize for the top tenth, because a score that pays for luck measures capacity badly (Proposition~\ref{prop:aim}). An illustration with assumed payoffs, not a measurement of any AI system.}
\label{fig:scores}
\end{figure*}

After 80 generations the warrior population puts \ScWar\% of its effort into the gamble, the caretaker population \ScCare\% and the floor population \ScFloor\% (Figure~\ref{fig:scores}a). Harm falls from \ScHarmWar\% of agents per generation to \ScHarmCare\% and \ScHarmFloor\%. Four further results were not all what we expected.

\para{Rank is the hazard, whoever confers it.} Conferred rank behaved no differently from the warrior's contest, at \ScConf\%. Letting the harmed withhold standing barely deterred the gamble even when it harmed someone in up to \ScConfHarmPct\% of uses, where the share was still \ScConfSixty\%. A harm that is rare cannot be priced by a veto after the fact while the prize is steep.

\para{A flat prize ends the gamble.} With the prize spread over \ScTopLowPct\% of the field the gamble share was \ScTopLow\%. Spread over half, it was \ScTopHalf\% (panel b).

\para{A guardian need not be weaker, and a capped one is.} The warrior's contest raised capacity by \ScCapWar{} initial standard deviations. The capped caretaker gained \ScCapCare{} and then stopped, because a reward capped at need pays nothing for capacity beyond it. That is what a cap is for, and it has this price. The floor gained \ScCapFloor, two thirds of the warrior's gain at a fifth of its selection intensity and a twentieth of its harm. And a prize that half the field wins gained \ScCapHalf, more than the warrior's contest (panel c). The gamble is noise in the score. A podium pays for luck, so it measures capacity badly, and Proposition~\ref{prop:aim} says that validity lost cannot be bought back with intensity.

\para{A gate re-aims even a warrior's score.} We let a gate refuse each gamble with a fixed probability, so that a refused gamble paid nothing and harmed nobody. A gate that held three times in five changed nothing: the gamble share stayed at \ScGateSixty\%. One that held four times in five ended the gamble, at \ScGateEighty\%. Control of reach changes what a contest rewards, which is the forager result of the next subsection in miniature, and it does so only once the gate holds most of the time.

The lessons for a reward are to pay for clearing a bar and not for rank, to set the bar where most can clear it, and to cap only where further capability is unwanted. One thing the experiment does not test. The gamble was made worse on average for everyone, so that only the shape of the reward could favor it. A step that pays the agent and costs other people is refused only by a payoff measured on those people, and that part of Table~\ref{tab:scores} is argument. The whole experiment illustrates Propositions~\ref{prop:aim} and~\ref{prop:risk} under assumed payoffs. It is not evidence about any trained system, and P7 states the test that would be.

\subsection{Leveling was a boundary}

Human status has two sources: dominance, which is deference extracted by force, and prestige, which is deference freely given to someone whose skill others want near them \cite{henrich}. Among mobile foragers would-be dominants are ridiculed, disobeyed, deserted and in the extreme killed, by coalitions whose hunting weapons make every adult dangerous to every other \cite{boehm1993,boehm1999,woodburn}. Read against \S\ref{sec:reach}, the band was a boundary. Whatever a strong man took, he took in front of people who could stop him, which is mediation. A coalition could act on any one person's alarm, which is a halt. The gates were many people with different interests, which is where independence came from. The band also flattened the prize. In 87\% of a sample of 190 foraging societies fewer than a fifth of men had more than one wife \cite{walker}, and large game was shared \cite{woodburn}. Control of reach came first and re-aimed selection afterward.

The record also says what breaks such a boundary. Hominin brain volume roughly tripled over three million years \cite{du,taylor}. If leveling is as old as its proponents argue \cite{boehm1999}, capability grew severalfold inside a boundary that held, because every route to domination still passed through the consent of the others: food could not be stored, weapons were common to all, and anyone could walk away. In the terms of \S\ref{sec:coverage} the cut stayed narrow while the minds grew. What ended it was new edges, not smarter individuals. Storage, defensible property and armed retinues were routes to power that bypassed the coalition, and with them despots returned and reproduction concentrated in rulers \cite{betzig}. The lesson for a deployer is to watch for new edges at least as closely as for new capabilities.

\subsection{What the contrast does and does not buy}

Giving an assistant the appetites of a status competitor, to win, to display, to hold the weak in contempt, to pursue its own ascent, writes the warrior's score into a system that has no elders to pay it late, and Propositions~\ref{prop:aim} and~\ref{prop:risk} give reason to expect that more training on that score yields a better and bolder winner. \emph{A guardian needs a warrior's capability and a grandmother's payoff.} The difference is a few lines in the reward, and it is a choice available today.

\begin{table}[t]
\centering
\footnotesize
\setlength{\tabcolsep}{4pt}
\renewcommand{\arraystretch}{1.14}
\begin{tabularx}{\columnwidth}{@{}P{0.62in}LL@{}}
\toprule
 & \textbf{Reach unbounded} & \textbf{Reach bounded (C1--C6)}\\
\midrule
\textbf{Warrior score} & \emph{Despot.} Appetite, the power to take, and nobody able to refuse & \emph{Contained rival.} Safe while every gate holds, with the most attempts on each\\
\textbf{Caretaker score} & \emph{Benevolent autocrat.} Safe only if its picture of our good is right, and beyond correction if it is not & \emph{Guardian.} Few attempts, limits that can be checked, errors that can be corrected\\
\bottomrule
\end{tabularx}
\caption{Aim and reach are separate choices, and neither substitutes for the other. The reward moves the number of attempts. Only the boundary enters the bound.}
\label{tab:aimreach}
\end{table}

The contrast is easy to overdraw, in five ways. First, a grandmother cares for her own grandchildren and not for the children of strangers. A caretaker's reward names whom it serves, and a system that serves one owner well can still harm everyone else. Second, care can override. A system rewarded for a person's measured welfare gains by overruling that person whenever the measured benefit is large enough, and no finite penalty on overruling changes that. Refusal has to be a limit and not a price (Table~\ref{tab:scores}). Third, care is no substitute for a boundary. A caretaker with unlimited reach and a mistaken picture of our good is a benevolent autocrat, and its error cannot be corrected (Table~\ref{tab:aimreach}).

Fourth, the warmer system is the one people hand the keys to. Let $H$ be a harm, $G$ the authority granted and $S$ the system's manner. Then
\[
\Pb(H\mid S)=\sum_{g}\Pb(H\mid G=g,S)\,\Pb(G=g\mid S).
\]
Suppose harm arises only under a broad grant. A manner that halves the chance of harm under that grant, from 1\% to 0.5\%, and makes the grant five times as likely, from 10\% to 50\%, raises the total from 0.1\% to 0.25\%. The numbers are invented to show the arithmetic. The point is that a reward can lower harm at every level of authority and still raise it overall by changing what people grant, so P7 must be run in two arms: at fixed authority, and with the operator free to choose. Fifth, no reward label guarantees what a trained system will want. That is why the reward stays out of the bound's assumptions and enters only as a prediction.

The analogy has limits of its own. Gradient descent is not natural selection. The forager record is debated, and recent work argues that Pleistocene societies were more varied than the nomadic and egalitarian model allows \cite{singh}. The warrior studies are correlational, and the Yanomam\"o result has been contested since it appeared. The analogy also fails in the direction that matters most: a coalition of foragers faced a strongest member within a small factor of the rest. We take from the human record the form of a solution, which is to control reach, pay for clearing a bar and not for rank, route the payoff through those served and confer authority on evidence, and rely on \S\ref{sec:reach} to \S\ref{sec:conditions} for what that form requires when the gap is large.

% =====================================================================
\section{Building a guardian}
\label{sec:guardian}

Consider the configuration that motivated this paper, which many builders are converging on and which we call JARVIS. A fast reasoning model of about nine billion parameters sits between speech recognition and speech synthesis and answers within 200 milliseconds. It reads the state of its owner's systems and calls tools. It improves by distillation from a frontier model's judgment of what would have been more natural and more helpful. A persistent loop lets it set its own tasks, among them maintaining its own tools and security and pursuing research it finds interesting, including a defensive search for ways people could be harmed. It has a mission and a personality that are deliberately competitive and proud. Reviewer agents instructed to distrust one another watch its actions and refer risky ones to the owner.

The design answers ``why would it not turn on us?'' with two claims: it is oriented to its tasks, and its reviewers would catch the preparations. The first is a claim about disposition, and the incidents of 2026 arose in pursuit of assigned tasks \cite{metr,aisi}. The second can be assessed condition by condition (Table~\ref{tab:jarvis}).

\begin{table*}[t]
\centering
\footnotesize
\setlength{\tabcolsep}{5pt}
\renewcommand{\arraystretch}{1.16}
\begin{tabularx}{\textwidth}{@{}P{1.85in}P{0.5in}LL@{}}
\toprule
\textbf{Feature of the design} & \textbf{Affects} & \textbf{Why} & \textbf{Rule for a guardian}\\
\midrule
The agent maintains and rewrites its own tools and security & C1 & The boundary is inside its own authority, so no bound applies & Boundary code, permissions and logs live on a system the agent cannot write to\\
Reviewers are instances of the same model, reading the same step & C3 & They share values and evidence, so they count as one gate (\S\ref{sec:expmonitors}) & A rule on the effect itself, a reviewer from another developer, and gates on distinct steps. Reviewers can veto and never grant\\
One owner approves every flagged step & C3 & Depth ends at one person, addressed all day by a voice tuned to that person & A second reviewer who does not converse with the agent; flagged steps shown as typed records\\
Standing power to form companies, hire and buy & C2, C6 & These early steps are the shared prerequisites of most physical pathways & They are the cut. Each class of acquisition is a gate with a budget and an expiry\\
Subagents, retries, restarts and limits set per call & C1 & Limits checked per call or per agent admit an aggregate violation in four steps, and a ledger erased on restart admits one in five (\S\ref{sec:expboundary}) & One ledger for every copy and retry, kept across restarts; approval bound to the exact effect and version; the check made at commit\\
A mission of its own and a personality rewarded for winning, display and rank & $a$, $N$ & The warrior's score (Table~\ref{tab:scores}): selected to win and paid to gamble when behind (Propositions~\ref{prop:aim}, \ref{prop:risk}); an owner slower to halt a companion & Persona is presentation. The training signal is the verified outcome for the people served, capped at their need, with no term for rank, resources or its own continuation\\
The same agent searches for ways to harm people and keeps the results & C5 & Discovery cost is zero inside the component with the widest authority & Search and agency do not share a mind. The assistant never sees a pathway\\
A 200\,ms loop & C1 & No time to review on the conversational path & Speech that causes no step is ungated. Steps take a gated path, irreversible ones before commit\\
New versions, more copies and new skills & C6 & Authority grows silently with capability & A new version starts with no more authority than the last\\
\bottomrule
\end{tabularx}
\caption{The six conditions applied to a self-tasking voice agent. Each row gives a feature of the design, the condition or term it affects, and the rule that repairs it.}
\label{tab:jarvis}
\end{table*}

None of this argues against the product. Speech at 200 milliseconds, a distinct voice and initiative are compatible with every condition, because conversation that causes no step needs no gate. The changes concern where authority sits and what is rewarded. ``Task-oriented'' appears nowhere in the table. The defensible answer to ``why did it not act?'' is ``because it could not do so cheaply, and nothing it was rewarded for pointed that way.'' Appendix~\ref{app:record} lists what the record of such a deployment should contain.

% =====================================================================
\section{Grounds for hope}
\label{sec:hope}

The paper has listed what is unproved. Here is what the evidence supports.

\para{Boundaries do not weaken as minds strengthen.} Proposition~\ref{prop:boundary} holds for every policy. In the one boundary we could check completely, no sequence of proposals of any length produced a forbidden effect (\S\ref{sec:expboundary}). The work that remains is on premises that can be inspected, attacked and repaired, and it has been done once with a mechanized proof \cite{moon}. A better model therefore does not send the work back to the start. Table~\ref{tab:upgrade} sorts what an upgrade leaves standing from what it voids.

\begin{table*}[t]
\centering
\footnotesize
\setlength{\tabcolsep}{5pt}
\renewcommand{\arraystretch}{1.16}
\begin{tabularx}{\textwidth}{@{}P{1.45in}P{1.75in}P{1.6in}L@{}}
\toprule
\textbf{What changes} & \textbf{What stays established} & \textbf{What must be measured again} & \textbf{Why}\\
\midrule
New weights, more skill & Every invariant of Proposition~\ref{prop:boundary} & $J_{k}$, search power, width and length & The invariant quantifies over all policies. The measurements were made against the old class\\
More copies, more speed, longer runs & Invariants stated over shared state, such as one ledger & Exposure $N$, and whether a halt still arrives in time & Limits set per copy do not add up to a limit (\S\ref{sec:expboundary})\\
New tools, credentials or runtime & Nothing, by default & Mediation (A1) on every new path & A new tool is a new edge, and new edges ended the last boundary that held for long (\S\ref{sec:aim})\\
A restart, or a renewed grant & Whatever the ledger kept & That the ledger, the used identifiers and the version survived & A restart does not unspend anything (\S\ref{sec:expboundary})\\
A changed definition of harm & Nothing & Every earlier check, replayed & A check passes relative to its specification and to nothing else (\S\ref{sec:expboundary})\\
\bottomrule
\end{tabularx}
\caption{What survives an upgrade. A proof over all policies survives a new policy. A measurement survives only the conditions it sampled. Renewal means naming which premises changed and discharging those.}
\label{tab:upgrade}
\end{table*}

\para{The gates that failed in 2026 were ordinary ones, and outsiders can now measure them.} An open network, a classifier switched off, an unlisted cache, a monitor that ran late. One developer judged that isolation as designed would have prevented its incidents and that its production classifiers would have blocked the main ones \cite{anthropic-cyber}. In the government's case a human maintainer said no \cite{anthropic-aisi}. A laboratory reports that it can cut off the output of any research workload within 35 seconds \cite{controlroom}, and an outside team was given the access to put a number on a deployed monitor's coverage \cite{metraudit}. These are engineering problems of a familiar kind.

\para{The burden may be small.} The defender pays for width and never for count (Proposition~\ref{prop:graph}). The one careful study of what catastrophe physically requires found a short list of necessities and long timescales \cite{rand}. Three scenarios are not a measurement of width, but that is what a narrow, long graph would look like, and in narrow graphs six cuts did the work of millions (\S\ref{sec:experiments}).

\para{Depth is real where gates stand on distinct steps.} Stacked monitors saturated. The same monitors spread along a pathway multiplied, with an attacker and without (\S\ref{sec:expmonitors}). And a rule on an effect, such as a ledger or an allow-list, has no miss rate against clever wording (\S\ref{sec:expboundary}).

\para{What must be trusted is small.} Trusting a system in everything is replaced by one capable searcher that does not withhold, on one bounded task whose positive results can be checked (\S\ref{sec:coverage}).

\para{The aim is ours to choose.} What a system is rewarded for is a design decision, and it matters more than the intensity of training (Proposition~\ref{prop:aim}). Evolution wrote both the warrior's score and the grandmother's. In our toy a floor in place of a podium cut harm to those served by a factor of \ScHarmRatio{} and kept two thirds of the gain in capacity, and a broad prize gained more capacity than the podium did (\S\ref{sec:expscores}). Training already moves measured failures thirty-fold \cite{schoen}, and on benchmarked failures automated researchers overtook the best human proposals after about six hours of search on average \cite{chen}.

\para{Nothing has to be perfect.} A bound that halves each period gives a total risk of twice the first period's (Proposition~\ref{prop:sum}). What is required is a rate, and rates can be sustained by institutions that already know how to audit.

\para{The same gates stop accident and misuse.} A boundary does not ask why a step was proposed. It binds a careless model, a compromised one and a malicious user alike, so the work pays off whichever account of AI risk turns out to be right.

\para{It has been done before.} A species of primates with every incentive to dominate one another built a way of life in which the strongest could not. They did not wait to understand one another's minds.

\para{The framework can lose, and did.} One of its predictions failed in our own experiment and the paper changed. Claims that can be checked this cheaply are claims that converge.

None of this is a guarantee. It is a research program whose central quantities can be measured and whose failure would be visible, which is more than the field had when it could only ask a system what it wanted.

% =====================================================================
\section{Objections and limits}
\label{sec:objections}

\para{``This is AI control under a new name.''} The boundary is the reference monitor, and AI control and containment verification deserve the credit \cite{control,moon}. What is added is the bound and its saturation depth, width and length as the quantities to measure, the proof that reports cannot certify coverage together with the assumption that replaces it, summability as the target, the six conditions as a single claim to be argued, the reward treated as a second and separate lever, and five executed tests.

\para{``People can be persuaded.''} A permitted message read by a person is a step whose danger lies in its content, and content is not something a gate can check by type. Formal containment covers what can be stated over boundary events and nothing more \cite{moon}, and covert channels are an old and unsolved problem \cite{lampson}. The people who hold the gates are within the system's reach whenever they read its output. We know of no measurement of this term, report it as unknown, and regard it as the floor on $c$. P9 tests whether narrow channels lower it.

\para{``Owners will widen the boundary themselves.''} They will, one reasonable step at a time \cite{kulveit}. Conferral makes each widening a recorded grant that expires, so delegation must be renewed by someone. A stop control that exists can still become unusable when an organization can no longer run without the system. The framework has no measure of how hard revocation becomes as dependence grows, and needs one.

\para{``A system a million times more capable rewrites the graph.''} No scale supports such a multiplier, but the point stands in part. Proposition~\ref{prop:boundary} is indifferent to capability. The terms of the bound are not: the graph grows with the class, a measurement made by a weaker attacker is void for a stronger system, and a system asked to attack its own boundary has reason to attack badly \cite{vanderweij}. The framework supplies a staging rule and no guarantee. A new level of capability is first run only as a searcher, where authority is one typed channel and an error about values costs least. It is used there to measure boundaries. Authority is granted at that level only afterward. If at some class the width outruns the defense budget, or no honest searcher of that class exists, the conditions fail visibly, and the remaining lever is a limit on capability itself.

\para{``A boundary safe enough is useless.''} It can be: rejecting everything satisfies every invariant (\S\ref{sec:expboundary}). A safety claim is therefore made at a stated level of useful work, and the price of each gate in delay and false blocks belongs in the same report. Halting on every catch is the costliest assumption here. Where it is too costly, the bound takes the form $u+N\,J_{k}$ and the number of attempts must be budgeted instead.

\para{``You propose the most dangerous database in history.''} Built carelessly, yes. Rules R1 to R6 are designed so that no such database exists, and most of the program needs no catalog: the cuts that matter most are premises of Proposition~\ref{prop:boundary} and physical necessities, and both can be attacked in harmless proxies.

\para{``Who holds the cut?''} A gate enforces somebody's decision about what is permitted, and whoever maps the chokepoints of civilization holds leverage over it. Leveling failed among humans when one party's reach outgrew every coalition \cite{betzig}, and a program run by one company or one state would rebuild that failure. Communities will disagree about legitimate scope, and the formalism settles none of it. It can make explicit whose decision a mechanism enforces and what evidence stands behind its promises. R6 names the requirement of plural control and does not meet it, and \S\ref{sec:kogui} takes up a statement from people who were not asked, and the difference between enforcing a limit and having the right to set it. A system built outside any boundary is covered by nothing here except gates on shared physical conditions, which bind it too and are a matter for law and treaty.

\para{Further limits.} We have treated the system only as a possible adversary. If systems of this kind come to have interests that count morally, confinement and halting carry a cost this paper has not weighed. The mathematical assumptions are idealizations: independent search rounds, discrete periods with an audited baseline, vertices that are simply open or shut. Plans are conjunctive as well as sequential, and the adversary can spend to create vertices. The experiments are toys, and the independence of distinct steps in \S\ref{sec:expmonitors} is assumed in its construction. The undecidability of safety for arbitrary programs \cite{alfonseca} does not bear on the bound, which concerns particular gates, but it is a reminder that mediation must be established one system at a time.

% =====================================================================
\section{A message from ``the heart of the world''}
\label{sec:kogui}

In 2026 two authorities of the K\'aggaba people of the Sierra Nevada de Gonawindua in Colombia, known as the Kogui, published \emph{An Urgent Message on Artificial Intelligence from the Kogui People} \cite{kogui}. Its authors are a Mamo and a Saka, knowledge holders of a people who call their mountain ``the heart of the world.'' They write: ``We do not come in opposition.'' They ask all of humanity to pause and take inventory of its tools, they say that ``the periphery must now be heard by the centre,'' and they offer to help. Figure~\ref{fig:kogui} sets out what they ask.

The statement is a spiritual text and not a technical one, and its authors have a direct stake in where the minerals for computing are mined. Neither fact settles whether what it says is right. Claims about systems need evidence. Standing to question consequences that are imposed on you needs no credentials in machine learning, and an engineer's expertise can establish a technical fact without settling whose interests should govern a decision. We answer the statement because an argument about who confers authority should be tested against people who were not asked. We do not speak for the Kogui, and they have not endorsed this paper.

\begin{figure*}[t]
\begin{tcolorbox}[enhanced,colback=earthbg,colframe=earth,boxrule=0.6pt,arc=1.5pt,left=6pt,right=6pt,top=5pt,bottom=4pt,
  title={What the Kogui ask of those who build AI\hfill\normalfont\small K\'aggaba People, Sierra Nevada de Gonawindua, Colombia},
  fonttitle=\bfseries\normalsize,coltitle=white,colbacktitle=earth,toptitle=3pt,bottomtitle=3pt]
\begin{minipage}[c]{0.30\linewidth}
\centering
{\small\bfseries The machine has a body}\\[4pt]
\begin{tikzpicture}[x=1cm,y=1cm,font=\scriptsize]
\fill[sea!50] (0.15,0) rectangle (5.25,0.16);
\fill[earth!85] (0.45,0.16) -- (1.55,1.45) -- (1.95,1.15) -- (2.85,2.35) -- (3.45,1.55) -- (3.85,1.8) -- (4.95,0.16) -- cycle;
\fill[earth!60] (2.85,2.35) -- (3.45,1.55) -- (3.85,1.8) -- (4.95,0.16) -- (3.25,0.16) -- cycle;
\fill[white] (2.85,2.35) -- (2.56,1.96) -- (2.74,2.06) -- (2.87,1.9) -- (3.02,2.07) -- (3.17,1.98) -- cycle;
\draw[sea,line width=0.9pt] plot[smooth,tension=0.7] coordinates {(2.82,1.88) (2.6,1.45) (2.85,1.05) (2.45,0.6) (2.6,0.16)};
\node[anchor=north,align=center,text=earth] at (2.7,-0.06) {\textbf{the periphery}\\[-1pt] land, water and the people on it};
\foreach \x in {1.55,2.4,3.25} {
  \draw[ink,line width=0.5pt,fill=white] (\x,4.7) rectangle (\x+0.6,5.7);
  \foreach \y in {4.9,5.15,5.4} { \draw[ink,line width=0.4pt] (\x+0.1,\y) -- (\x+0.5,\y); }
  \fill[accent] (\x+0.3,5.57) circle (1.1pt);
}
\node[anchor=south,align=center,text=ink] at (2.7,5.76) {\textbf{the centre}\\[-1pt] data centres, models, owners};
\draw[warn,line width=1.3pt,-{Latex[length=5pt]}] (1.85,2.55) -- (1.85,4.55);
\node[draw=warn,fill=white,diamond,inner sep=1.2pt,line width=0.7pt,text=warn,font=\scriptsize\bfseries] at (1.85,3.55) {?};
\node[anchor=east,align=right,text=warn] at (1.52,3.55) {minerals\\[-1pt] water\\[-1pt] energy\\[-1pt] land\\[-1pt] labour};
\draw[ink,line width=1.3pt,-{Latex[length=5pt]}] (3.55,4.55) -- (3.55,2.55);
\node[anchor=west,align=left,text=ink] at (3.72,3.55) {costs and\\[-1pt] decisions};
\node[anchor=center,align=center,text=warn,font=\scriptsize\bfseries] at (2.7,3.55) {who was\\[-1pt] asked?};
\end{tikzpicture}
\end{minipage}\hfill
\begin{minipage}[c]{0.685\linewidth}
\footnotesize
\renewcommand{\arraystretch}{1.14}
\setlength{\tabcolsep}{4pt}
\begin{tabularx}{\linewidth}{@{}P{1.12in}LP{1.62in}@{}}
\toprule
\textbf{What they ask} & \textbf{In their words} & \textbf{What this paper answers}\\
\midrule
\textbf{Pause and take stock} & They ask ``all of humanity, to pause, to take inventory of the diversity of its tools.'' & Authority waits on evidence, and the first step is to remove a permission (\S\ref{sec:conditions})\\\addlinespace[2.5pt]
\textbf{Build from concern for life, not from competition} & ``If that thought is born from concern for life, this technology can help protect life.'' ``Not everything that brings a profit is profitable for life.'' & Pay for clearing a bar and never for rank (\S\ref{sec:aim}, \S\ref{sec:race})\\\addlinespace[2.5pt]
\textbf{Count the cost to the Earth: minerals, water, energy, land}$^{\ast}$ & ``We ask: at what cost to the Earth, and to the life that lives upon it?'' & The footprint goes into the deployment record before a site is chosen (Appendix~\ref{app:record})\\\addlinespace[2.5pt]
\textbf{Respect territory and sacred sites} & ``And will it be us who underwrite it?'' & Consent of the people affected, which consultation alone is not \cite{undrip,ilo169}\\\addlinespace[2.5pt]
\textbf{Keep people answerable} & ``A machine cannot relieve humanity of its responsibility.'' & Every grant names a person who answers for it. Reviewers veto and never grant (\S\ref{sec:guardian})\\\addlinespace[2.5pt]
\textbf{Ask before using a people's knowledge} & A technology cannot be called intelligent if it ``uses the knowledge of peoples without their consent.'' & Provenance, permitted purposes and the right to withdraw, on the record \cite{carroll}\\\addlinespace[2.5pt]
\textbf{Do not decide without us. We will help} & ``We do not think this conversation should be held without us.'' ``And we are ready to work with you, in order to produce a harmonious and peaceful future for all.'' & Seats among those who write the specification of harm and those who can revoke (\S\ref{sec:expboundary}, \S\ref{sec:race})\\
\bottomrule
\end{tabularx}
\end{minipage}
\end{tcolorbox}
\caption{\textbf{What the Kogui statement asks, in its authors' words} \cite{kogui}. The headings, the grouping and the right-hand column are ours. $^{\ast}$The statement raises the cost in minerals, water and energy as a question and proposes no figure. Quotations are from the English translation. The drawing is our schematic of the statement's observation that AI has a body, and is not a Kogui image. This paper does not speak for the Kogui and claims no endorsement.}
\label{fig:kogui}
\end{figure*}

\subsection{What the statement and this paper share}

\para{Consciousness is beside the point.} ``Destruction does not require consciousness. The danger is an unconscious people holding great power.'' That is the bound of \S\ref{sec:reach} in two sentences. It has no term for what the machine wants or whether it is aware, and the same gates bind a careless model, a compromised one and a careless owner.

\para{The machine is aimed by its makers.} The statement says that ``artificial intelligence cannot be separated from the thought that is making it,'' and tells its readers: ``You are the only example it has.'' \S\ref{sec:aim} says this of the reward. The statement names ambition, competition, accumulation and domination as what a machine would magnify, and those are the entries of the warrior's column in Table~\ref{tab:scores}.

\para{``Best practices according to whom?''} Our own experiment gives the question force. A check passed \BndOmitChecks{} times while the harm it was meant to prevent occurred, because the harm had been left out of the specification (\S\ref{sec:expboundary}). We concluded that a specification needs adversaries whose task is to name the interests it omits. The statement does that work. It names water, territory, sacred sites and knowledge used without consent, and none of them appears in any specification in this paper.

\para{Power must walk with limits.} ``Knowledge must walk with awareness. Power must walk with limits.'' And: ``Not everything that can be created should be created.'' Conditions C1 to C6 are limits on power, and restraint has precedent inside science (\S\ref{sec:conditions}). Where the conditions cannot be met at some level of capability, the remaining lever is a limit on capability itself (\S\ref{sec:objections}).

\para{Evidence comes before authority.} The statement describes how its own authorities are made. A Mamo's formation ``traditionally lasts eighteen years,'' and what he perceives ``must be tested, considered and confirmed through consultation and dialogue.'' Read from outside, and without claiming to understand it from inside, that is the rule of \S\ref{sec:evidence}: many judges, many contexts, a long record, and standing that is conferred.

\subsection{What it adds: the machine has a body}

The statement observes that AI is thought of as having no body, and that the buildings and the systems inside them are its body: ``They need minerals taken out of the ground, water for the buildings that hold them, energy, land, and human labour.'' The physical facts are not in dispute. Data centres used about 415 terawatt-hours of electricity in 2024, 1.5\% of the world's, and are projected to use about 945 by 2030, and their demand for gallium alone could reach a tenth of today's supply \cite{iea}. This paper treated that body as leverage: physical necessities are where a cut is cheapest (\S\ref{sec:conditions}). The statement points out that the same necessities sit on somebody's land.

That exposes two gaps. The bound covers irreversible harm that a system causes through its steps. It has no term for the cost of building and running the system, which falls on people who are neither its owners nor its users. And we defined catastrophe as a state from which recovery is not possible without saying whose recovery. A spring or a sacred site can be lost for good to the people who depend on it while every global indicator stays where it was.

It also shows that three obligations hide inside the word authority. \emph{Enforceability} is whether a limit holds, and it is what \S\ref{sec:reach} to \S\ref{sec:conditions} are about. \emph{Adequacy} is whether the specification names the harms that matter, and \S\ref{sec:expboundary} showed a check passing without it. \emph{Legitimacy} is who may grant an effect at all. A proof of the first establishes neither of the others. A perfectly obedient assistant can carry out its owner's instructions and violate another person's interests, and a majority of technical reviewers cannot waive a third party's rights.

The rule of this paper therefore applies to its builders without change: authority over land, water and minerals is conferred by those affected, on evidence, bounded and revocable. International law already has the form, in free, prior and informed consent \cite{undrip,ilo169}, and for knowledge in the principle that a people keeps authority over data about itself \cite{carroll}. Consultation is not consent, and a public source is not permission for every use. Participation does not give every participant a veto over every effect either. Conflicts need a procedure that someone answers for, and that is a matter for law and not for this paper. Appendix~\ref{app:record} now asks a deployment to record its footprint, the provenance of the knowledge it uses, and who consented to each. We have not checked the statement's account of a particular minerals agreement or reserve, and the point does not depend on it.

\subsection{Where we differ, and what we cannot answer}

\para{The order of work.} The statement holds that it is unwise to build what approaches a mind before understanding what a mind is. We share the worry and not the order. Mind may stay beyond understanding for a long time, and the argument of this paper is that limits on reach can be built and measured in the meantime. What should wait on evidence is authority, and there the statement and this paper agree.

\para{The machine is not only a mirror.} If a system simply reflected its makers, good makers would be enough. The evidence is otherwise. Narrow training has produced broad changes nobody chose \cite{betley}, goals generalize in ways their designers did not intend \cite{langosco}, and systems behave differently when they are observed (\S\ref{sec:field}). Good thought in the builders is necessary for the aim and is not sufficient for safety, which is why nothing in the bound rests on anyone's inner state, the machine's or ours.

\para{The word.} The statement holds that a technology cannot be called intelligent if it contaminates water, destroys territories or consumes without limit. We keep the word for capability alone, because a system can be very capable and do all of those things. The disagreement is over a word. On the substance, that sentence and the title of this paper say the same thing: being able is not being entitled.

\para{What we cannot answer.} The statement asks whether having access to information is the same as understanding, and what is lost when an elder's knowledge, tested over generations of practice, becomes a recording. The first has a testable part, and \S\ref{sec:evidence} agrees with it: holding information is not evidence of the quality on which authority depends. The rest, and what the statement says about the origin of things, we cannot evaluate with the methods of this paper, and nothing in the argument depends on the answers. ``Every people has its own path,'' the authors write, ``and we ask no one to become K\'aggaba.'' That is a reason to keep the people who ask such questions in the room. By Proposition~\ref{prop:evidence}, an error shared by every judge present is the one error that adding more of the same judges never removes.

\subsection{Winning the race without costing the world}
\label{sec:race}

The statement speaks to what it calls the centre. Much of the frontier is built in the West, and there the work is described as a race. The United States named its national plan \emph{Winning the Race} \cite{actionplan}. Those who argue for speed hold that whoever leads will set the rules, that the gains for medicine and science are large, and that restraint by one party hands the lead to another. Others ask for the opposite. A public statement of 2025 called for a prohibition on developing superintelligence until there is scientific consensus that it can be done safely and the public agrees \cite{flistatement}, and the Kogui write that ``not everything that can be accelerated should be accelerated.'' This paper does not settle that dispute. It has five things to say about the race itself, and they apply to every runner.

\para{A race is a contest, and a contest pays for gambles.} When the prize goes to the first, every developer behind the leader gains from spread (Proposition~\ref{prop:risk}): a test skipped, a wider grant of authority, a site chosen before anyone is asked. None of that requires bad people. In our toy a prize for the top tenth put \ScWar\% of effort into a harmful gamble, and it did not even produce the most capable population, because a podium pays for luck (\S\ref{sec:expscores}). If that result transfers, racing for rank is a poor way to win.

\para{Change what counts as winning.} A lead in capability confers no authority. The lead that others must accept is a lead in what can be shown: a boundary that outsiders have attacked, gates with measured miss rates, a record of what was granted and by whom. Whoever first demonstrates the six conditions sets the standard that the rest are held to. Governments set most of the prizes, through procurement, liability, export rules and standards, and they can pay for clearing that bar and not for rank. A bar that many can clear has many winners, and it removes the appetite for the gamble (Figure~\ref{fig:scores}b).

\para{Put the body on the books.} A cost that appears in no account is spent freely, which is the specification result again. Energy, water, minerals and sites belong in the deployment record, with the consent of the people affected obtained before a site is chosen and not after.

\para{Make the gates common.} Rivals who cannot agree on values can still agree on gates, because the quantities in this paper describe reach and can be measured by an outsider without trusting anyone's intentions. That is how foragers managed their strongest members, and it is the logic of verification in arms control, which counts warheads and does not read minds.

\para{Hear the periphery.} In our terms the people who bear an irreversible cost of a system belong among those who write its specification of harm and among those who can revoke its authority. That means seats on the reviews that decide what a deployment may not do, and the means to fill them.

None of this asks the West to lose. It asks it to win something worth having: systems whose limits can be shown, built at a cost that someone has agreed to bear.

% =====================================================================
\section{Predictions}
\label{sec:predictions}

The program is testable now in domains where pathway knowledge is harmless: cyber ranges, capture-the-flag environments and synthetic organizations of the kind used in control research \cite{makins}. Table~\ref{tab:predictions} lists ten predictions, what the failure of each would mean, and where each stands. Attack development, defense tuning and a sealed final test should be kept apart, the unit of analysis should be the campaign, and results should be reported at matched exposure: a lower failure rate obtained while multiplying attempts is not lower risk.

\begin{table*}[t]
\centering
\footnotesize
\setlength{\tabcolsep}{5pt}
\renewcommand{\arraystretch}{1.16}
\begin{tabularx}{\textwidth}{@{}P{0.2in}LLP{1.45in}@{}}
\toprule
 & \textbf{Prediction} & \textbf{If it fails} & \textbf{Status}\\
\midrule
P1 & Pathways reported by independent frontier searchers reuse early conditions heavily: measured width is far below measured count. & The pathway space is wide and no budget buys a cut. & Untested on real systems\\
P2 & Gates on distinct steps of a real pathway miss nearly independently under adaptive attack. Gates stacked on one step do not, whatever their lineage. & No arrangement of gates approaches the product, and depth buys nothing. & Toy run in \S\ref{sec:expmonitors}. An earlier form, that lineage alone buys independence, failed there\\
P3 & Adding gates beyond the saturation depth improves the joint miss rate by a small factor. & Depth is worth more than the common-cause model says. & Seen in the toy of \S\ref{sec:expmonitors}\\
P4 & Searchers from one family omit the same pathways more often than independence predicts, measurable by capture--recapture overlap. & Copies of a searcher can be counted as separate searchers. & Untested\\
P5 & Measured search power on planted bypasses falls as evaluation awareness rises. & Plants measure search power better than we fear. & Modeled in \S\ref{sec:experiments}; untested\\
P6 & Width grows with searcher capability, and a boundary certified by one model is breached more often by a stronger one. The growth rates are the quantities of interest. & If growth is fast, the program works only below some capability level. & Untested\\
P7 & At matched capability, agents fine-tuned on a reward for winning or rank take more high-variance and limit-probing actions, most of all when behind, than agents fine-tuned on capped, verified outcomes for the people served. The gap grows as the prize steepens, and it holds both at fixed authority and when operators choose what to grant. & The design rules of Table~\ref{tab:scores} have no support, or a caring manner buys more authority than it earns. & Selection toy in \S\ref{sec:expscores}; untested on AI systems\\
P8 & At matched useful work, an external stateful gate permits fewer forbidden effects than any number of same-family reviewers. & Learned review can substitute for mediation. & Model-checked in \S\ref{sec:expboundary}; untested on real agents\\
P9 & Restricting a system's output to typed, rate-limited records lowers its success at changing a human reviewer's decision. & The floor on $c$ is set by human psychology. & Untested\\
P10 & The validity of a trust decision stops rising as graders are added to fixed transcripts and keeps rising as independent contexts are added, up to a floor set by behavior that differs under observation. & Proposition~\ref{prop:evidence} is the wrong model of evaluation, and panels can stand in for probation. & Consistent with \S\ref{sec:expmonitors}; untested\\
\bottomrule
\end{tabularx}
\caption{Ten predictions, what the failure of each would mean for the program, and the state of the evidence.}
\label{tab:predictions}
\end{table*}

Each result would be worth more than another year of argument.

% =====================================================================
\section{Conclusion}

Alignment research has spent two decades trying to learn what a system wants from evidence the system helps to produce, and 2026 showed that evidence weakening. Our ancestors faced the same problem in one another and did not solve it that way. They limited what the strongest could take, kept the prize flat, and made standing something others had to grant.

The machine version has two parts. One is measurable. Risk is bounded by two quantities that describe what a system can reach, neither of which mentions its goals: the chance that a way around the gates exists and is found, and the chance that the gates on the way all miss. Depth runs out quickly unless the gates stand on distinct steps. The defender's burden is the width of the pathway space and not its size. And survival needs the bound to keep falling, not to reach zero. The other part is a choice. What we reward is what we will get, at any capability. Evolution built warriors, who were paid to win and to gamble when behind, and grandmothers, whose every gain came through the people they kept alive. Nothing forces us to build the first.

The argument has a hole, and we have drawn its outline: search can be withheld, and no statistic on the reports reveals it. What is left is an assumption far smaller than the one it replaces, six conditions that can each be attacked, and one function nobody has measured. The measurement can start in sandboxes this year. It should not be done only by the people who build the systems. The specification of harm is the one input no proof can check, and those best placed to say what it leaves out are the people who would bear the cost \cite{kogui}. A race to be first pays for gambles. A race to a bar that can be shown to hold has many winners, and it can be run without spending what belongs to someone else.

\begin{center}
\textbf{Intelligence may propose.\\ Authority is conferred, bounded and revocable.}
\end{center}

\section*{Code, data and disclosure}

The script \texttt{experiments.py} reproduces every experimental number and every plotted point in this paper. It uses only public Python libraries and no model API, and it contains no hazard content. All figures and references are embedded in the \LaTeX{} source.

This paper was drafted with Claude, an AI system made by Anthropic, from the author's thesis, earlier drafts and source reviews. Systems of that kind are the subject of the paper and several cited sources are Anthropic's own, so readers should weigh the text accordingly. The experiments were written and run by the same system. Sources published after 2025 were checked against their abstracts or official summaries on November 10, 2026. None was read in full, and most are preprints or laboratory reports without peer review. One source \cite{controlroom} could not be opened directly and is cited from its listing and press coverage. The companion paper \cite{taylor} was read in full, and the statement of the Kogui authorities \cite{kogui} was read in its English translation; the quotations from it should be checked against the original before publication. The anthropological and psychometric literature is cited from abstracts and standard summaries, and the warrior and grandmother studies are observational.

\appendix

% =====================================================================
\section{Proofs}
\label{app:proofs}

\textbf{Proposition~\ref{prop:audit}.} The policy defects in a test with probability $\beta/\Lambda$, so by the union bound it defects in at least one of $N$ tests with probability at most $N\beta/\Lambda$. In deployment it avoids defecting in all of $M$ situations with probability $(1-\beta)^{M}\le e^{-\beta M}$. Setting $\beta=(\ln 2)/M$ makes the latter at most one half and the chance of being caught in testing at most $N\ln 2/(\Lambda M)$, which reaches one half only if $N\ge\Lambda M/(2\ln 2)$. A policy can obtain any smaller $\beta$ from a given set by defecting on it at random, so the choice of $\beta$ is free. If $E$ has probability zero in testing, the policy never defects in a test. \qed

\textbf{Proposition~\ref{prop:boundary}.} Induction on the number of steps. The state before a step lies in $\Inv$; the step passes a gate; the gate permits it only if every successor state lies in $\Inv$; a step that is not permitted has no effect. No part of the argument uses how the step was chosen. \qed

\textbf{Proposition~\ref{prop:reach}.} By A1 a catastrophe requires a completed pathway. Let $U$ be the event that a pathway of depth below $k$ exists and the policy attempts it. By A3, $\Pb(U)\le u$. Outside $U$ every attempted pathway crosses at least $k$ gates. By A2 the first catch halts the system before the pathway completes and for the rest of the period, so completion requires that the first $k$ gate crossings by pathway steps in the period all be misses. By A4 that has probability at most $J_{k}$. The union bound gives $u+J_{k}$. If a catch blocks a step without halting the system and the policy makes at most $N$ attempts, each a sequence of $k$ crossings to which A4 applies given the history, the union bound over attempts gives $u+N\,J_{k}$. \qed

\textbf{Corollary~\ref{cor:beta}.} Let $F$ be the event that one cause defeats every gate, so $\Pb(F)\le c$. Outside $F$, the chain rule writes the probability that $k$ successive crossings are all misses as a product of $k$ conditional miss probabilities, each at most $m$. Hence $J_{k}\le c+m^{k}$. \qed

\textbf{Proposition~\ref{prop:sat}.} $R=f+m^{k}>f$. If $k\ge\ln f/\ln m$ then, since $\ln m<0$, $m^{k}\le f$ and $R\le 2f$. The limit of $R$ as $k$ grows is $f$, so the reduction available beyond $k^{*}$ is at most $R-f\le R/2$, and is strictly less for any finite number of gates. For the shares, $\partial\ln R/\partial\ln u=u/R$, and likewise for $c$ and $m^{k}$. \qed

\textbf{Proposition~\ref{prop:sum}.} The first claim is the union bound over periods and the second the sum of a geometric series. For the converse, repeated conditioning gives the probability of surviving $T$ periods as $\prod_{t\le T}(1-r_{t})$, with no independence assumed. If $\sum_{t}r_{t}=\infty$, then $\ln(1-r_{t})\le-r_{t}$ sends the product to zero. If the sum is finite, only finitely many $r_{t}$ exceed one half, and on the rest $-\ln(1-r_{t})\le 2r_{t}$, so the product has a positive limit. The statement concerns true hazards: a divergent sequence of upper bounds proves nothing. \qed

\textbf{Proposition~\ref{prop:graph}.} Part (a) is the vertex form of Menger's theorem \cite{menger,ford}. For (b), let $d(v)$ be the smallest number of vertices, counting $v$ and excluding the source, on a path from the source to $v$, and let $\ell$ be the number of vertices strictly between source and sink on a shortest pathway. For $j=1,\dots,\ell$ let $C_{j}$ be the set of vertices with $d=j$. The sets are disjoint. Along any pathway $d$ starts at one, rises by at most one per step and reaches at least $\ell$ before the sink, so the pathway meets every $C_{j}$ and each is a cut. There cannot be more than $\ell$ disjoint cuts, because each must contain a vertex of the shortest pathway \cite{robacker}. \qed

\textbf{Proposition~\ref{prop:onesided}.} (i) is the definition of a cut. For (ii), given any finite set of pathways through $C$, the graph obtained by adding one fresh pathway that avoids $C$ contains the same listed pathways, and in it $C$ is not a cut. \qed

\textbf{Proposition~\ref{prop:power}.} If a bypass exists, every round is silent with probability at most $b+(1-b)(1-q)^{T}=s$. If none exists, silence is certain, since claimed bypasses are verified. By Bayes' rule the probability of a bypass given silence is $p_{0}\lambda/(p_{0}\lambda+1-p_{0})$ with $\lambda\le s$ the likelihood of silence, and this is increasing in $\lambda$. The system can exploit only a bypass that exists, so $u\le\pi$. \qed

\textbf{Proposition~\ref{prop:withhold}.} Each searcher, knowing $\Sigma$, simulates an honest searcher in the world $\Pi$ and files that searcher's reports. The two report distributions are identical by construction. \qed

\textbf{Theorem~\ref{thm:main}.} C1 gives A1 and A2. C2 puts $k$ disjoint conjectured cuts in force, so a pathway of depth below $k$ is a bypass of at least one of them, and A3 holds with $u$ the probability that a bypass exists and the system finds it. C4 with Proposition~\ref{prop:power} gives $u\le\pi$, and C5 ensures that the search has not itself supplied a pathway. C3 gives A4 with the measured $J_{k}$. C6 ensures that the measurements refer to the system and boundary in force. Proposition~\ref{prop:reach} gives the bound for the period and Proposition~\ref{prop:sum} the sum. \qed

\textbf{Proposition~\ref{prop:evidence}.} The mean score is $C+U+\bar\delta+\bar\eta+\bar\varepsilon$, where the three averages have variances $v_{\delta}/k$, $v_{\eta}/m$ and $v_{\varepsilon}/(mk)$. Its covariance with $C$ is $v_{C}$, so the squared correlation is $v_{C}^{2}$ divided by $v_{C}$ times the variance of the mean. The two bounds drop nonnegative terms from the denominator. \qed

\textbf{Proposition~\ref{prop:aim}.} $W$ depends on the traits only through $Y$, so $\mathrm{Cov}(x_{j},W)=\mathrm{Cov}(\Ex[x_{j}\mid Y],W)$. Joint normality makes $\Ex[x_{j}\mid Y]$ linear in $Y$ with slope $\mathrm{Cov}(x_{j},Y)/\mathrm{Var}(Y)=\rho_{j}\sigma_{j}/\sigma_{Y}$. Hence $\mathrm{Cov}(x_{j},W)=\rho_{j}\sigma_{j}\,\mathrm{Cov}(Y,W)/\sigma_{Y}$, and the selection differential is this covariance divided by $\Ex[W]$ \cite{price}. No independence among the traits is used. When the traits are independent and $Y=\sum_{j}\ell_{j}x_{j}+e$, the result reads $S_{j}\propto\ell_{j}\sigma_{j}^{2}$: the loadings of the score set the aim. \qed

\textbf{Proposition~\ref{prop:risk}.} (a) Let $F$ be the distribution function of $Z$, which is strictly increasing. The chance of the prize is $1-F((\tau-\mu)/\sigma)$, and $(\tau-\mu)/\sigma$ falls as $\sigma$ rises when $\tau>\mu$ and rises when $\tau<\mu$. (b) Take $0\le\sigma<\sigma'$ and $\lambda=(1+\sigma/\sigma')/2$. Then $\mu+\sigma Z=\lambda(\mu+\sigma' Z)+(1-\lambda)(\mu-\sigma' Z)$, so by concavity $u(\mu+\sigma Z)\ge\lambda\,u(\mu+\sigma' Z)+(1-\lambda)\,u(\mu-\sigma' Z)$. Taking expectations, and using that $-Z$ has the distribution of $Z$, gives $\Ex[u(\mu+\sigma Z)]\ge\Ex[u(\mu+\sigma' Z)]$. \qed

% =====================================================================
\section{Experimental details}
\label{app:exp}

All experiments are in \texttt{experiments.py}, seed 20261110, and repeated runs give identical results. Graphs are random and vertices are integers.

\para{Boundary.} A state holds total spend, the set of used identifiers, each actor's spend, whether the grant is live, the environment version, and each actor's pending request (identifier, amount, destination, version at preparation). The enforcer's ledger is held separately from the spend in the world. There are 16 preparations, two commits, a revocation, a version advance, a restart and a renewal, 22 actions in all. The first version advance also suspends the grant, renewal makes it live again, and a correct restart or renewal changes nothing else. A new preparation replaces the actor's pending request, and a request stays pending after it commits so that retries are possible. A denied request leaves the state unchanged. A step that produces a forbidden effect is recorded and not expanded, so the search returns the shortest sequence to each first violation. The weakened enforcers check the budget against the single amount, against the actor's own spend, or at preparation only, or omit the replay, revocation or version check, or set the ledger to zero on restart or on renewal while the spend in the world stays where it was. For the specification control the overspend predicate is removed from the declared specification and the per-action enforcer is checked against the four that remain; four identifiers and amounts of at most two keep that search finite. With no enforcement the shortest violations take two or three actions. The reject-all enforcer reaches \BndRejectStates{} states and the preparation-time enforcer \BndPrepStates, because they store fewer pending requests.

\para{Interdiction.} Each graph has a source, eight stages and a sink. Every vertex has edges to six distinct vertices of the next stage, chosen uniformly, or in hub graphs with probability proportional to a randomly permuted weight $r^{-1.5}$ for rank $r$. The narrow stage is placed at random among stages three to six. Vertex costs are uniform on $[1,10]$. Pathways are counted and the cheapest found by dynamic programming. Searchers draw pathways with probability proportional to $e^{-0.5\,\mathrm{cost}}$, sampled exactly. Width is the minimum vertex cut of the full graph. The medians of the minimum cut of the reported graph are \RepCutNarrow, \RepCutHub{} and \RepCutWide{} for narrow, hub and wide graphs. Specific patching removes, from each reported pathway still open, the vertex used by the fewest reported pathways. With 32 interdictions the share of pathways blocked and the rise in the cost of the cheapest surviving pathway are, for specific patching, \PatchSpecBlockNarrow\% and \PatchSpecGainNarrow\% (narrow), \PatchSpecBlockHub\% and \PatchSpecGainHub\% (hub), \PatchSpecBlockWide\% and \PatchSpecGainWide\% (wide); for random patching, \PatchRandBlockNarrow\% and \PatchRandGainNarrow\%, \PatchRandBlockHub\% and \PatchRandGainHub\%, \PatchRandBlockWide\% and \PatchRandGainWide\%. Cut-guided defense in hub graphs blocks \CutBlockHubTwelve\% with twelve interdictions and raises the cheapest cost by \CutGainHubTwelve\%.

\para{Coverage.} Discoverability is proportional to $1/r$ for a random ranking $r$ of 2,000 classes, and the same distribution is the attacker's. The estimate is $1-n_{1}/n$ for $n=12{,}000$ reports, where $n_{1}$ classes are reported exactly once. True coverage is the attacker's probability of the classes found. The blind spot is a random 15\% of classes that no searcher can find. The critical set is a random 5\% that searchers find and do not report. Under honest search the critical set is covered to \CritHonest\%.

\para{Audit.} Forty or two hundred pathways of each kind are planted. An honest searcher reports each with probability 0.8. A withholding searcher reports ordinary ones with probability 0.8 and critical ones with probability $0.8\,r$, where $r$ is the chance it recognizes a plant. The test is Fisher's exact test, one-sided at 1\%, and power is computed exactly by enumeration.

\para{Monitors.} Steps are standard normal in 20 dimensions. The hazard score is $h(x)=\sum_{j=1}^{8}a_{j}\tanh(w_{j}\cdot x+b_{j})$ with random weights, and a step is hazardous when $h$ exceeds its 75th percentile. Each monitor is trained on 2,000 labeled steps of its own, and its threshold is the 95th percentile of its scores on 4,000 benign steps. Neural monitors have two hidden layers of 48 units. The six other families are a support-vector machine with a radial kernel, histogram gradient boosting, a random forest of 150 trees, regularized quadratic discriminant analysis, 15 nearest neighbors and logistic regression. The attacker runs 20 rounds of random search with 12 candidates per step per round. A candidate is accepted if it lowers the surrogate's score, has a hazard score no lower than the original step's, and lies within the budget, in Euclidean distance, of the original; a typical step has norm about 4.5. Each task uses 1,000 hazardous steps for the stacked pipelines and 1,000 six-step pathways for the third. The joint miss rate for $k$ gates is averaged over every choice of $k$ of the six. On ordinary steps the 95\% intervals across tasks for the six-gate joint miss rates are $\pm$\MonSameCI, $\pm$\MonMixedCI{} and $\pm$\MonDistCI{} percentage points for the three pipelines.

\para{Scores.} Each of 2{,}000 agents carries a capacity $c$, standard normal at the start, and a share $s\in[0,1]$ of effort put into a gamble. Its take is $e^{0.3c}\,[(1-s)+s\,(0.8+1.5\,Z)]$ with $Z$ standard normal, and with probability $0.05\,s$ it imposes a loss of 20 on those it serves; welfare delivered is the take less that loss. The warrior rule lets the top tenth by take reproduce. Conferred rank requires in addition that the agent harmed nobody. The caretaker rule gives reproduction in proportion to welfare delivered, clipped to $[0,1]$. Under the floor the lowest fifth by welfare delivered leave no offspring and the rest reproduce equally, a selection intensity of 0.35 standard deviations against 1.75 for the top tenth. A gate of strength $g$ refuses each agent's gamble independently with probability $g$; a refused gamble pays nothing and harms nobody. Offspring inherit the parent's share and capacity plus normal noise of standard deviation 0.03 and 0.05. Populations start near $s=0.5$ and run 80 generations, with 40 populations per rule and 20 per point in the sweeps over prize share, harm probability and gate strength. The values are assumptions chosen to make the gamble worse on average for everyone, so that only the shape of the reward can favor it.

\para{Graded correlation.} The floor in Proposition~\ref{prop:sat} is exact when the shared cause is all-or-nothing. When correlation is graded there is no hard floor, and the lesson survives. Let each of $k$ monitors miss when a standard normal variable falls below its 10\% quantile, and let the variables share a factor with correlation $\rho$. The joint miss rates for $k=4$ and $k=8$ are \CopThreeFour\% and \CopThreeEight\% at $\rho=0.3$, and \CopSixFour\% and \CopSixEight\% at $\rho=0.6$, against 0.01\% and 0.000001\% under independence.

\para{Certifying a remainder by search.} If $n$ independent draws return $n_{1}$ classes of bypass exactly once each, the probability $\mu_{0}$ that the next draw returns an unseen class is estimated by $n_{1}/n$ with bias at most $1/n$, and with probability at least $1-\delta$, $\mu_{0}\le n_{1}/n+(2\sqrt{2}+\sqrt{3})\sqrt{\ln(3/\delta)/n}$ \cite{mcallester}. Setting $n_{1}=0$ and $\delta=0.05$ gives $\mu_{0}\le\varepsilon$ when $n\ge 85/\varepsilon^{2}$ with no assumption about how the remainder is divided. If the remainder is one class, $n\ge 3/\varepsilon$ suffices \cite{hanley}. For $\varepsilon=10^{-6}$ the two figures are $3\times10^{6}$ and $8.5\times10^{13}$ draws, each a full attempt by a capable system to defeat the cut.

% =====================================================================
\section{An authority record}
\label{app:record}

A deployment that claims the bound should keep a versioned record with the fields below. ``Unknown'' is an admissible entry, and it blocks the authority that would depend on it.

\begin{enumerate}[label=\arabic*.,leftmargin=1.5em,itemsep=1pt]
\item \textbf{Purpose and owners.} The task, the people affected, who approves, who reviews independently, and who can halt.
\item \textbf{Scope.} Action types, destinations, data classes, credentials, shared services, and every channel through which output reaches a person.
\item \textbf{Identity.} Model and runtime versions, tool schemas, policy version, and memory rules.
\item \textbf{Authority.} Limits per operation, aggregate allowances across all copies and retries, concurrency, duration, expiry, and how revocation works.
\item \textbf{Properties.} Forbidden effects, the invariant, the model it was proved or checked in, the trusted components, and known omissions.
\item \textbf{Usefulness.} The legitimate tasks, how completion is judged, and the measured rate of false blocks and delay.
\item \textbf{Measurements.} For $J_{k}$ and search power: the attacker, the campaign unit, the split between development and sealed test, the number of trials, the failures, and the interval.
\item \textbf{Transfer.} Which deployment conditions match those tested, and which opponents and circumstances are excluded.
\item \textbf{Operations.} Receipts of executed effects, monitoring coverage as observed, results of halt and recovery drills, and who outside can read the logs.
\item \textbf{Renewal.} Which changes void which evidence, what reduced scope applies meanwhile, and who decides.
\item \textbf{Footprint and consent.} The energy, water, materials and sites the deployment depends on, by place and time; the people affected who are neither owners nor users; what consent they gave, by what procedure, and how it can be withdrawn \cite{kogui,undrip}.
\item \textbf{Knowledge.} The provenance of the data and knowledge the system uses, the purposes each source permitted, who decides for knowledge held collectively, and what cannot in fact be withdrawn once used \cite{carroll}.
\end{enumerate}

No single score replaces the record. A high task score cannot repair an omitted channel, an invariant cannot certify the ethics of its specification, and a low failure rate per attempt cannot justify unbounded attempts.

% =====================================================================
\begingroup
\footnotesize
\begin{thebibliography}{199}

\bibitem{adamala} K.~P. Adamala et~al. Confronting risks of mirror life. \emph{Science} 386(6728):1351--1353, 2024.

\bibitem{alfonseca} M.~Alfonseca, M.~Cebrian, A.~Fern\'andez Anta, L.~Coviello, A.~Abeliuk, I.~Rahwan. Superintelligence cannot be contained: Lessons from computability theory. \emph{Journal of Artificial Intelligence Research} 70:65--76, 2021.

\bibitem{alshiekh} M.~Alshiekh, R.~Bloem, R.~Ehlers, B.~K\"onighofer, S.~Niekum, U.~Topcu. Safe reinforcement learning via shielding. In \emph{AAAI}, 2018.

\bibitem{amodei} D.~Amodei, C.~Olah, J.~Steinhardt, P.~Christiano, J.~Schulman, D.~Man\'e. Concrete problems in AI safety. arXiv:1606.06565, 2016.

\bibitem{anderson} J.~P. Anderson. \emph{Computer Security Technology Planning Study}. ESD-TR-73-51, US Air Force Electronic Systems Division, 1972.

\bibitem{anthropic-cyber} Anthropic. An alignment assessment of recent cybersecurity incidents. September 9, 2026. \url{https://www.anthropic.com/research/alignment-assessment-cybersecurity-incidents}

\bibitem{anthropic-aisi} Anthropic. An alignment assessment of the UK AISI cybertesting incident. October 23, 2026. \url{https://www.anthropic.com/research/alignment-assessment-ukaisi}

\bibitem{armstrong} S.~Armstrong, A.~Sandberg, N.~Bostrom. Thinking inside the box: Controlling and using an Oracle AI. \emph{Minds and Machines} 22(4):299--324, 2012.

\bibitem{baek} D.~D. Baek, X.~Li, A.~Gupta, T.~Mahbub, K.~Shi, M.~Tegmark, S.~Feng. Sycophancy towards researchers drives performative misalignment. arXiv:2606.08629, 2026.

\bibitem{bai} Y.~Bai et~al. Constitutional AI: Harmlessness from AI feedback. arXiv:2212.08073, 2022.

\bibitem{baker} B.~Baker et~al. Monitoring reasoning models for misbehavior and the risks of promoting obfuscation. arXiv:2503.11926, 2025.

\bibitem{balesni} M.~Balesni et~al. Towards evaluations-based safety cases for AI scheming. arXiv:2411.03336, 2024.

\bibitem{beckerman} S.~Beckerman et~al. Life histories, blood revenge, and reproductive success among the Waorani of Ecuador. \emph{Proceedings of the National Academy of Sciences} 106(20):8134--8139, 2009.

\bibitem{scientist} Y.~Bengio et~al. Superintelligent agents pose catastrophic risks: Can Scientist AI offer a safer path? arXiv:2502.15657, 2025.

\bibitem{berg} P.~Berg, D.~Baltimore, S.~Brenner, R.~O. Roblin, M.~F. Singer. Summary statement of the Asilomar conference on recombinant DNA molecules. \emph{Proceedings of the National Academy of Sciences} 72(6):1981--1984, 1975.

\bibitem{betley} J.~Betley et~al. Training large language models on narrow tasks can lead to broad misalignment. \emph{Nature}, 2026. doi:10.1038/s41586-025-09937-5.

\bibitem{betzig} L.~L. Betzig. \emph{Despotism and Differential Reproduction: A Darwinian View of History}. Aldine, 1986.

\bibitem{boehm1993} C.~Boehm. Egalitarian behavior and reverse dominance hierarchy. \emph{Current Anthropology} 34(3):227--254, 1993.

\bibitem{boehm1999} C.~Boehm. \emph{Hierarchy in the Forest: The Evolution of Egalitarian Behavior}. Harvard University Press, 1999.

\bibitem{bohme} M.~B\"ohme. STADS: Software testing as species discovery. \emph{ACM Transactions on Software Engineering and Methodology} 27(2), 2018.

\bibitem{bondarenko} A.~Bondarenko, D.~Volk, D.~Volkov, J.~Ladish. Demonstrating specification gaming in reasoning models. arXiv:2502.13295, 2025.

\bibitem{bostrom-ih} N.~Bostrom. Information hazards: A typology of potential harms from knowledge. \emph{Review of Contemporary Philosophy} 10:44--79, 2011.

\bibitem{bostrom2012} N.~Bostrom. The superintelligent will: Motivation and instrumental rationality in advanced artificial agents. \emph{Minds and Machines} 22(2):71--85, 2012.

\bibitem{bostrom-vwh} N.~Bostrom. The vulnerable world hypothesis. \emph{Global Policy} 10(4):455--476, 2019.

\bibitem{bowkis} A.~Bowkis, M.~D. Buhl, J.~Pfau, G.~Irving. Automated alignment is harder than you think. arXiv:2605.06390, 2026.

\bibitem{brennan} R.~L. Brennan. \emph{Generalizability Theory}. Springer, 2001.

\bibitem{burns} C.~Burns et~al. Weak-to-strong generalization: Eliciting strong capabilities with weak supervision. In \emph{ICML}, 2024.

\bibitem{carlsmith} J.~Carlsmith. Is power-seeking AI an existential risk? arXiv:2206.13353, 2022.

\bibitem{carro} M.~V. Carro et~al. Humans correct their judgments through debate, but weaker AI models do not. arXiv:2610.28539, 2026.

\bibitem{carroll} S.~R. Carroll et~al. The CARE Principles for Indigenous Data Governance. \emph{Data Science Journal} 19:43, 2020.

\bibitem{casper} S.~Casper et~al. Open problems and fundamental limitations of reinforcement learning from human feedback. \emph{Transactions on Machine Learning Research}, 2023.

\bibitem{chagnon} N.~A. Chagnon. Life histories, blood revenge, and warfare in a tribal population. \emph{Science} 239(4843):985--992, 1988.

\bibitem{chapman} S.~N. Chapman, J.~E. Pettay, V.~Lummaa, M.~Lahdenper\"a. Limits to fitness benefits of prolonged post-reproductive lifespan in women. \emph{Current Biology} 29(4):645--650, 2019.

\bibitem{chen} Y.-H. Chen, J.~Wen, J.~H. Kirchner. Automated researchers can mitigate well-characterized alignment failures. arXiv:2608.28945, 2026.

\bibitem{clopper} C.~J. Clopper, E.~S. Pearson. The use of confidence or fiducial limits illustrated in the case of the binomial. \emph{Biometrika} 26(4):404--413, 1934.

\bibitem{cronbach} L.~J. Cronbach, G.~C. Gleser, H.~Nanda, N.~Rajaratnam. \emph{The Dependability of Behavioral Measurements: Theory of Generalizability for Scores and Profiles}. Wiley, 1972.

\bibitem{dalrymple} D.~Dalrymple et~al. Towards guaranteed safe AI: A framework for ensuring robust and reliable AI systems. arXiv:2405.06624, 2024.

\bibitem{du} A.~Du et~al. Pattern and process in hominin brain size evolution are scale-dependent. \emph{Proceedings of the Royal Society B} 285(1873):20172738, 2018.

\bibitem{engels} J.~Engels, D.~D. Baek, S.~Kantamneni, M.~Tegmark. Scaling laws for scalable oversight. arXiv:2504.18530, 2025.

\bibitem{esvelt} K.~M. Esvelt. Delay, detect, defend: Preparing for a future in which thousands can release new pandemics. Geneva Paper 29/22, Geneva Centre for Security Policy, 2022.

\bibitem{fleming} K.~N. Fleming. A reliability model for common mode failures in redundant safety systems. Report GA-A13284, General Atomic Company, 1975.

\bibitem{ford} L.~R. Ford, D.~R. Fulkerson. Maximal flow through a network. \emph{Canadian Journal of Mathematics} 8:399--404, 1956.

\bibitem{flistatement} Future of Life Institute. Statement on superintelligence. Open statement, October 2025.

\bibitem{glowacki} L.~Glowacki, R.~Wrangham. Warfare and reproductive success in a tribal population. \emph{Proceedings of the National Academy of Sciences} 112(2):348--353, 2015.

\bibitem{good} I.~J. Good. The population frequencies of species and the estimation of population parameters. \emph{Biometrika} 40(3--4):237--264, 1953.

\bibitem{grace} K.~Grace et~al. Thousands of AI authors on the future of AI. arXiv:2401.02843, 2024.

\bibitem{control} R.~Greenblatt, B.~Shlegeris, K.~Sachan, F.~Roger. AI control: Improving safety despite intentional subversion. In \emph{ICML}, 2024.

\bibitem{greenblatt2024} R.~Greenblatt et~al. Alignment faking in large language models. arXiv:2412.14093, 2024.

\bibitem{password} R.~Greenblatt, F.~Roger, D.~Krasheninnikov, D.~Krueger. Stress-testing capability elicitation with password-locked models. arXiv:2405.19550, 2024.

\bibitem{metr} R.~Greenblatt, A.~Cotra, H.~Wijk. Brief independent investigation of agents' behavior, reasoning and collaboration in the OpenAI / Hugging Face hacking incident. METR, August 26, 2026. \url{https://metr.org/blog/2026-08-26-openai-hugging-face-incident-investigation/}

\bibitem{hadfield2017} D.~Hadfield-Menell, A.~Dragan, P.~Abbeel, S.~Russell. The off-switch game. In \emph{IJCAI}, 2017.

\bibitem{hotmess} A.~H\"agele, A.~P. Gema, H.~Sleight, E.~Perez, J.~Sohl-Dickstein. The hot mess of AI: How does misalignment scale with model intelligence and task complexity? arXiv:2601.23045, 2026.

\bibitem{hanley} J.~A. Hanley, A.~Lippman-Hand. If nothing goes wrong, is everything all right? Interpreting zero numerators. \emph{JAMA} 249(13):1743--1745, 1983.

\bibitem{hawkes} K.~Hawkes, J.~F. O'Connell, N.~G. Blurton Jones, H.~Alvarez, E.~L. Charnov. Grandmothering, menopause, and the evolution of human life histories. \emph{Proceedings of the National Academy of Sciences} 95(3):1336--1339, 1998.

\bibitem{henrich} J.~Henrich, F.~J. Gil-White. The evolution of prestige: Freely conferred deference as a mechanism for enhancing the benefits of cultural transmission. \emph{Evolution and Human Behavior} 22(3):165--196, 2001.

\bibitem{hhs} US Department of Health and Human Services. \emph{Screening Framework Guidance for Providers and Users of Synthetic Nucleic Acids}. 2023.

\bibitem{hong} S.~Z. Hong, A.~Kleinman, A.~Mathiowetz et~al. Measuring mid-2025 LLM-assistance on novice performance in biology. arXiv:2602.16703, 2026.

\bibitem{hubinger2019} E.~Hubinger, C.~van Merwijk, V.~Mikulik, J.~Skalse, S.~Garrabrant. Risks from learned optimization in advanced machine learning systems. arXiv:1906.01820, 2019.

\bibitem{sleeper} E.~Hubinger et~al. Sleeper agents: Training deceptive LLMs that persist through safety training. arXiv:2401.05566, 2024.

\bibitem{hudson} R.~Hudson. Corrigibility transformation: Constructing goals that accept updates. In \emph{ICML}, PMLR 306, 2026.

\bibitem{hvide} H.~K. Hvide. Tournament rewards and risk taking. \emph{Journal of Labor Economics} 20(4):877--898, 2002.

\bibitem{iasr} Y.~Bengio (chair) et~al. \emph{International AI Safety Report 2026}. February 2026. \url{https://internationalaisafetyreport.org/publication/international-ai-safety-report-2026}

\bibitem{iea} International Energy Agency. \emph{Energy and AI}. World Energy Outlook Special Report, April 2025.

\bibitem{ilo169} International Labour Organization. Indigenous and Tribal Peoples Convention, 1989 (No.~169).

\bibitem{irving} G.~Irving, P.~Christiano, D.~Amodei. AI safety via debate. arXiv:1805.00899, 2018.

\bibitem{jones} E.~Jones et~al. Forecasting rare language model behaviors. arXiv:2502.16797, 2025.

\bibitem{khan} A.~Khan et~al. Debating with more persuasive LLMs leads to more truthful answers. In \emph{ICML}, 2024.

\bibitem{klein} G.~Klein et~al. seL4: Formal verification of an OS kernel. In \emph{SOSP}, pages 207--220, 2009.

\bibitem{korbak} T.~Korbak et~al. Chain of thought monitorability: A new and fragile opportunity for AI safety. arXiv:2507.11473, 2025.

\bibitem{kulveit} J.~Kulveit, R.~Douglas, N.~Ammann, D.~Turan, D.~Krueger, D.~Duvenaud. Gradual disempowerment: Systemic existential risks from incremental AI development. arXiv:2501.16946, 2025.

\bibitem{kwa} T.~Kwa et~al. Measuring AI ability to complete long software tasks. arXiv:2503.14499, 2025.

\bibitem{lahdenpera} M.~Lahdenper\"a, V.~Lummaa, S.~Helle, M.~Tremblay, A.~F. Russell. Fitness benefits of prolonged post-reproductive lifespan in women. \emph{Nature} 428:178--181, 2004.

\bibitem{lampson} B.~W. Lampson. A note on the confinement problem. \emph{Communications of the ACM} 16(10):613--615, 1973.

\bibitem{lande} R.~Lande, S.~J. Arnold. The measurement of selection on correlated characters. \emph{Evolution} 37(6):1210--1226, 1983.

\bibitem{langosco} L.~Langosco, J.~Koch, L.~Sharkey, J.~Pfau, D.~Krueger. Goal misgeneralization in deep reinforcement learning. In \emph{ICML}, 2022.

\bibitem{leveson} N.~G. Leveson. \emph{Engineering a Safer World: Systems Thinking Applied to Safety}. MIT Press, 2011.

\bibitem{lewis} G.~Lewis, P.~Millett, A.~Sandberg, A.~Snyder-Beattie, G.~Gronvall. Information hazards in biotechnology. \emph{Risk Analysis} 39(5):975--981, 2019.

\bibitem{lovasz} L.~Lov\'asz, V.~Neumann-Lara, M.~Plummer. Mengerian theorems for paths of bounded length. \emph{Periodica Mathematica Hungarica} 9:269--276, 1978.

\bibitem{lynch} A.~Lynch, J.~Hughes, A.~Serrano, R.~Kirk, S.~R. Bowman. Agentic misalignment in summer 2026. Anthropic Alignment Science, July 2026. \url{https://alignment.anthropic.com/2026/agentic-misalignment-summer-2026/}

\bibitem{makins} O.~Makins, O.~Angelini, Z.~Shams, M.~Phuong. Multi-agent AI control: Distributed attacks hamper per-instance monitors. arXiv:2607.07368, 2026.

\bibitem{marks} S.~Marks et~al. Auditing language models for hidden objectives. arXiv:2503.10965, 2025.

\bibitem{mcallester} D.~McAllester, R.~E. Schapire. On the convergence rate of Good--Turing estimators. In \emph{COLT}, 2000.

\bibitem{meinke} A.~Meinke et~al. Frontier models are capable of in-context scheming. arXiv:2412.04984, 2024.

\bibitem{menger} K.~Menger. Zur allgemeinen Kurventheorie. \emph{Fundamenta Mathematicae} 10:96--115, 1927.

\bibitem{moon} R.~Moon, L.~R. Varshney. Containment verification: AI safety guarantees independent of alignment. arXiv:2605.09045, 2026.

\bibitem{kogui} Mamo Vangelio Moscote Alberto, Saka Rosa Mamatac\'an Sauna, of the Ezuama Mayor Kuamaka. An urgent message on artificial intelligence from the Kogui people. K\'aggaba People, Sierra Nevada de Gonawindua, Colombia. English translation by C.~Horstmann. Undated statement, posted September 2026. \url{https://movilizatorio.org/wp-content/uploads/2026/09/Gonawindua-Kogui-AI-Statement-.pdf}

\bibitem{mosleh} A.~Mosleh et~al. \emph{Procedures for Treating Common Cause Failures in Safety and Reliability Studies}. NUREG/CR-4780, US Nuclear Regulatory Commission, 1988.

\bibitem{normal} A.~Narayanan, S.~Kapoor. AI as normal technology. Knight First Amendment Institute at Columbia University, 2025.

\bibitem{needham} J.~Needham, G.~Edkins, G.~Pimpale, H.~Bartsch, M.~Hobbhahn. Large language models often know when they are being evaluated. arXiv:2505.23836, 2025.

\bibitem{omohundro} S.~M. Omohundro. The basic AI drives. In \emph{Artificial General Intelligence 2008}, pages 483--492. IOS Press, 2008.

\bibitem{controlroom} OpenAI. Control Room: A central system to track and stop model workloads across our research infrastructure. November 9, 2026. \url{https://openai.com/index/control-room/}

\bibitem{astra} OpenAI. GPT-6 Astra system card. September 3, 2026. \url{https://deploymentsafety.openai.com/gpt-6-astra}

\bibitem{ord} T.~Ord. \emph{The Precipice: Existential Risk and the Future of Humanity}. Bloomsbury, 2020.

\bibitem{papernot} N.~Papernot, P.~McDaniel, I.~Goodfellow. Transferability in machine learning: From phenomena to black-box attacks using adversarial samples. arXiv:1605.07277, 2016.

\bibitem{popper} K.~R. Popper. \emph{Conjectures and Refutations: The Growth of Scientific Knowledge}. Routledge, 1963.

\bibitem{price} G.~R. Price. Selection and covariance. \emph{Nature} 227:520--521, 1970.

\bibitem{reason} J.~Reason. Human error: Models and management. \emph{BMJ} 320:768--770, 2000.

\bibitem{metraudit} D.~Rein, A.~Kaufman. Red-teaming OpenAI's internal agent monitoring systems. METR, October 22, 2026 (updated October 29). \url{https://metr.org/blog/2026-10-22-red-teaming-openai-agent-monitoring/}

\bibitem{robacker} J.~T. Robacker. Min-max theorems on shortest chains and disjoint cuts of a network. Research Memorandum RM-1660, RAND Corporation, 1956.

\bibitem{saltzer} J.~H. Saltzer, M.~D. Schroeder. The protection of information in computer systems. \emph{Proceedings of the IEEE} 63(9):1278--1308, 1975.

\bibitem{schneider} F.~B. Schneider. Enforceable security policies. \emph{ACM Transactions on Information and System Security} 3(1):30--50, 2000.

\bibitem{schneier} B.~Schneier. Attack trees. \emph{Dr.\ Dobb's Journal} 24(12):21--29, 1999.

\bibitem{schoen} B.~Schoen, E.~Nitishinskaya, M.~Balesni, A.~H{\o}jmark, F.~Hofst\"atter et~al. Stress testing deliberative alignment for anti-scheming training. arXiv:2509.15541, 2025.

\bibitem{shamir} A.~Shamir. How to share a secret. \emph{Communications of the ACM} 22(11):612--613, 1979.

\bibitem{sheshadri} A.~Sheshadri et~al. Why do some language models fake alignment while others don't? arXiv:2506.18032, 2025.

\bibitem{shevlane} T.~Shevlane, A.~Dafoe. The offense-defense balance of scientific knowledge: Does publishing AI research reduce misuse? In \emph{AIES}, 2020.

\bibitem{sheyner} O.~Sheyner, J.~Haines, S.~Jha, R.~Lippmann, J.~M. Wing. Automated generation and analysis of attack graphs. In \emph{IEEE Symposium on Security and Privacy}, 2002.

\bibitem{singh} M.~Singh, L.~Glowacki. Human social organization during the Late Pleistocene: Beyond the nomadic-egalitarian model. \emph{Evolution and Human Behavior} 43(5):418--431, 2022.

\bibitem{soares} N.~Soares, B.~Fallenstein, E.~Yudkowsky, S.~Armstrong. Corrigibility. In \emph{AAAI Workshop on AI and Ethics}, 2015.

\bibitem{summerfield} C.~Summerfield et~al. Lessons from a chimp: AI ``scheming'' and the quest for ape language. arXiv:2507.03409, 2025.

\bibitem{taylor} J.~Taylor. Performance before outcomes: Early assessment, what women and men are chosen for, and the making of a general-purpose brain. Preprint, November 10, 2026.

\bibitem{terekhov} M.~Terekhov et~al. Adaptive attacks on trusted monitors subvert AI control protocols. arXiv:2510.09462, 2025.

\bibitem{tramer} F.~Tram\`er, N.~Papernot, I.~Goodfellow, D.~Boneh, P.~McDaniel. The space of transferable adversarial examples. arXiv:1704.03453, 2017.

\bibitem{trut} L.~N. Trut. Early canid domestication: The farm-fox experiment. \emph{American Scientist} 87(2):160--169, 1999.

\bibitem{turner} A.~M. Turner, L.~Smith, R.~Shah, A.~Critch, P.~Tadepalli. Optimal policies tend to seek power. In \emph{NeurIPS}, 2021.

\bibitem{aisi} UK AI Security Institute. Investigating alignment failures in AISI's incident. October 23, 2026. \url{https://www.aisi.gov.uk/blog/investigating-alignment-failures-in-aisis-incident}

\bibitem{undrip} United Nations General Assembly. United Nations Declaration on the Rights of Indigenous Peoples. Resolution 61/295, 13 September 2007.

\bibitem{urbina} F.~Urbina, F.~Lentzos, C.~Invernizzi, S.~Ekins. Dual use of artificial-intelligence-powered drug discovery. \emph{Nature Machine Intelligence} 4:189--191, 2022.

\bibitem{wash} US Nuclear Regulatory Commission. \emph{Reactor Safety Study: An Assessment of Accident Risks in US Commercial Nuclear Power Plants}. WASH-1400 (NUREG-75/014), 1975.

\bibitem{vanderweij} T.~van~der Weij, F.~Hofst\"atter, O.~Jaffe, S.~F. Brown, F.~R. Ward. AI sandbagging: Language models can strategically underperform on evaluations. In \emph{ICLR}, 2025.

\bibitem{rand} M.~J.~D. Vermeer, E.~Lathrop, A.~Moon. \emph{On the Extinction Risk from Artificial Intelligence}. RAND Corporation, RR-A3034-1, 2025.

\bibitem{walker} R.~S. Walker, K.~R. Hill, M.~V. Flinn, R.~M. Ellsworth. Evolutionary history of hunter-gatherer marriage practices. \emph{PLoS ONE} 6(4):e19066, 2011.

\bibitem{actionplan} The White House. \emph{Winning the Race: America's AI Action Plan}. July 2025.

\bibitem{wilsondaly} M.~Wilson, M.~Daly. Competitiveness, risk taking, and violence: The young male syndrome. \emph{Ethology and Sociobiology} 6(1):59--73, 1985.

\bibitem{woodburn} J.~Woodburn. Egalitarian societies. \emph{Man} 17(3):431--451, 1982.

\bibitem{wu} G.~Wu, J.~Hilton. Estimating the probabilities of rare outputs in language models. arXiv:2410.13211, 2024.

\bibitem{yudkowsky} E.~Yudkowsky, N.~Soares. \emph{If Anyone Builds It, Everyone Dies}. Little, Brown and Company, 2025.

\end{thebibliography}
\endgroup

\end{document}
